\RequirePackage[thinlines]{easybmat} %-- muss aufgrund eines Bugs vor etex und tikz geladen werden
\documentclass[a4paper, twoside, headsepline, index=totoc,toc=listof, fontsize=10pt, cleardoublepage=empty, headinclude, DIV=13, BCOR=13mm, titlepage,]{scrartcl}
\usepackage{scrtime} % KOMA, Uhrzeit ermoeglicht
\usepackage{scrpage2} % wie fancyhdr, nur optimiert auf KOMA-Skript, leicht andere Syntax
\usepackage{etoolbox}
\usepackage{letltxmacro}

% Compiler abhaengige IFs
\usepackage{ifxetex,ifluatex}
\newif\ifxetexorluatex
\ifxetex
  \xetexorluatextrue
\else
  \ifluatex
    \xetexorluatextrue
  \else
    \xetexorluatexfalse
  \fi
\fi

%--Farbdefinitionen
%-- muss vor tikz geladen werden
\usepackage[usenames, table, x11names]{xcolor}
\definecolor{dark_gray}{gray}{0.45}
\definecolor{light_gray}{gray}{0.6}

%--Zum Zeichnen
%-- muss vor polyglossia bzw. babel geladen werden
\usepackage{tikz}
\usepackage{tikz-cd}
\usetikzlibrary{external}
\tikzset{>=latex}
\usetikzlibrary{shapes,arrows,intersections}
\usetikzlibrary{calc,3d}
\usetikzlibrary{decorations.pathreplacing,decorations.markings}
\usetikzlibrary{angles}

%-- Konfiguration von tikzexternalize
\tikzexternalize[prefix=tikz/, up to date check=diff]

\ifxetexorluatex
\tikzset{external/system call={xelatex \tikzexternalcheckshellescape %-- verwende LuaLaTeX, wegen dynamischer Speicherallokation
    -halt-on-error -interaction=batchmode -jobname "\image" "\texsource"}}
\else
\tikzset{external/system call={pdflatex \tikzexternalcheckshellescape 
    -halt-on-error -interaction=batchmode -jobname "\image" "\texsource"}}
\fi

%-- tikzexternalize fuer tikzcd deaktivieren, da inkompatibel
\AtBeginEnvironment{tikzcd}{\tikzexternaldisable}
\AtEndEnvironment{tikzcd}{\tikzexternalenable}

%-- um Inkompatibilitaeten von quotes und polyglossia bzw. babel zu vermeiden
\tikzset{
  every picture/.append style={
    execute at begin picture={\shorthandoff{"}},
    execute at end picture={\shorthandon{"}}
  }
}
\usetikzlibrary{quotes}
\usepackage{pgfplots}
\usepgfplotslibrary{colormaps}
\pgfplotsset{compat=1.10}



%-- Mathe
\usepackage{mathtools} % beinhaltet amsmath
\mathtoolsset{showonlyrefs, centercolon}
\newtagform{brackets}[\textbf]{[}{]}
\usetagform{brackets}
\usepackage{wasysym}
\usepackage{amssymb} %zusätzliche Symbole
\usepackage{latexsym} %zusätzliche Symbole
\usepackage{stmaryrd} %für Blitz
\usepackage{nicefrac} %schräge Brüche
\usepackage{cancel} %Befehle zum Durchstreichen
\usepackage{mathdots}
% \usepackage{bm}
\DeclareSymbolFont{bbold}{U}{bbold}{m}{n}
\DeclareSymbolFontAlphabet{\mathbbold}{bbold}
\newcommand{\mathds}[1]{\mathbb{#1}} %Um Kompatibilitaet mit frueheren Benutzung von dsfont herzustellen
\newcommand{\ind}{\mathbbold{1}} % charakteristische-Funktion-Eins
\def\mathul#1#2{\color{#1}\underline{{\color{black}#2}}\color{black}} %farbiges Untersteichen im Mathe-Modus

%-- Alles was mit Schrift und XeteX zu tun hat
\ifxetexorluatex
    % XeLaTeX or LuaTeX
	\ifxetex
		\usepackage{mathspec} %beinhaltet fontspec 
	\else
		\usepackage[no-math]{fontspec}
	\fi
    \usepackage{polyglossia} %babel-ersatz
    \setmainlanguage[spelling=new,babelshorthands=false]{german}
	\newcommand\glqq{"}
	\newcommand\grqq{"}
	\defaultfontfeatures{Mapping=tex-text, WordSpace={1.4}} %
    \setmainfont[Ligatures=Common, BoldFont={* Bold}, ItalicFont={* Light Italic},SmallCapsFont={Linux Libertine O}, SmallCapsFeatures={Letters=SmallCaps}]{Source Sans Pro}
	\setsansfont[Scale=MatchLowercase,Ligatures=Common, BoldFont={* Medium}]{Ubuntu}
	\ifxetex
		\setallmonofonts[Scale=MatchLowercase, ItalicFont={*}]{Consolas} 
	\else
		\setmonofont[Scale=MatchLowercase, ItalicFont={*}]{Consolas}
	\fi
	\usepackage{xltxtra}
	\usepackage{fontawesome}
	\usepackage{microtype}
\else
    % default: pdfLaTeX
    \usepackage[ngerman]{babel}
    \usepackage[T1]{fontenc}
    \usepackage[utf8]{inputenc}
	\usepackage{textcomp} %verhindert ein paar Fehler bei den Fonts
    \usepackage[babel=true, tracking=true]{microtype}
	\usepackage{lmodern}
	\renewcommand{\familydefault}{\sfdefault} 
\fi


%--Mixed
\usepackage[neverdecrease]{paralist}
\usepackage[german=quotes]{csquotes}
\usepackage{makeidx}
\usepackage{booktabs}
\usepackage{wrapfig}
\usepackage{float}
\usepackage[margin=10pt, font=small, labelfont={sf, bf}, format=plain, indention=1em]{caption}
\captionsetup[wrapfigure]{name=Abb. }
\usepackage{stackrel}
\usepackage{ifthen}
\usepackage{multicol}
\usepackage{qrcode}
\flushbottom

%--Hyphenation
\hyphenation{di-men-si-o-na-le}


%--Unterstreichung
\usepackage[normalem]{ulem}
\setlength{\ULdepth}{1.8pt}

%--Indexverarbeitung
\newcommand{\bet}[1]{\textbf{#1}}
\newcommand{\Index}[1]{\textbf{#1}\index{#1}}
\makeindex
\setindexpreamble{{\noindent \itshape Die \emph{Seitenzahlen} sind mit Hyperlinks zu den entsprechenden Seiten versehen, also anklickbar} \ifxetexorluatex \faHandUp\fi \par \bigskip}
\renewcommand{\indexpagestyle}{scrheadings}


%--Marginnote & todonotes
\usepackage{marginnote}
\renewcommand*{\marginfont}{\color{dark_gray} \itshape \footnotesize }
\usepackage[textsize=small]{todonotes}

%--Todonotes mit tikz/externalize kompatibel machen
\LetLtxMacro{\oldtodo}{\todo}
\renewcommand{\todo}[2][]{\tikzexternaldisable\oldtodo[#1]{#2}\tikzexternalenable}
\LetLtxMacro{\oldmissingfigure}{\missingfigure}
\renewcommand{\missingfigure}[2][]{\tikzexternaldisable\oldmissingfigure[{#1}]{#2}\tikzexternalenable}

%--Konfiguration von Hyperref pdfstartview=FitH, 
\usepackage[hidelinks, pdfpagelabels,  bookmarksopen=true, bookmarksnumbered=true, linkcolor=black, urlcolor=SkyBlue2, plainpages=false,pagebackref, citecolor=black, hypertexnames=true, pdfauthor={Jannes Bantje}, pdfborderstyle={/S/U}, linkbordercolor=SkyBlue2, colorlinks=false]{hyperref}

% \let\orighref\href
\ifxetexorluatex
\newcommand{\hrefsym}[2]{\href{#1}{\texttt{#2\,\faExternalLink}}}
\newcommand{\hrefsymmail}[2]{\href{#1}{\texttt{\faEnvelope\,#2}}}
\renewcommand{\url}[1]{\hrefsym{#1}{\nolinkurl{#1}}}
\fi
\providecommand{\hrefsym}{\href}
% \usepackage{cleveref}
% \crefformat{footnote}{#2\footnotemark[#1]#3}

%--Römische Zahlen
\newcommand{\RM}[1]{\MakeUppercase{\romannumeral #1{}}}

%--Punkte (nach hyperref)
\usepackage{ellipsis}

%%--Abkürzungen etc.
\newcommand{\light}{\text{\Large $\lightning$}}

%-- Definitionen von weiteren Mathe-Befehlen
\DeclareMathOperator{\re}{Re} %Realteil
\DeclareMathOperator{\im}{Im} %Imaginaerteil
\DeclareMathOperator{\id}{id} %identische Abbildung
\DeclareMathOperator{\Sp}{Sp} %Spur
\DeclareMathOperator{\sgn}{sgn} %Signum
\DeclareMathOperator{\alt}{Alt} %Alternierende n-linearFormen
\DeclareMathOperator{\End}{End} %Endomorphismen
\DeclareMathOperator{\Vol}{Vol} %Volumen
\DeclareMathOperator{\dom}{dom} %Domain
\DeclareMathOperator{\Hom}{Hom} %Homomorphismen
\DeclareMathOperator{\bild}{Bild} %Bild
\DeclareMathOperator{\Kern}{Kern }
\DeclareMathOperator{\rg}{rg} %Rang
\DeclareMathOperator{\Rg}{Rg} %Rang
\DeclareMathOperator{\diam}{diam} %Durchmesser
\DeclareMathOperator{\dist}{dist} %Distanz
\DeclareMathOperator{\grad}{grad} %Gradient
\DeclareMathOperator{\dive}{div} %Gradient
\DeclareMathOperator{\rot}{rot} %Rotation
\DeclareMathOperator{\hess}{Hess} %Hesse-Matrix
\DeclareMathOperator{\koker}{Koker} %Kokern
\DeclareMathOperator{\aut}{Aut} %Automorphismen
\DeclareMathOperator{\ord}{ord} %Ordnung
\DeclareMathOperator{\ggT}{ggT} %ggT
\DeclareMathOperator{\kgV}{kgV} %kgV
\DeclareMathOperator{\Gr}{Gr} %Gerade
\DeclareMathOperator{\Kr}{Kr} %Kreis
\DeclareMathOperator{\Char}{char} %Charakteristik
\DeclareMathOperator{\Aut}{Aut} %Automorphismen
% \DeclareMathOperator{\D}{D} %formale Ableitung
\newcommand{\D}{\ensuremath{\mathrm{D}\mkern-1.5mu}}
\newcommand{\mathd}{\ensuremath{\mathrm{d}\mkern-1.0mu}}
\newcommand{\Tmap}{\ensuremath{\mathrm{T}\mkern-1.0mu}}
\newcommand{\opL}{\ensuremath{\mathrm{L}\mkern-0.6mu}}
\DeclareMathOperator{\cov}{cov} %Kovarianz
\DeclareMathOperator{\Gal}{Gal} %Galois
\DeclareMathOperator{\supp}{supp} %Träger
\DeclareMathOperator{\Sym}{Sym} %Symmetrische Gruppe
\DeclareMathOperator{\Zyl}{Zyl} %Zylinder
\DeclareMathOperator{\Mod}{Mod} %Moduln
\DeclareMathOperator{\EPK}{EPK} %Einpunktkompaktifizierung
\DeclareMathOperator{\conj}{conj} %Konjugation
\DeclareMathOperator{\ann}{ann} %Annulator
\DeclareMathOperator{\tor}{tor} %Torsionsmodul
\DeclareMathOperator{\Int}{Int} %Interior/Inneres
\DeclareMathOperator{\Br}{Br} %Brauerergruppe
\DeclareMathOperator{\GL}{GL} %allgemeine lineare Gruppe
\DeclareMathOperator{\SL}{SL} %spezielle lineare Gruppe
\DeclareMathOperator{\SU}{SU} %spezielle unitäre Gruppe
\DeclareMathOperator{\Tr}{Tr} %Spur/Trace
\DeclareMathOperator{\conv}{conv} %Konvexe Teilmenge
% \DeclareMathOperator{\tr}{tr}
\newcommand{\tr}{\mathrm{tr}}
\newcommand{\w}{\mathrm{w}}
\newcommand{\weakT}[1]{\ensuremath{\mathcal{T}_{#1}^{\mkern+1.0mu\text{\raisebox{0.4ex}{$\mathrm{w}$}}}}}
\newcommand{\weakTstar}[1]{\ensuremath{\mathcal{T}_{#1}^{\mkern+1.0mu\text{\raisebox{0.4ex}{$\mathrm{w}$}}^*}}}
\newcommand{\rand}[1]{\ensuremath{\partial^{\scriptscriptstyle #1}}}


% -- Kategorien
\DeclareMathOperator{\Ob}{Ob} %Objekte
\DeclareMathOperator{\Mor}{Mor} %Morphismen

% \DeclareMathOperator{\TOP}{TOP} %Topologische Räume
\newcommand{\TOP}{\textsc{Top}}
% \DeclareMathOperator{\HTOP}{HTOP} %Topologische Räume + Homotopieäquivalenz
\newcommand{\HTOP}{\textsc{HTop}}
% \DeclareMathOperator{\VR}{VR} %Vektorräume
\newcommand{\VR}{\textsc{VR}}
% \DeclareMathOperator{\MOD}{MOD} %Moduln
\newcommand{\MOD}{\textsc{Mod}}
% \DeclareMathOperator{\MENGEN}{MENGEN} %Mengen
\newcommand{\MENGEN}{\textsc{Mengen}}
% \DeclareMathOperator{\MAN}{MAN} %Mannigfaltigkeiten
\newcommand{\MAN}{\textsc{Man}}
% \DeclareMathOperator{\GRUPPEN}{GRUPPEN} %Gruppen
\newcommand{\GRUPPEN}{\textsc{Gruppen}}
% \DeclareMathOperator{\ABELGRUPPEN}{ABEL.GRUPPEN} %abelsche Gruppen
\newcommand{\ABELGRUPPEN}{\textsc{Abel.Gruppen}}
% \DeclareMathOperator{\KAT}{KAT} %Kategorien
\newcommand{\KAT}{\textsc{Kat}}
% \DeclareMathOperator{\FUN}{FUN} %Funktoren
\newcommand{\FUN}{\textsc{Fun}}


\newcommand{\op}{\mathrm{op}}

\newcommand{\Ker}{\ensuremath{\Kern \,}}
\newcommand{\zirkel}{\rotatebox[origin = c]{-90}{$\varangle$}}
\newcommand{\wirkung}{\!\curvearrowright\!}

\newcommand{\oE}{\OE}

%--Beweisende
\newcommand{\bewende}{\ifmmode \tag*{$\square$} \else \hfill $\square$ \fi}

%--Volumen
\newcommand{\vol}[1]{\ensuremath{\Vol_{#1}}}

%--Skalarprodukt
\DeclarePairedDelimiterX\sprod[2]{\langle}{\rangle}{#1\,\delimsize\vert\,#2}

%--Betrag, Gaußklammer
\DeclarePairedDelimiter{\abs}{\lvert}{\rvert}
\DeclarePairedDelimiter{\ceil}{\lceil}{\rceil}

%--Norm
\DeclarePairedDelimiter\doppelstrich{\Vert}{\Vert}
\newcommand{\norm}[2][\relax]{
\ifx#1\relax \ensuremath{\doppelstrich*{#2}}
\else \ensuremath{\doppelstrich*{#2}_{#1}}
\fi}


%--Umklammern
\DeclarePairedDelimiter\enbrace{(}{)}
\DeclarePairedDelimiter\benbrace{[}{]}
\DeclarePairedDelimiter\lenbrace{<}{>}

%--Mengen
\DeclarePairedDelimiterX\mengenA[1]{\lbrace}{\rbrace}{#1}
\DeclarePairedDelimiterX\mengenB[2]{\lbrace}{\rbrace}{#1\, \delimsize\vert \, #2}

\makeatletter
\newcommand{\set}[2][\relax]{
\ifx#1\relax \ensuremath{
\mengenA*{#2}}
\else \ensuremath{%
  \mengenB*{#1}{#2}}
\fi}
\makeatother


%--offen und abgeschlossen
\newcommand{\off}{\! \stackrel[\text{offen}]{}{\subset} \!}
\newcommand{\abg}{\! \stackrel[\text{abg.}]{}{\subset} \!}

%--Differential
\newcommand{\diff}[2]{\ensuremath{\frac{{\partial #1}}{{\partial #2}} }}
\newcommand{\diffd}[2]{\ensuremath{\frac{\mathd #1}{\mathd #2} }}

%--Diagonale Linien für Matrizen
\newcommand\tikzmark[1]{%
  \tikzexternaldisable\tikz[overlay,remember picture,baseline] \node [anchor=base] (#1) {};\tikzexternalenable}

\newcommand\MyLine[3][]{%
  \tikzexternaldisable\begin{tikzpicture}[overlay,remember picture]
    \draw[#1] (#2.south) -- (#3.north);
  \end{tikzpicture}\tikzexternalenable}
  
%--Jordankasten
\newcommand{\Jord}{ \tikzexternaldisable \begin{BMAT}{|ccc|}{|ccc|} %
				0 & & \\ %
				1\tikzmark{z} &  \diagdown & \\ %
				& \tikzmark{u}1 \MyLine[thick]{z}{u} &  0 %
			\end{BMAT} \tikzexternalenable}
			
%--Jordankasten mit Lambda
\newcommand{\JordLa}[1][\empty]{%
	\tikzexternaldisable\ifthenelse{\equal{#1}{\empty}}
		{\begin{BMAT}{|ccc|}{|ccc|} %
				\lambda & & \\ %
				1\tikzmark{z} &  \diagdown & \\ %
				& \tikzmark{u}1 \MyLine[thick]{z}{u} &  \lambda  %
			\end{BMAT}}
		{\begin{BMAT}{|ccc|}{|ccc|} %
				\lambda_{#1}  & & \\ %
				1\tikzmark{z} &  \diagdown & \\ %
				& \tikzmark{u}1 \MyLine[thick]{z}{u} &  \lambda_{#1}  %
			\end{BMAT}}\tikzexternalenable}

%--Inhaltsverzeichnis
\usepackage[tocindentauto]{tocstyle}
\usetocstyle{KOMAlike}			

\newcommand{\fach}{Differentialgeometrie \RM{1}.}
\newcommand{\semester}{WiSe 2014}
\newcommand{\homepage}{https://wwwmath.uni-muenster.de/u/frederik.witt/diffgeoI.html}
\newcommand{\prof}{Frederik Witt}
\newcommand{\verfasser}{Arne Grauer}
\providecommand{\verfasser}{Jannes Bantje}

%--Konfiguration von scrheadings
\setheadsepline{1pt}[\color{light_gray}]
\pagestyle{scrheadings}
%\fancyhf[H, F]{}
\clearscrheadfoot

\providecommand{\shortFach}{\fach}
\lehead{ \includegraphics[height=0.5 cm,keepaspectratio]{../!config/Bilder/Logo_WWU_Muenster_light_gray.pdf}}
\rehead{\rule{0cm}{0.5cm}\footnotesize \sffamily \color{light_gray} \verfasser -- Mitschrift \shortFach}
\lohead{\rule{0cm}{0.5cm} \footnotesize \sffamily \color{light_gray} Stand: \today \; \thistime[:]}
\rohead{\includegraphics[height=0.5 cm,keepaspectratio]{../!config/Bilder/fb10logo_gray.pdf}}


\ofoot[{ \color{dark_gray} \LARGE \sffamily \thepage}]{{ \color{dark_gray} \LARGE \sffamily \thepage}} %hier wir auch der plain Stil bearbeitet!
\automark{section}
\ifoot{ \color{dark_gray} \small \leftmark}

%--Metadaten
\author{\verfasser}
\titlehead{\includegraphics[height=1.5cm, keepaspectratio]{../!config/Bilder/Logo_WWU_Muenster.pdf}%
\hfill \includegraphics[height=1.3cm, keepaspectratio]{../!config/Bilder/fb10logo.pdf}}
\title{Skript \fach}
\subtitle{Mitschrift der Vorlesung  \enquote{\fach} von \prof}

\publishers{
Aktuelle Version verfügbar bei: \bigskip\\
\normalsize
\begin{minipage}[t]{0.45\textwidth}
	\centering{
	\qrcode[height=3.5cm, version=4]{https://github.com/JaMeZ-B/latex-wwu}
	\medskip\\
	\includegraphics[height=0.4cm, keepaspectratio]{../!config/Bilder/github_octo.pdf} 
	\includegraphics[height=0.4cm, keepaspectratio]{../!config/Bilder/GitHub_Logo.pdf} (inklusive Sourcecode)\\
	{\small\url{https://github.com/JaMeZ-B/latex-wwu}}}
\end{minipage}
\hfill
\begin{minipage}[t]{0.45\textwidth}
	\centering{
	\qrcode[height=3.5cm, version=4]{btsync://B6WH2DISQ5QVYIRYIEZSF4ZR2IDVKPN3I?n=latex_share}
	\medskip\\
	\includegraphics[height=0.4cm, keepaspectratio]{../!config/Bilder/bt_sync_logo.pdf} 
	{\large \textbf{Bittorrent} Sync} \\
	\texttt{B6WH2DISQ5QVYIRYIEZSF4ZR2IDVKPN3I}
	}
\end{minipage}
}



\begin{document}
\pagenumbering{Roman}
\maketitle
\begin{abstract}
\section*{Vorwort --- Mitarbeit am Skript}
Dieses Dokument ist eine Mitschrift aus der Vorlesung \enquote{\fach, \semester}, gelesen von \prof. Der Inhalt entspricht weitestgehend dem Tafelanschrieb. Für die
Korrektheit des Inhalts übernehme ich keinerlei Garantie! Für Bemerkungen und Korrekturen -- und seien es nur Rechtschreibfehler -- bin ich sehr dankbar. 
Korrekturen lassen sich prinzipiell auf drei Wegen einreichen: 
\begin{itemize}
	\item Persönliches Ansprechen in der Uni, Mails an \hrefsymmail{mailto:\mail}{\mail} (gerne auch mit annotieren PDFs) 
	\item \emph{Direktes} Mitarbeiten am Skript: Den Quellcode poste ich auf GitHub (siehe oben), also stehen vielfältige Möglichkeiten der Zusammenarbeit zur Verfügung:
	Zum Beispiel durch Kommentare am Code über die Website und die Kombination Fork + Pull Request. Wer sich verdient macht oder ein Skript zu einer Vorlesung, die 
	ich nicht besuche, beisteuern will, dem gewähre ich gerne auch Schreibzugriff.
	
	Beachten sollte man dabei, dass dazu ein Account bei \url{github.com} notwendig ist, der allerdings ohne Angabe von persönlichen Daten angelegt werden kann. 
	Wer bei GitHub (bzw. dem zugrunde liegenden Open-Source-Programm \enquote{\texttt{git}}) -- verständlicherweise -- Hilfe beim Einstieg braucht, dem helfe ich gerne 
	weiter. Es gibt aber auch zahlreiche empfehlenswerte Tutorials im Internet.\footnote{zB. \url{https://try.github.io/levels/1/challenges/1}, ist auf Englisch, aber dafür 
	interaktives LearningByDoing}
	\item \emph{Indirektes} Mitarbeiten: \TeX-Dateien per Mail verschicken. 
	
	Dies ist nur dann sinnvoll, wenn man einen ganzen Abschnitt ändern möchte (zB. einen alternativen Beweis geben), da ich die Änderungen dann per Hand einbauen muss! Ich freue mich aber auch über solche Beiträge!
\end{itemize}
\section*{Vorlesungshomepage}
{\centering 
\begin{minipage}[c][][c]{\textwidth}
	\centering \qrcode[height=3cm]{\homepage} \medskip\\
	\footnotesize \url{\homepage}
\end{minipage}
\par}

\section*{Über die Differentialgeometrie}
\begin{description}
	\item[1. Differentialgeometrie] Mannigfaltigkeiten $\Rightarrow $ Differentialtopologie
	\item[2. Riemannsche Geometrie] Längen und Winkel messen $\leadsto$ Geodäte (lokal) kürzeste Verbindung zwischen 2 Punkten einer Mannigfaltigkeit.
	
	Krümmung
	
	Allgemeine Relativitätstheorie: Krümmung = Gravitation ist für solche Effekte verantwortlich
\end{description}
\end{abstract}

\setcounter{page}{1}
\tableofcontents
\cleardoubleoddemptypage
\pagenumbering{arabic}
\setcounter{page}{1}


\section{Differenzierbare Mannigfaltigkeiten} % (fold)
\label{sec:1}
Konvention: glatt = $C^\infty$.\\
Wiederholung: Diffeomorphismus $f : U \subseteq \mathds{R}^n \to \mathds{R}^n$ glatt, bijektiv und $f ^{-1}$ glatt.

\subsection{Untermannigfaltigkeiten} % (fold)
\label{sub:11}

\subsubsection[Definition: Submersion, Immersion und Étale]{Definition} % (fold)
\label{ssub:11}
Eine Abbildung $f : U \subseteq \mathds{R}^n \to \mathds{R}^m$ glatt heißt \Index{Submersion}/\Index{Immersion}/\Index{Étale} in $x \in U$, falls
$\mathd_x f : \mathds{R}^n \to \mathds{R}^m$ surjektiv/injektiv/Isomorphismus ist. \marginnote{Insbesondere $n\ge m$/$n \le m$/$n=m$}. $f$ heißt schlechthin
\bet{Submersion}/\bet{Immersion}/\bet{Étale}, falls $f$ Submersion/Immersion/Étale für alle $x \in U$ ist.
% subsubsection 111 (end)

\subsubsection[Bemerkung: Submersiv/immersiv/étale sind offene Bedingungen]{Bemerkung} % (fold)
\label{ssub:12}
Submersiv/immersiv/étale sind offene Bedingungen, d.h. ist $f$ Submersion/Immersion/Étale in $x$, so auch in einer hinreichend kleinen Umgebung von $x$ 
(vergleiche auch Satz 15 unten)
% subsubsection 12 (end)

\subsubsection{Satz über inverse Funktionen (inverse function theorem)} % (fold)
\label{ssub:13}
Sei $f : U \subseteq \mathds{R}^n \to \mathds{R}^m$ étale in $x$. Dann existiert eine offene Menge $x \in V \subset U \subset \mathds{R}^n$, sodass $f\big|_{V}$ ein
Diffeomorphismus auf seinem Bild ist, d.h. $f\big|_{V} : V \xrightarrow{\sim} f(V) $ ist bijektiv und $f\big|_V$, $(f\big|_V)^{-1}$ glatt.
% subsubsection 13 (end)

\subsubsection{Satz über implizite Funktionen (implicit function theorem)} % (fold)
\label{ssub:14}
Sei $f : U \subset \mathds{R}^{n+k} \to \mathds{R}^n$ eine Submersion in $0$ mit $f(0)=0$. Dann existiert eine glatte Funktion $g$ auf einer Umgebung von $0 \in \mathds{R}^n$
mit Werten in $\mathds{R}^k$, d.h.
\[
	g=g(x_1, \ldots , x_n) = \begin{pmatrix}
		g_1(x_1, \ldots ,x_n) \\
		\vdots \\
		g_k(x_1, \ldots ,x_n)
	\end{pmatrix}
\]
mit
\[
	f \enbrace[\big]{x_1, \ldots ,x_n , g_1(x_1, \ldots , x_n), \ldots , g_k(x_1, \ldots ,x_n)} = f \enbrace[\big]{x,g(x)} = 0
\]
% subsubsection 14 (end)

\subsubsection{Satz (Normalenformen von Submersionen/Immersionen)} % (fold)
\label{ssub:15}
\begin{enumerate}[a)]
	\item Sei $f : U \subseteq \mathds{R}^{n+k} \to \mathds{R}^n$ eine Submersion mit $f(0)=0$. Dann existiert ein Diffeomorphismus 
	$\varphi : 0 \in V \subset \mathds{R}^{n+k} \to \mathds{R}^{n+k}$ mit $\varphi(V) \subset U$ und 
	$f \circ \varphi(x_1, \ldots , x_{n+k}) = (x_1, \ldots , x_n) = \pi_{\mathds{R}^n}(x_1, \ldots , x_{n+k})$, d.h. modulo einer Diffeomorphismus ist jede Submersion 
	äquivalent zu einer Projektion.
	\item Sei $f : U \subset \mathds{R}^n \to \mathds{R}^{n+k}$ eine Immersion in $0$, $f(0)=0$. Dann existiert ein Diffeomorphismus 
	\[
		\varphi : 0 \in V \subset \mathds{R}^{n+k} \to \mathds{R}^{n+k}
	\]
	mit $\varphi(0)=0$ und 
	\[
		\varphi \circ f (x_1, \ldots ,x _n) = \enbrace*{x_1, \ldots , x_n, 0 , \ldots , 0} = i(x_1, \ldots , x_n) 
	\]
	d.h. modulo Diffeomorphismen ist $f$ äquivalent zur kanonischen Einbettung von $\mathds{R}^n \hookrightarrow \mathds{R}^{n+k}$.
\end{enumerate}
\minisec{Beweis}
\begin{enumerate}[(i)]
	\item Wähle Koordinaten $x_1, \ldots , x_n$ mit $\mathd_0 f = (A \, B)$ mit $A \in \mathds{R}^{n \times n}$, $B \in \mathds{R}^{k \times n}$ mit $\det A \not= 0$.
	Definiere $F : U \subset \mathds{R}^n \times \mathds{R}^k \to \mathds{R}^n \times \mathds{R}^k $ durch $(p,q) \mapsto \enbrace*{f(p,q), q}$. Dann gilt
	\[
		\mathd_{(0,0)} F = \begin{pmatrix}
			\mathd_0 f & \\
			0 & \id_{k \times k}
		\end{pmatrix} = \begin{pmatrix}
			A & B \\
			0 & \id_{k \times k} 
		\end{pmatrix}
		\Longrightarrow \det \mathd_0 F \not= 0
	\]
	Mit \ref{ssub:13} folgt: $\varphi = F ^{-1}$ Diffeomorphismus um $0 \in \mathds{R}^{n+k}$ herum.
	\item Definiere $F : U \times \mathds{R}^k \to \mathds{R}^n \times \mathds{R}^k$. $F(p,q) = f(p)+ (0,q)$.
	\[
		\mathd_{(0,0)} F = \begin{pmatrix}
			\mathd_0 f & 0 \\
			0 & \id_{k \times k}
		\end{pmatrix},
	\]
	also $\det \mathd_{(0,0)} F \not= 0$. $\Rightarrow  \exists \varphi = F ^{-1}$ um $0$ mit $\varphi(f(p))=\varphi \circ F (p,0) = (p,0) = i(p)$ \bewende
\end{enumerate}
\minisec{Bemerkung}
Durch Komposition mit Translationen können Submersionen/Immersionen mit $F(x)=y$ (nicht notwendigerweise $=0$) behandelt werden.
% subsubsection 15 (end)

\subsubsection[Definition: Untermannigfaltigkeit]{Definition} % (fold)
\label{ssub:16}
$M \subset \mathds{R}^{n+k}$ heißt $n$-dimensionale \Index{Untermannigfaltigkeit}, falls $\forall x \in M : \exists x \in U \subset \mathds{R}^{n+k}$ Umgebung, sowie eine
Submersion $f : U \to \mathds{R}^k$ mit $U \cap M = f ^{-1}(0)$. \marginnote{Verallgemeinerung des $n$-dimensionalen Unterraumes}
% subsubsection 16 (end)

\subsubsection[Beispiele für Untermannigfaltigkeiten]{Beispiele} % (fold)
\label{ssub:17}
\begin{enumerate}[(i)]
	\item Die $n$-\Index{Sphäre} $S^n := \set[x=(x_1, \ldots ,x_{n+1}) \in \mathds{R}^{n+1}]{\abs*{x}^2=1 } $ ist eine $n$-dimensionale Untermannigfaltigkeit des 
	$\mathds{R}^{n+1}$. Definiere $f : \mathds{R}^{n+1} \to \mathds{R}$ durch $f(x_1, \ldots , x_{n+1}) = x_1^2 + \ldots + x_{n+1}^2 -1$. $\Rightarrow f ^{-1}(0) = S^n$ nach
	Definition.
	\[
		\mathd_x f = (2 x_1, \ldots , 2 x_{n+1}) : \mathds{R}^{n+1} \to \mathds{R}
	\]
	\item Die hyperbolischen Räume $H_c^n := \set[x=(x_1, \ldots ,x_{n+1}) \in \mathds{R}^{n+1}]{-x_1^2 + \sum_{i=2}^{n+1}x_i^2 = c}$\begin{figure}[H]
\centering{
\begin{tikzpicture}[xscale=0.6, yscale=0.8]
	    \begin{axis}[view={40}{20}]
	       \addplot3[surf,
	       colormap/cool,
	       samples=60,
	       domain=-2*pi:2*pi,y domain=-2*pi:2*pi,
	       z buffer=sort,
		   shader=faceted,]
			({0.5*sqrt(x^2+1)*cos(deg(y))},
			 {0.5*sqrt(x^2+1)*sin(deg(y))},
			{x});
	    \end{axis}
\end{tikzpicture}
\begin{tikzpicture}[xscale=0.6, yscale=0.8]
	    \begin{axis}[view={40}{20}]
	       \addplot3[surf,
	       colormap/cool,
	       samples=60,
	       domain=1:2*pi,y domain=-2*pi:2*pi,
	       z buffer=sort,
		   shader=faceted,]
			({0.5*sqrt(x^2-1)*cos(deg(y))},
			 {0.5*sqrt(x^2-1)*sin(deg(y))},
			{x});
			\addplot3[surf,
	       colormap/cool,
	       samples=60,
	       domain=-2*pi:-1,y domain=-2*pi:2*pi,
	       z buffer=sort,
		   shader=faceted,]
			({0.5*sqrt(x^2-1)*cos(deg(y))},
			 {0.5*sqrt(x^2-1)*sin(deg(y))},
			{x});
	    \end{axis}
\end{tikzpicture}
\begin{tikzpicture}[xscale=0.6, yscale=0.8]
	    \begin{axis}[view={40}{20}]
	       \addplot3[surf,
	       colormap/cool,
	       samples=60,
	       domain=-2*pi:2*pi,y domain=-2*pi:2*pi,
	       z buffer=sort,
		   shader=faceted,]
			({0.5*sqrt(x^2)*cos(deg(y))},
			 {0.5*sqrt(x^2)*sin(deg(y))},
			{x});
	    \end{axis}
\end{tikzpicture}
\caption{hyperbolischer Raum für $c>0$, $c<0$, $c=0$}}
\end{figure}
	Sei $f:\mathds{R}^{n+1} \to \mathds{R}$ definiert durch $f_c(x) = -x_1^2 + \sum_{i=2}^{n+1}x_i^2 - c$.\\
	$\Rightarrow$ $f$ ist glatt und $f_c^{-1}(0) = H_c^n$\\
	Es bleibt zu zeigen, das $f$ eine Submersion ist.
	\[
		\mathd_x f = (-2x_1,2x_2, \ldots,2x_{n+1}): \mathds{R}^{n+1} \to \mathds{R}	
	\]
	Es gilt $\mathd_x f(\nu) = 0 \forall \nu \in \mathds{R}^{n+1} \Leftrightarrow x_i = 0\ \forall i=1, \ldots,n+1$.\\
Das heißt $\exists i: x_i \neq 0 \Rightarrow \mathd_x f\ \text{Submersion}$, wobei $\exists i: x_i \neq 0$ gilt, falls $x \in H_c^n, c \neq 0$. Also ist für $c \neq 0$ $H_c^n$ eine $n$-dimensionale Untermannigfaltigkeit. Für $c = 0$ ist $H_0^n \backslash \set{0}$ eine $n$-dimensionale Untermannigfaltigkeit.
	\item Der $n$-Torus $T^n = \set[x=(x_1,x_2, \ldots,x_{2n-1},x_{2n})\in \mathds{R}^{2n}]{x_1^2 + x_2^2 = 1, \ldots, x_{2n-1}^2 + x_{2n}^2 = 1}$
	\[
		f:\mathds{R}^{2n} \to \mathds{R},\ f(x) = 
		\begin{pmatrix}
			x_1^2 + x_2^2 - 1\\
			\vdots\\
			x_{2n-1}^2 + x_{2n}^2 - 1		
		\end{pmatrix}
	\]
	$\Rightarrow$ $f$ glatt, $f^{-1}(0)=T^n$
	\[
		\mathd_x f =
		\begin{pmatrix}
			2x_1 & 2x_2 & 0 & 0 &\ldots & 0\\
			0 & 0 & 2x_3 &2x_4 &  &0\\
			\vdots &&& \ddots &\ddots & \\
			0 &&&& 2x_{2n-1} & 2x_{2n}
		\end{pmatrix}
	\]
	$\Rightarrow$ $\mathd_x f$ ist eine Submersion.
	\item Die orthogonale Gruppe $O(n) \subset \mathds{R}^{n \times n}$ ist eine Untermannigfaltigkeit der Dimension $\frac{n(n-1)}{2}$.
	\[
		O(n) = \set[A \in \mathds{R}^{n \times n}]{AA^T = \id_{\mathds{R}^n} = f^{-1}(0)},\ f(A) = AA^T - \id_{\mathds{R}^n}
	\]
	Es gilt $f(A)^T = {A^T}^T \cdot A^T - \id^T = AA^T - \id = f(A)$. Also $f:\mathds{R}^{n \times n} \to \text{Sym}(\mathds{R}^{n \times n}) = \set[A \subseteq \mathds{R}^{n \times n}]{A^T = A}$.\\
	$\Rightarrow$ $f$ glatt, $f^{-1}(0) = O(n)$
	\[
		\mathd_A f:\mathds{R}^{n \times n} \to \text{Sym}(\mathds{R}^{n \times n}),\ \mathd_A f = HA^T+AH^T
	\]
	Sei $S \in \text{Sym}(\mathds{R}^{n \times})$. Setze $H = \frac{SA}{2} \Rightarrow \mathd_A f(\frac{SA}{2}) = \frac{1}{2}(SAA^T+ASA^T) = \frac{1}{2}(S +  S) = S$\\
	$\Rightarrow$ $O(n)$ ist eine Untermannigfaltigkeit
\end{enumerate}
% subsubsection 17 (end)

\subsubsection[Bemerkung: ]{Bemerkung} % (fold)
\label{ssub:18}
$M$ kann von der Form $M = f^{-1}(0)$ sein, $f$ glatt, aber keine Submersion. Zum Beispiel
	\[
		M = \set[(x,x)\in \mathds{R}^2]{x \in \mathds{R}} = f^{-1}(0),\ f(x,y) = x^3 +y^3	
	\]
%subsubsection18 (end)
	
\subsubsection[Satz: lokale Parametrisierung von Untermannigfaltigkeiten]{Satz}
\label{ssub:19}
Es sind äquivalent:
\begin{enumerate}[(i)]
	\item $M^n\subset \mathds{R}^{n+k}$ ist $n$-dim Untermannigfaltigkeit
	\item $\forall x \in M^n\ \exists$ Umgebungen $U$ und $V \subset \mathds{R}^{n+k}$ von $x$ und $0$, und ein Diffeomorphismus $\phi: U \to V$ mit $\phi(M\cap U) \to \underbrace{(\mathds{R}^n \times \set{0}\cap V)}_{=: W}$
	\item $\forall x \in M^n\ \exists$ Umgebungen $U\subset \mathds{R}^{n+k}$ und $W\subset \mathds{R}^n$, sowie eine glatte Abbildung $g:W \to \mathds{R}^{n+k}$, sodass $g$ ein Homöomorphismus von $W$ auf $M\cap U$ und eine Immersion ist. $g$ nennt man eine \Index{lokale Parametrisierung}
\end{enumerate}
\begin{figure}[H]
\centering{

\begin{tikzpicture}[xscale=1,yscale=0.7]
	%\draw[help lines] (0,0) grid (10,9);
	\draw[<->] (1,9) node[left]{$\mathds{R}^k$} -- (1,5) node[left]{$\mathds{R}^{n+k}$} -- (5,5) node[below]{$\mathds{R}^n$};
	\draw[<->] (1,4) -- (1,1) -- (5,1) node[below]{$\mathds{R}^n\times \set{0}$};
	\draw (9,5) node[below]{$\mathds{R}^k$} -- (9,5) -- (9,9);
	\draw[fill=OrangeRed2] (9,7) node[right] {$0$} circle[radius=0.05cm];

	
	\draw (1.5,5.5) .. controls +(75:0.5) and +(220:0.25) .. (2.25,6.5)
	 		.. controls +(40:0.25) and +(205:0.25) .. (3,7)
		 	.. controls +(25:1) and +(255:0.5) .. (4.5,8.5) node[right]{$M^n$};
	\draw[SteelBlue2] (1.89,6.16) .. controls +(47.5:0.4) and +(215:0.25) .. (2.65,6.8) node[right] {$U\cap M$};
	\draw[fill=OrangeRed2] (2.25,6.5) circle[radius=0.05cm];
	\draw[SeaGreen4,dashed] (2.25,6.5) circle[radius=0.5cm];
	\node[SeaGreen4,above] at (2.25,7) {$U$};
	
	\draw[->] (5.5,7) .. controls +(30:0.5) and +(180:0.5) .. (7,7.4) node[above] {$f$}
			 .. controls +(0:0.5) and +(150:0.5) .. (8.5,7);
			 
	\draw[SteelBlue2] (2,1) -- (3,1) node[below right] {$W$};	 
	\draw[fill=OrangeRed2] (2.5,1) circle[radius=0.05cm];
	\draw[SeaGreen4,dashed] (2.5,1) circle[radius=0.5cm];
	\node[SeaGreen4,above] at (2.5,1.5) {$V$};
	
	\draw[->] (2.25,1.5) .. controls +(110:1) and +(270:1) .. (1.75,3.75) node[right] {$\varphi$ Satz 5}
			 .. controls +(90:1) and +(250:1) .. (2.25,5.9);
			 
	\draw[->] (4.5,2.5) -- (8.5,6);
\end{tikzpicture}
\caption{Beweis von \ref{ssub:19}}}
\end{figure}
%subsubsection 19 (end)

\subsubsection[Beispiele zur lokalen Parametrisierung]{Beispiele}
\label{ssub:110}
\begin{enumerate}[(i)]
	\item Der $2$-Torus $T^2$ ist das Bild von
	\[
		g:\mathds{R}^2 \to \mathds{R}^4,\ g(x,y) = \enbrace[\big]{\cos(x),\sin(x),\cos(y),\sin(y)}\subset T^2 \subset \mathds{R}^4
	\]
	$g$ ist lokale Parametrisierung:
	\[
		\mathd_{(x,y)} g = 
		\begin{pmatrix}
		-\sin(x) & 0\\
		\cos(x) & 0\\
		0 & -\sin(x)\\
		0 & \cos(y)
		\end{pmatrix}
	\]
	$\mathd_{(x,y)} g$ hat somit immer vollen Rang 2, also ist $\mathd_{(x,y)} g$ injektiv. Es folgt somit, dass $g$ eine Immersion ist.
	\begin{figure}[H]
		\centering
		\begin{minipage}{0.4\textwidth}
		\begin{tikzcd}[column sep=tiny,row sep=small]
			\mathds{R}^2 \arrow{rr} \arrow{dr}{p} && T^2\subset \mathds{R}^4\\
			& \enbrace{\nicefrac{\mathds{R}}{2\pi \mathds{Z}}}^2 \simeq S^1\times S^1 \arrow{ur}{\hat{g}} &
		\end{tikzcd}
		\caption{Diagramm zum lok. Homöomorphismus von g}
		\end{minipage}
		\hfill
		\color{SeaGreen4}\fbox{
		\begin{minipage}{0.49\textwidth}
			Sei $X$ ein topologischer Raum , $\sim$ eine Äquivalenzrelation. So ist $\nicefrac{X}{\sim}$ ein topologischer Raum versehen mit der Quotiententopologie, also $U\subset\ \nicefrac{X}{\sim}$ offen $\Leftrightarrow\ p^{-1}(U)$ offen. Somit folgt für $X \overset{p}{\to} \nicefrac{X}{\sim}$, dass $p$ stetig und auch ein lokaler Hoöomorphismus ist.
		\end{minipage}}	
	\end{figure}
	$\hat{g}$ ist stetig (Quotiententopologie), $\hat{g}$ bijektiv und $T^2$ kompakt. Damit folgt, dass $\hat{g}$ ein Homöomorphimus und somit $g$ ein lokaler Homöomorphismus ist.
	\item 
	Betrachte
	\[
		g:\mathds{R}^2 \to \mathds{R}^3\text{, } g(x,y) = \enbrace[\big]{\sin(x)\cos(y),\sin(x)\sin(y),\cos(x)} \in S^2\subset \mathds{R}^3
	\]
	$g$ ist ein lokaler Homöomorphismus, jedoch ist $\mathd_{(x,y)}$ nicht injektiv, das heißt keine lokale Parametrisierung nahe der Pole. Also müssen wir eine andere Parametrisierung wählen, zum Beispiel $g(x,y) = \enbrace{x,y,\sqrt{1-x^2-y^2}}$
\end{enumerate}
%subsubsection 110 (end)

\subsubsection[Bemerkung]{Bemerkung}
\label{ssub:111}
\begin{enumerate}[(i)]
	\item 
		Betrachte:
		\begin{figure}[H]
		\centering{
			\begin{tikzpicture}[scale = 0.6]
				%\draw[help lines] (0,0) grid (12,4);
				\draw (1,2) node[below]{$\infty$} -- (3,2) node[below]{$\mathds{R}$} -- (5,2) node[below]{$-\infty$};
				\draw[color = blue, thick, ->] (2,2) -- (3,2) node[above]{$\nu,\ \mathd_x g \neq 0$} -- (4,2);
				\draw[->] (5.5,2) .. controls +(25:0.125) and +(180:0.125) .. (6,2.125) node[above]{$g$} .. controls +(0:0.125) and +(155:0.125) .. (6.5,2);
				\draw[<->] (7,4) -- (7,0) -- (11,0);
				\draw[<->] (9,2) .. controls +(90:1) and +(180:0.5) .. (10,3.5)
				.. controls +(0:0.5) and +(90:0.5) .. (10.75,2.75)
				.. controls +(270:0.5) and +(0:1) .. (9.5,2)
				.. controls +(180:0.5) and +(0:0.5) .. (8.5,2)
				.. controls +(180:1) and +(90:0.5) .. (7.25,1.25)
				.. controls +(270:0.5) and +(180:0.5) .. (8,0.5)
				.. controls +(0:0.5) and +(270:1) .. (9,2);
				\draw[color = blue, thick, ->] (9.5,2) .. controls +(0:1) and +(270:0.5) .. (10.75,2.75) node[right]{$\nu$};
			\end{tikzpicture}
			}		
		\end{figure}
		$g$ ist immersiv, hat aber keine Umkehrfunktion und ist somit kein Homöomophismus.
		\item
		Betrachte $C = \set[(t^2,t^3)\in \mathds{R}^2]{t\in \mathds{R}}$
		\begin{figure}[H]
		\centering
			\begin{tikzpicture}[yscale = 0.4, xscale= 0.6]
				%\draw[help lines] (0,0) grid (12,8);
				\draw (1,1) -- (1,7);
				\draw (1,4) -- (7,4);
				\draw (1,4)  .. controls +(3:3) and +(225:1) .. (6,7); 
				\draw (1,4)  .. controls +(357:3) and +(135:1) .. (6,1);
	
				\draw (10,1) node[right]{$\mathds{R}$} -- (10,7);
				\draw[->] (9.5,4) -- (8.5,4) node[above]{$g$} -- (7.5,4);
		
				\draw[SteelBlue2,->] (0.1,4.9) node[above]{$\mathd_0 g = 0$} -- (0.9,4.1);
			\end{tikzpicture}
		\end{figure}
		$g$ ist keine lokale Parametrisierung, da $g$ nicht immersiv in $0$ ist. $g(t) = \enbrace[\big]{g_1(t),g_2(t)}\in \mathds{R}^2 \Rightarrow g_1^2(t) = g_2^3(t) \Rightarrow \mathd_0 g = 0$
\end{enumerate}
%subsubsection 111 (end)

\subsubsection[Satz: Wechsel lok. Parametrisierungen]{Satz}
\label{ssub:112}
Sei $M^n$ eine $n$-dim glatte Untermannigfaltigkeit, $x\in M$, und seien $U_1,U_2$ Umgebungen von $x$ in $\mathds{R}^{k+n}$ mit lok. Parametrisierung $g_1$ und $g_2$ von $U_1\cap M$ und $U_2\cap M$ definiert auf $\Omega_1$ und $\Omega_2$ (das heißt $g_i: \Omega_i \to \mathds{R}^{n+k}$). Dann ist
\[
	g_2^{-1} \circ g_1:\Omega_1 \cap g_1^{-1}(U_2) \overset{g_1}{\to} g_1(\Omega_1) \cap U_2 \overset{g_2^{-1}}{\to} \Omega_2 \cap g_2^{-1}(U_1)
\]
ein $C^{\infty}$-Diffeomorphismus.
\begin{figure}[H]
\centering{
\begin{tikzpicture}[scale=0.6]
	%\draw[help lines] (0,0) grid (12,8);
	
	\draw (1.5,5) .. controls +(140:1) and +(200:1) .. (2,7.5) node[above]{$\Omega_1$}
	 .. controls +(20:1) and +(135:1) .. (4,7)
	 .. controls +(315:1) and +(40:1) .. (3.5,5)
	 .. controls +(220:1) and +(320:1) .. (1.5,5) node[left]{$\Omega_1\cap g_1^{-1}(U_2)$};
	 
	\draw[fill=SteelBlue2,fill opacity=0.2] (3.5,5) .. controls +(220:1) and +(320:1) .. (1.5,5) .. controls +(20:0.5) and +(160:0.5) .. (3.5,5);

	\draw (1.5,1) node[left]{$\Omega_2$} .. controls +(150:1) and +(210:1) .. (2,3.5) node[left]{$\Omega_2\cap g_2^{-1}(U_1)$}
	 .. controls +(30:0.7) and +(145:0.7) .. (3.5,3.5)
	 .. controls +(325:1) and +(30:1) .. (3.5,1)
	 .. controls +(210:1) and +(330:1) .. (1.5,1);
	 
	\draw[fill=SteelBlue2,fill opacity=0.2] (2,3.5) .. controls +(30:0.7) and +(145:0.7) .. (3.5,3.5) .. controls +(200:0.5) and +(340:0.5) .. (2,3.5);
	
	\draw[->] (2.4,4.5) .. controls +(250:0.125) and +(90:0.125) .. (2.3,4.1) node[right]{$g_2^{-1}\circ g_1$} .. controls +(270:0.125) and +(110:0.125) .. (2.4,3.7);
	
	\draw (8.5,2) .. controls +(150:1) and +(210:1) .. (9,4.5)
	 .. controls +(30:0.7) and +(145:0.7) .. (10.5,4.5)
	 .. controls +(325:1) and +(30:1) .. (10.5,2) node[below right]{$U_2\cap M$}
	 .. controls +(210:1) and +(330:1) .. (8.5,2);
	 
	\draw (9,4.5) .. controls +(150:1) and +(210:1) .. (9,7)
	 .. controls +(30:0.7) and +(145:0.7) .. (10.5,7) node[right]{$U_1\cap M$}
	 .. controls +(325:1) and +(30:1) .. (10.5,4.5)
	 .. controls +(210:0.7) and +(330:0.7) .. (9,4.5);
	 
	\draw[fill=SteelBlue2,fill opacity=0.2] (9,4.5) .. controls +(30:0.7) and +(145:0.7) .. (10.5,4.5) .. controls +(210:0.7) and +(330:0.7) .. (9,4.5); 

	\draw[->] (5,6) -- (6.25,6) node[above]{$g_1$} -- (7.5,6);
	\draw[->] (5,3) -- (6.25,3) node[above]{$g_2$} -- (7.5,3);
\end{tikzpicture}
\caption{Veranschaulichung des Wechsel lok. Parametrisierung}
}
\end{figure}
\minisec{Bemerkung}
	\begin{itemize}
		\item $h:\Omega_1 \to \Omega_2$ Homöomorphismus, offene Menge in $\Omega_1$ $\leftrightarrow$ offene Menge in $\Omega_2$
		\item $\varphi:\Omega_1 \to \Omega_2$ Diffeomorphismus, falls Homöom. und $C^{\infty}(\Omega_1)$ $\leftrightarrow$ $C^{\infty}(\Omega_2)$
	\end{itemize}
%subsubsection 112 (end)
% subsection 11 (end)

\subsection{Abstrakte Mannigfaltigkeit}
\label{sub:12}

\subsubsection[Definition: Karte und Atlas]{Definition}
\label{ssub:113}
Sei $M$ ein topologischer Raum. Ein \Index{Atlas} ist eine Familie $\set{(U_i,\varphi_i)}_i\in I$ mit
\begin{enumerate}[(i)]
	\item $U_i \subset M$ offen, $M = \bigcup_{i \in I} U_i$ ($\set{U_i}$ offene Überdeckung)
	\item $\set{\varphi_i:U_i \to \Omega_i \subset \mathds{R}^n}$ ist eine Familie von Homöomorphismen mit:
	\[
		U_{ij} = U_i \cap U_j \neq \emptyset \text{, dann ist } \varphi_{ij} = \varphi_i \circ \varphi_j^{-1}: \varphi_j(U_{ij}) \to \varphi_i(U_{ij})
	\]
	ein $C^{\infty}$-Diffeomorphismus.
\end{enumerate}
Ein Paar $(U_i,\varphi_i)$ heißt \Index{Karte} und $\varphi_{ij}$ heißt Übergangsfunktion oder \index{Kartenwechsel}.
%subsubsection 113(end)

\subsubsection[Definition: Abstrakte differenzierbare Mannigfaltigkeit]{Definiton}
\label{ssub:114}
Sei $M$ ein topologischer Raum.
\begin{enumerate}[(i)]
	\item Zwei Atlanten $\mathcal{A}_1 = \set{(U_i,\varphi_i)}_{i\in I}$ und $\mathcal{A}_2 = \set{(V_j,\psi_j)}_{j\in J}$ sind äquivalent, genau dann wenn $\forall i\in I,j\in J$ mit $U_i\cap V_j\neq \emptyset$
	\[
		\varphi_i\circ \psi_j^{-1}:\psi(U_i\cap V_j) \to \varphi(U_i\cap V_j)
	\]
	ein $C^{\infty}$-Diffeomorphismus, das heißt die Vereinigung und $\mathcal{A}_1 \cup \mathcal{A}_2$ ist wieder ein Atlas von $M$.
	\item \Index{Diffenrenzierbare Struktur} auf $M$ ist die Äquivalenzklasse eines Atlanten. Ein Hausdorffraum $M$ mit einer differenzierbaren Struktur heißt \Index{(abstrakte) differenzierbare Mannigfaltigkeit} oder Mannigfaltigkeit.
	\item Nehmen alle Karten eines Atlanten Werte in $\mathds{R}^n$ an, so heißt $n$ die Dimension der Mannigfaltigkeit.
\end{enumerate}
% subsubsection 114(end)

\subsubsection[Beispiele für differenzierbare Mannigfaltigkeiten]{Beispiele}	
\label{ssub:115}
\begin{enumerate}[(i)]
	\item $\mathds{R}^n$ ist Mannigfaltigkeit der Dimension $n$, Atlas: $\set{\mathds{R}^n,\id_{\mathds{R}^n}}$ und $U\subset \mathds{R}^n$ ist Mannigfaltigkeit der Dimension $n$, Atlas: $\set{U,\id_{\mathds{R}^n}\big\vert_U}$
	\item Untermannigfaltigkeit im Sinne von Abschnitt $1.1$ sind abstrakte Mannigfaltigkeiten.
	\item Mannigfaltigkeit $N\subset M$ ist eine Untermannigfaltigkeit von $M$, falls für alle $y\in N$ eine Karte $(U,\varphi)$ von M um $y$ existiert, sodass $\varphi(U\cap N)$ Untermannigfaltigkeit von $\mathds{R}^n$ im Sinne von Abschnitt $1.1$ ist.\\
Dann gilt zum Beispiel $U \overset{offen}{\subset} \mathds{R}^n$ ist Untermannigfaltigkeit im "abstrakten" Sinne. Untermannigfaltigkeit $N\subset \mathds{R}^n$ im Sinne von $1.1$ ist Untermannigfaltigkeit im abstrakten Sinne.
	\item Die $n$-Sphäre, z.B $n = 1$
	\begin{figure}[H]
		\tikzsetnextfilename{stereographic_projection}
		\centering{\begin{tikzpicture}[scale=0.8]
			\draw[name path=real] (-5.2,0) node[above right]{$\mathds{R}^n$} -- (5.2,0);
			\draw[name path=sphere, thick] (0,0) circle[radius=2cm];
			
			%einzelner Punkt mit Beschriftung
			\draw[name path=g1,lightgray, thick, dashed] (0,2) -- (-4,0);
			\path[name intersections={of=g1 and real, by=in1}];
			\path[name intersections={of=g1 and sphere}];
			\draw[fill=SteelBlue2] (in1) circle[radius=0.06cm] node[below, color=SteelBlue3]{$h^+(p)$};
			\draw[fill=SteelBlue2] (intersection-2) circle[radius=0.06cm] node[above left, color=SteelBlue3]{$p$};
			
			\begin{scope}
			\foreach \x in {-3, -2.2, 2, 2.5, 3.4, 5} %
			{	
				\draw[name path=g1,lightgray, thick, dashed] (0,2) -- (\x,0);
				\path[name intersections={of=g1 and real, by=in1}];
				\path[name intersections={of=g1 and sphere}];
				\draw[fill=SteelBlue2] (in1) circle[radius=0.06cm];
				\draw[fill=SteelBlue2] (intersection-1) circle[radius=0.06cm];
				\draw[fill=SteelBlue2] (intersection-2) circle[radius=0.06cm];
			}
			\end{scope}
			\foreach \x in {-3.1,-2.1, -0.7, 0.2, 0.7, 1.8, 2.9} %
			{	
				\path[name path=g1] (0,2) -- (\x,-2);
				\path[name intersections={of=g1 and real, by=in1}];
				\path[name intersections={of=g1 and sphere}];
				\draw[name path=g1,lightgray, thick, dashed] (0,2) -- (intersection-2);
				\draw[fill=SteelBlue2] (in1) circle[radius=0.06cm];
				\draw[fill=SteelBlue2] (intersection-2) circle[radius=0.06cm];
			}
			\draw[fill=OrangeRed2] (0,2) circle[radius=0.07cm] node[above]{$N$};
			\draw[fill=OrangeRed2] (0,-2) circle[radius=0.07cm] node[below]{$S$};
		\end{tikzpicture}
		\caption{Stereografische Projektion}}
	\end{figure}
		Allgemein: $U_N = S^n\backslash \set{N},\ U_S = S^n\backslash \set{S}\subset \mathds{R}^{n+1}$
		\begin{align*}
			\varphi_N:\ U_N \to \mathds{R}^n,\ \enbrace{x_1, \ldots , x_{n+1}} \mapsto \frac{1}{1-x_{n+1}} \enbrace{x_1, \ldots, x_n}\\
			\varphi_S:\ U_S \to \mathds{R}^n,\ \enbrace{x_1, \ldots , x_{n+1}} \mapsto \frac{1}{1+x_{n+1}} \enbrace{x_1, \ldots, x_n}\\
			\varphi_S \circ \varphi_N^{-1}:\ \mathds{R}^n\backslash \set{0} \to \mathds{R}^n\backslash \set{0},\ x\in \mathds{R}^n\backslash \set{0} \mapsto \frac{x}{\|x\|^2}\\
		\end{align*}
		\minisec{Bemerkung} $\set{\enbrace*{S^n,\id_{\mathds{R}^{n+1}}\vert_{S^n}}}$ ist kein Atlas, denn $S^n$ ist nicht offen im $\mathds{R}^{n+1}$.
		\item Der $n$-Torus $T^n$. Betrachte $f:\mathds{R}^n \to \mathds{C}^n,\ \enbrace{x_1, \ldots, x_n} \mapsto \enbrace{\exp(ix_1), \ldots, \exp(ix_n)}\in \underbrace{S^1\times \ldots \times S^1}_{n-\text{mal}} \subset \mathds{C}\times \ldots \times \mathds{C}$, so folgt $f(\mathds{R}^n) = T^n,\ f(a_1,\cdot,a_n) = (e^{ia_1},\cdot,e^{ia_n})\in T^n$.
\begin{figure}[H]
	\centering{
	\begin{tikzpicture}[yscale=0.6,xscale=0.8]
		\begin{axis}[
			z post scale=3,
			xmin=-1.5, xmax=3,
			zmin=-22,
			view={15}{18},
			hide axis,
			]
		%\draw[dashed, gray] (axis cs: 1,0,-19) -- (axis cs: 1,0,28);
		\addplot3[
			domain=0:2*pi,
			smooth,
			samples=50,
			samples y=1,
			ultra thick,
			postaction={decorate},
			decoration={markings,
				mark=at position 0.25 with {\arrow{stealth}},
				mark=at position 0.75 with {\arrow{stealth}}}]
			({cos(deg(x))}, {sin(deg(x))}, -17);
		\addplot3[
			domain=-0.2*pi:0.2*pi,
			smooth,
			samples=50,
			samples y=1,
			ultra thick,
			Firebrick2,
			dashed,
			-stealth,
			]
			({cos(deg(x))}, {sin(deg(x))}, -17);
		\node[pin=0:$e^{ia}$, minimum size=0.09cm, draw=black, thick, fill=Chartreuse3, inner sep=0.008cm, circle] at (axis cs: 1,0,-17){};
		\addplot3[
			domain=-1*pi:9*pi,
			smooth,
			samples=100,
			samples y=1,
			ultra thick,
			postaction={decorate},
			decoration={markings,
			mark=between positions 0.05 and 0.99 step 0.1 with {\arrow{stealth}}},%
			]
			({cos (deg(x) )},{sin(deg(x))}, {x} );
		\addplot3[
			domain=3.8*pi:4.2*pi,
			smooth,
			samples=50,
			samples y=1,
			ultra thick,
			dashed,
			-stealth,
			Firebrick2,
			]
			({cos (deg(x) )},{sin(deg(x))}, {x} );
		\draw[-to, very thick] (axis cs:0,0,-8) -- ++(axis direction cs: 0,0,-3.5) node[midway,left]{$\exp$};
		\foreach \t in {-2,0,2} {\edef\temp{\noexpand
			\node[pin=0:a + $\t \pi$, minimum size=0.11cm,draw, thick,circle, fill=Chartreuse3, inner sep=0.01cm] at (axis cs: 1,0,\t*pi+4*pi){}; 
		} \temp}
		\node at (axis cs: 0,0,31) {$\mathds{R}$};
		\end{axis}
	\end{tikzpicture}
	\caption{Veranschaulichung der Wahl der Karten}}
\end{figure}
Schränke $f\vert_{\enbrace{a_1-\epsilon,a_1+\epsilon}\times \ldots \times \enbrace{a_n-\epsilon,a_n+\epsilon}}$ auf sein Bild ein. $\varphi_a := f^{-1} : \underset{= U_a \subset T}{f(I_a)} \to I_a \overset{\text{offen}}{\subset} \mathds{R}^n$. Somit ist $\set{\enbrace{U_a,\varphi_a}}_{a\in \mathds{R}^n}$ eine Familie von Karten, die $T^n$ überdeckt. Es bleibt zu zeigen, dass $\varphi_{ab} = \varphi_a \circ \varphi_b$ ein Diffeomorphismus ist, was wir an dieser Stelle dem Leser überlassen. Somit folgt, dass $\set{\enbrace{U_a,\varphi_a}}_{a\in \mathds{R}^n}$ ein Atlas ist.
		\item Betrachte den Hyperbolischen Raum $H^n$.
		\missingfigure{Veranschaulichung der Wahl der Karten}
		$H^n$ ist eine Mannigfaltigkeit mit einer Karte, gegeben durch stereographische Projektion. Man sagt daher, dass $H^n$ diffeomorph zu $\mathds{R}^n$
		\item Würfel $Q\subset \mathds{R}^n$
		\begin{figure}[H]
			\centering{			
			\begin{tikzpicture}[scale = 0.6]
				%\draw[help lines] (0,0) grid (6,6);
				\draw[<->] (0,3) -- (6,3);
				\draw[<->] (3,0) -- (3,6);
				\draw (1,1) -- (5,1) -- (5,5) -- (1,5) -- (1,1);
				\draw (3,3) node[below right]{$S^n$} circle(2);
				\foreach \x/\y in {1/2,1/4,2/5,4/5,5/2,5/4,4/1,2/1}
				{
				\draw[dashed] (\x,\y) -- (3,3);
				}
				\node at (2,3.6) {$\sigma$};
			\end{tikzpicture}	
			\caption{Der Würfel $Q$ ist keine Untermannigfaltigkeit des $\mathds{R}^n$!}
			}
		\end{figure}
		$Q$ ist jedoch eine abstrakte Mannigfaltigkeit, denn $Q$ ist ein Hausdorffraum und bijektiv zu $S^n$ durch Projektion $sigma$. So ist $\mathcal{A} := \set{\enbrace*{\sigma^{-1}(U_N),\varphi_N \circ \sigma},\enbrace*{\sigma^{-1}(U_S),\varphi_S \circ \sigma}}$ ein Atlas, da $\varphi_2 \circ \varphi_1^{-1} = \varphi_S \circ \sigma \circ \sigma^{-1} \circ \varphi_N$ und somit ist $Q$ eine abstrakte Mannigfaltigkeit.
		\item Sei $\mathds{K} = \mathds{R},\ \mathds{C}$. So ist der \Index{projektive Raum}
		\begin{align*}
			\mathds{K}\mathds{P}^n =& \text{Menge der Ursprungsgraden in $\mathds{K}^{n+1}$}\\
			=& \nicefrac{\mathds{K}^{n+1}\backslash \set{0}}{\sim},\ x,y\in \mathds{K}^{n+1}\backslash \set{0}\ \text{seien äquivalent},\ x\sim y \Leftrightarrow \exists \lambda\in \mathds{K}\backslash \set{0}: x = \lambda y\\
			=& \nicefrac{\mathds{K}^{n+1}\backslash \set{0}}{\mathds{K}^*}\ \text{wobei}\\
			&[x]\in \nicefrac{\mathds{K}^{n+1}\backslash \set{0}}{\mathds{K}^*} \mapsto\ <x>_{\mathds{K}}\ \text{lineare Erzeugnis von $x$ in $\mathds{K}^{n+1}$}
		\end{align*}
		Wir schreiben $[x] = [x_0:\ldots:x_n]$, falls $x=(x_0,\ldots,x_n)\in \mathds{K}^{n+1}$, "homogene Koordinaten".
		\minisec{Bemerkung} $[x_0:\ldots:x_n] = [\lambda x_0:\ldots:\lambda x_n],\ \lambda \in \mathds{K}^{n+1}$
		\begin{enumerate}[(1)]
			\item Betrachte $p:\ x\in \mathds{K}^{n+1}\backslash \set{0} \mapsto [x] = p(x)\in \nicefrac{\mathds{K}^{n+1}\backslash \set{0}}{\sim}$ so folgt, dass $\mathds{K}\mathds{P}^n$ versehen mit der Quotiententopologie ein Hausdorff-Raum ist.
			\begin{figure}[H]
			\centering{
	\begin{tikzpicture}[scale=1]
		\draw[->] (0,0) -- (2,0) node[right]{$x_2$};
		\draw (0,0) -- (0,1);
		\draw[->] (0,0) -- (0,0,2) node[left]{$x_1$};
		\draw (0.25,-0.5,-0.5) -- (-0.5,1,1);
		\draw (-0.15,-0.5,-0.3) -- (0.3,1,0.6);
		\draw (-0.3,-0.5,0.35) -- (0.6,1,-0.7);
		\draw (0.55,-0.5,0.15) -- (-1.1,1,-0.3);
		%\draw ( ,-0.5, ) -- ( ,1, );
		\draw[SeaGreen4, thick, fill=SeaGreen4!20, fill opacity=0.4] (2,1,2) -- (2,1,-2) -- (-2,1,-2) -- (-2,1,2) -- cycle;
		\draw (-0.5,1,1) -- (-1.25,2.5,2.5);
		\draw (0.3,1,0.6) -- (0.75,2.5,1.5);
		\draw (0.6,1,-0.7) -- (1.5,2.5,-1.75);
		\draw (-1.1,1,-0.3) -- (-2.75,2.5,-0.75);
		\draw[fill, Firebrick2] (-0.5,1,1) circle[radius=0.03];
		\draw[fill, Firebrick2] (0.3,1,0.6) circle[radius=0.03];
		\draw[fill, Firebrick2] (0.6,1,-0.7) circle[radius=0.03];
		\draw[fill, Firebrick2] (-1.1,1,-0.3) circle[radius=0.03];
		\draw[->] (0,1) -- (0,2) node[left]{$x_3$};
		\draw[fill, SeaGreen4] (0,1) circle[radius=0.03];
		\node[below right,SeaGreen4] at (2,1,0) {$\mathds{R}^2\times\set{1}$} ;
		%\draw (0,1) -- (0,2);
	\end{tikzpicture}
	\caption{Der 2-dimensionale projektive Raum $\mathds{R}P^2$ (projektive Ebene)}
	\begin{tikzpicture}[scale=0.8]
		\draw[->, thick] (-1,0) -- (3,0) node[right]{$x_1$};
		\draw[->, thick] (0,-0.2) -- (0,1.5) node[above]{$x_2$};
		\draw[name path=ebene, thick, SeaGreen4] (-1.5,1) -- (3,1) node[below]{$\mathds{R}\times \set{1}$};
		\foreach \phi / \rad in {30 / 3, 55 / 2, 78 / 1.7, 109 / 1.7, 134 / 2}
		{
			\draw[name path=gerade] (\phi+180:0.3) -- (0,0) -- (\phi:\rad);
			\path[name intersections={of=ebene and gerade, by=in1}];
			\draw[fill=Firebrick2] (in1) circle[radius=0.05cm];
		}	
	\end{tikzpicture}
	\caption{Der eindimensionale projektive Raum $\mathds{R}P^1$}}
	\end{figure}
	Aus Abbildung 3 bzw. 2 folgt somit $\mathds{R}\mathds{P}^n = \mathds{R} \cup \set{\infty}  \cong S^2$ und analog $\mathds{R}^2 \cup \set{N} \cong S^2$. Um zu zeigen, dass $\mathds{K}\mathds{P}^n$ eine abstrakte Mannigfaltigkeit ist, muss die Existenz eines Atlanten gezeigt werden. Betrachte $U_i = \set[[x_0:\ldots:x_n]\in \mathds{K}\mathds{P}^n]{x_i \neq 0}$. So folgt, dass $\bigcup_{i=0}^n = \mathds{K}\mathds{P}^n$, da $(0,\ldots,0)\not \in \mathds{K}^{n+1}\backslash \set{0}$, also $\not \exists: [0:\ldots:0]$. Da $[x_0:\ldots:x_n]\in U_i \Rightarrow x_i \neq 0$, ist $\varphi_i: [x_0:\ldots:x_n]=[\frac{x_0}{x_i}:\ldots:\frac{x_n}{x_i}] \overset{\varphi_i}{\to} [\frac{x_0}{x_i}:\ldots:\hat{\frac{x_i}{x_i}}:\ldots\frac{x_n}{x_i}]$ wohldefiniert.
	Es bleibt als Übung zu zeigen, dass $\varphi_i \circ \varphi^{-1}_j$ ein Diffeomorphismus ist.
		\end{enumerate}
\end{enumerate}	
% subsubsection 115(end)

\subsubsection[Bemerkungen zu abstrakten Mannigfaltigkeiten]{Bemerkung}
\label{ssub:116}
\begin{enumerate}[(i)]
	\item Die Vereinigung aller äquivalenten Atlanten ist wieder ein Atlas, der sogenannte maximale Atlas (eindeutig bestimmt). In der Praxis arbeitet man mit konkreten Atlanten.
	\item Allgemein kann man $C^p$-Atlanten betrachten, das heißt $\varphi_{ij}$ $C^p$-Diffeomorphismen $\to$ $C^p$-Mannigfaltigkeiten. Ebenso kann man Karten mit allgemeinen Zielräumen betrachten, wo der Begriff der Differenzierbarkeit/Diffeomorphismus Sinn hat. Zum Beispiel Banach- und Hilbertraum $\to$ Banach- und Hilbertmannigfaltigkeiten oder Frécheraum $\to$ Fréchemannigfaltigkeiten.
	\item Sei $M$ ein topologischer Raum mit differenzierbarer Struktur, so ist er nicht automatisch Hausdorff! Zum Beispiel
\begin{figure}[H]
\centering{
			\begin{tikzpicture}[yscale=0.6]
				%\draw[help lines] (0,0) grid (8,5);
				\draw (1,4) -- (7,4);
				\draw[thick] (4,3.9) -- (4,4.1) node[below right]{$0$};
				\filldraw[SteelBlue4] (4,4.5) circle[radius=0.04cm];
				\filldraw[OrangeRed2] (4,3.5) circle[radius=0.04cm];
				
				\draw[color = SteelBlue4] (1,1.5) -- (7,1.5) node[below right]{$\mathds{R}$};
				\draw[color = OrangeRed2] (1,1) -- (7,1) node[below right]{$\mathds{R}$};
				\draw[thick,->] (4,2) -- (4,2.5) node[right]{$\nicefrac{\mathds{R}\cup\mathds{R}}{\sim}$ mit $(x,0)\sim (0,y) \Leftrightarrow x = y \neq 0$} -- (4,3);
			\end{tikzpicture}
		\caption{Beispiel einer Nichthaudorffmenge}}
\end{figure}
Benutze Atlas $\set{\enbrace{U_1=\set[(x,0)]{x\in \mathds{R}},\id},\enbrace{U_2=\set[(0,y)]{y\in \mathds{R}},\id}}$. Dies ist ein Atlas, da $U_1\cap U_2 = M\backslash \set{:}$.
\item Wir nehmen künftig zwei zusätzliche Bedingungen für Mannigfaltigkeiten an:
	\begin{enumerate}[(a)]
		\item $M$ zusammenhängend $\to$ jede Mannigfaltigkeit hat wohldefinierte Dimension schreibe $M^n$
		\item $M$ erfüllt das 2. Abzählbarkeitsaxiom. Für uns: $\exists$ abzählbarer Atlas.
	\end{enumerate}
\end{enumerate}
%subsubsection 116(end)

% subsection 12(end)

\subsection{Glatte Abbildung}
\label{sub:13}

\subsubsection[Definition: Glatte Funktionen zwischen zwei Mannigfaltigkeiten]{Definition}
\label{ssub:117}

Seien $M^n$, $N^m$ zwei Mannigfaltigkeiten und sei $f:M^n\to N^m$ eine stetige Abbildung. Dann heißt $f$ glatt, falls für alle $x$ Karten $\enbrace{U,\varphi} \subset M$ und $\enbrace{V,\psi} \subset N$ um $x$ beziehungsweise $f^{-1}(x)$ existieren mit $f(U)\subset V$ und 
\[
	\psi \circ f \circ \varphi^{-1}:\ \underset{\subset \mathds{R}^n}{\varphi(U)} \to \underset{\subset \mathds{R}^n}{\psi(V)}
\]
glatt ist. ("$f$ in Karten gelesen ist glatt")
\missingfigure{Veranschaulichung der glatten Funktionen}
Wir schreiben $C^{\infty}(M,N) = \set[f:M\to N]{f\ \text{glatt}}$ und als Spezialfall $C^{\infty}(M) = C^{\infty}(M,\mathds{R})$.
% subsubsection 117(end)

\subsubsection{Bemerkung}
\label{ssub:118}
Glattheit ist unabhängig von der Wahl der Karten! ($\leadsto$ Diffeomorphismusbedingung!)
% subsubsection 118(end)

\subsubsection{Definition}
\label{ssub:119}
\begin{enumerate}[(i)]
\item Eine glatte Abbildung $f:M\to N$ heißt Submersion/Immersion/étale, falls für ein Paar (und damit alle Karten) $\psi \circ f \circ \varphi^{-1}$ eine Submersion/Immersion/étale ist.
\item $f$ ist Diffeomorphismus zwischen $M$ und $N$ ($M$ und $N$ sind "diffeomorph"), falls $f$ bijektiv ist und $f$ sowie $f^{-1}$ glatt ist.
\item $f:M\to N$ heißt Einbettung, falls $f$ Immersion und Homöomorphismus auf sein Bild ist.
\end{enumerate}
% subsubsection 119(end)

\subsubsection{Beispiele}
\label{ssub:120}
\begin{enumerate}[(i)]
\item Wenn $M$ eine Untermannigfaltigkeit des $\mathds{R}^{n+k}$, $f:\mathds{R}^{n+k}\to \mathds{R}^m$ glatte Funktion, dann ist auch $f\vert_M$ glatt.
\begin{itemize}
\item $\set{\mathds{R}^n,\id}$, $\set{\mathds{R}^m,\id}$ $\Rightarrow\ f:\mathds{R}^n\to \mathds{R}^m$ glatt im Sinne von Definition $\ref{ssub:117}$ $\Leftrightarrow$ glatt im üblichen Sinne
\item $x\in M$, Karte um $x$: $\varphi: U\subset \mathds{R}^{n+k} \underset{\text{Diffeo.}}{\to} V = \varphi(U)\subset \mathds{R}^{n+k}$ mit $\varphi(U\cap M) = U\cap (\mathds{R}^n\times \set{0})$ (aus $\ref{ssub:119}$), so gilt
\[
f\circ \varphi^{-1}: V\cap (\mathds{R}^n\times \set{0}) \xrightarrow{\varphi^{-1} \vert_{V\cap (\mathds{R}^n\times \set{0})}\in C^{\infty}(\mathds{R}^n,\mathds{R}^{n+k})} 
M\cap U \xrightarrow{f \in C^{\infty}(\mathds{R}^{n+k},\mathds{R}^m)} \mathds{R}^m
\]
Somit folgt, dass $f\vert_M$ glatt ist. 
\missingfigure{Veranschaulichung von f}
\end{itemize}
\item $S^2$ ist diffeomorph zu $\mathds{C}\mathds{P}^1 = \set{\text{Ursprungsgeraden in}\ \mathds{C}^2} = \nicefrac{\mathds{C}\backslash \set{0}}{\mathds{C}^*}$. Definiere 
\[
f: S^2\subset \mathds{R}^3 \to \mathds{C}\mathds{P}^1,\ \enbrace{x,y,z} \mapsto
\begin{cases}
z\neq -1:\ [\frac{x+iy}{1+z}:1]\\
z = -1:\ [1:0]
\end{cases}
\]
Somit folgt
\begin{enumerate}[(1)]
\item $f$ ist bijektiv
\item Wir "lesen nun in Karten"(, wobei $U_1 = \set[ z_1:z_2 ]{z_2\neq 0} \overset{\psi_2}{\to} \frac{z_1}{z_2}$:
\[
\psi \circ f \circ \varphi_S^{-1}(x,y) = \underset{\in U_1}{\psi_2 [x+iy:1]} = x + iy \underset{\mathds{C} \simeq \mathds{R}}{=} (x,y)
\]
Also ist $\psi \circ f \circ \varphi_S^{-1} = \id$ und somit $f \vert_{S^2\backslash\set{S}}$ ist $C^{\infty}$ und $f^{-1}$ auch $C^{\infty}$. Nun brauchen wie nur noch eine Karte um S!. Zum Beispiel $\varphi_N$. Hier $\varphi_N(S) = (0,0)$. Es gilt:
\[
	f\circ \varphi_N^{-1}(x,y) = 
	\begin{cases}
		(x,y) \neq (0,0): [\varphi_S \circ \varphi_N^{-1}(x,y):1 \overset{\mathds{C} \simeq \mathds{R}}{=} [\frac{x+iy}{x^2+y^2}:1] \in U_0]\\
	[1:0] \in U_0 
	\end{cases}
\]
Somit folgt:
\[
	\psi_1 \circ f \circ \varphi_N^{-1}(x,y) =
	\begin{cases}
		\frac{x^2+y^2}{x+iy} = x-iy\\
		0
	\end{cases}
	= (x,-y)\ \text{glatt}
\]
\end{enumerate}
\end{enumerate}
% subsubsection 120(end)

\subsubsection{Definition}
\label{ssub:121}
\begin{enumerate}[(i)]
	\item Eine Mannigfaltigkeit erfüllt das zweite Abzählbarkeitsaxiom, falls es einen abzählbaren Atlas besitzt.
	\item Sein $M^n$ eine Mannigfaltigkeit. Eine Familie $\set{f_i:M\to \mathds{R}}$ glatter Funktionen heißt \Index{Zerlegung der Eins}, falls
	\begin{itemize}
		\item $\supp f_i = \overline{\set[x\in M]{f_i(x) \neq 0}}$ "Träger von $f_i$" sind kompakt und bilden "lokal endliche Familie" (das heißt, sei $K\subset M$ kompakt, so gilt $\supp f_i \cap K = \emptyset$ außer für endlich viele $i\in I$)
		\item $\forall x \sum_{i\in I} f_i(x) = 1$ (wohldefiniert, also $\sum$ endlich, da lokal endlich)
	\end{itemize}
\end{enumerate}
% subsubsection 121(end)

\subsubsection{Theorem}
\label{ssub:122}
Sei $\set{U_i}_{i\in I}$ eine Überdeckung von $M$. Ist $M$ zweifach abzählbar, so existiert eine Zerlegung der Eins, die $\mathcal{U}$ untergeordnet ist, das heißt $\supp f_i \subset U_i$. Dies ist nützlich um beliebig definierte Objekte zu globalen Objekten zu verkleben.
\missingfigure{Untergeordnete Zerlegung der Eins}
% subsubsection 122(end)

\subsubsection{Whitney Theorem}
\label{ssub:123}
Sei $M^n$ eine zweifach abzählbare Mannigfaltigkeit. Dann existiert eine Einbettung $f:M^n\to \mathds{R}^{2n}$ von M (bzw. Immersion $f:M^n\to \mathds{R}^{2n-1}$) als Untermannigfaltigkeit von $\mathds{R}^{2n}$ (bzw. $\mathds{R}^{2n-1}$)
% subsubsection 123(end)

\subsubsection{Bemerkung}
\label{ssub:124}
\begin{enumerate}[(i)]
\item Können jede "abstrakte" Mannigfaltigkit $M^n$ als eine Untermannigfaltigkeit des $\mathds{R}^{2n-1}$ realisieren.
\minisec{Aber}
Die Einbettung ist nicht kanonisch (oder "natürlich"), das heißt gewisse Eigenschaften der Untermannigfaltigkeit entstehen durch die Einbettung und sind keine intrinsische Einbettung von $M$!
\item Die Dimension $2n-1$ des Bildraumes ist optimal im folgenden Sinne
\[
\exists\ \text{Mannigfaltigkeit}\ M^n,\ \text{die sich nicht in}\ \mathds{R}^N\ \text{mit}\ N<2n-1\ \text{einbetten lassen.}
\]
\end{enumerate}
\minisec{Beweisskizze}
Wir nehmen an, das die Mannigfaltigkeit $M$ kompakt ist. So folgt mit \ref{ssub:135}, das ein endlicher Atlas $\set{\enbrace{U_i,\varphi_i : U_i \to \Omega_i\subset \mathds{R}^n}}_{i=1}^n$ mit untergeordneter Zerlegung der Eins $\set{f_i:M\to \mathds{R}}_{i=1}^n,\ \supp f_i \subset U_i$ existiert. Definiere $\psi_i:M\to \mathds{R}^n$ durch $\psi_i(x) = \varphi_i(x)f_i(x)$, so folgt $\psi_i(x) \equiv 0$ für $x\in U_i^c$ und $\psi_i$ ist glatt. Dann definiert $\psi: M\to \mathds{R}^{nr+r},\ \psi(x)=\enbrace{\psi_1(x),\ldots,\psi_r(x),f_1(x),\ldots,f_r(x)}$ eine Einbettung.\\
Von jetzt an ist $M^n$ stets eine glatte zweifach abzählbare Mannigfaltigkeit.
% subsubsection 124(end)

% subsection 14(end)

\subsection{Das Tangentialbündel einer Untermannigfaltigkeit}
\label{sub:14}

\subsubsection{Definition}
\label{ssub:125}
Sei $M^n\subset \mathds{R}^{n+k}$ eine $n$-dimensionale Untermannigfaltigkeit. Sei $x\in M$. $v\in \mathds{R}^{n+k}$ ist tangential an $M$ in $x$, falls eine $C^{\infty}$-Abbildung $c: 0\in I\subset \mathds{R} \to \mathds{R}^{n+k}$ existiert mit
\begin{itemize}
\item $c(I)\subset M$
\item $c(0) = x$
\item $c'(t) = v$
\end{itemize}
\minisec{Notation}
$T_x M^n = $ Menge der Tangentialvektoren $v$ in $x$.
% subsubsection 125(end)

\subsubsection[Satz: Vektorraumstruktur des Tangentialraumes]{Satz}
\label{ssub:126}
$T_x M^n$ ist ein $n$-dim Vektorraum$\vert_{\mathds{R}}$ ($\cong \mathds{R}^n$)
\minisec{Beweis}
Sei $M$ eine Untermannigfaltigkeit, $x\in M$, so existiert Umgebung $U\subset \mathds{R}^{n+k}$ offen, $\varphi:U\to V\subset \mathds{R}^{n+k}$ diffeomorph mit $\varphi(U\cap M) = V\cap(\mathds{R}^n\times \set{0})$. Sei $c: I\to \mathds{R}^{n+k}$ ein Weg mit $c(I)\subset M\cap U$, so ist $\varphi \circ c$ ein Weg im $V\cap(\mathds{R}^n\times \set{0})$. Es folgt, dass $T_x M \cong \mathd_{\varphi(x)} \varphi^{-1}(\mathds{R}^n\times \set{0})$, das heißt $T_x M$ übernimmt die Vektorraumstruktur von $\mathds{R}^n \times \set{0}$ 
\missingfigure{Beispiel eines Tangentialraumes am $S^2$}
% subsubsection 126

\subsubsection{Satz}
\label{ssub:127}
Sei $M^n\subset \mathds{R}^{n+k}$ eine Untermannigfaltigkeit der Dimension $n$, $x\in M$ und $U\subset \mathds{R}^{n+k}$ offene Umgebung um $x$.
\begin{enumerate}[(i)]
\item $U\cap M = f^{-1}(0)$ für $f:\mathds{R}^{n+k} \to \mathds{R}^k$ eine Submersion, so ist $T_x M^n = \ker \mathd_x f$.\\
Sei zum Beispiel $M = S^2$. So ist $f(x,y,z)=x^2+y^2+z^2-1$ eine Submersion mit $S^2 = f^{-1}(0)$. Betrachte nun den Tangentialraum am Nordpol $N$. Es gilt $\mathd_N f = (0,0,z): \mathds{R}^3 \to \mathds{R}$, also gilt für den Tangentialraum: 
\[
T_N S^2 = \ker \mathd_N f = \set{
\begin{pmatrix}
x\\
y\\
0
\end{pmatrix}
\in \mathds{R}^3}
\]
\item Sei $\varphi: U\to V\subset \mathds{R}^{n+k}$ ein Diffeomorphismus mit $\varphi(U\cap M) = V\cap (\mathds{R}^n \times \set{0})$, dann ist $T_x M = \mathd_{\varphi(x)} \varphi^{-1}(\mathds{R}^n \times \set{0})$
\item Ist $\enbrace{\Omega,g}$ eine Parametrisierung von $U\cap M$, so ist der Tangentialraum $T_{g(y)} M = \mathd_y g(\mathds{R}^n)$
\end{enumerate} 
\minisec{Beweis}
\begin{enumerate}[(i)]
\item Sei $v\in T_x M$, das heißt $v = c^(0),\ c(0) = x$. So gilt $f\circ c(t) = 0$, da $c(t)\subset M\cap U = f^{-1}(0)$. Somit folgt mit der Kettenregel, das $\mathd_{c(0)} f \circ c'(0) = \mathd_x f(v) = 0$. Also ist $v\in \ker d_x f$ und es gilt $T_x M \subset \ker d_x f$. Da $T_x M$ und $\ker \mathd_x f$ $n$-dimensional, folgt somit die Gleichheit.
\item Siehe \ref{ssub:142}
\item Sei $\tilde{c}\subset \Omega$ eine Kurve, $c := g\circ \tilde{c}$ Kurve durch $c(0) = g\enbrace*{\tilde{c}(0)}\in M$. So ist $c'(0) = d_{\tilde{c}(0)} g(\tilde{c'(0)})\in T_{c(0)} M$. Also gilt $d_y g(\mathds{R}^n) \subset T_{g(y)} M$ und die Gleichheit folgt wie in (i) durch die Dimension.
\end{enumerate}
% subsubsection 127(end)
% subsection 14(end)


\subsection{Das Tangentialbündel einer abstrakten Mannigfaltigkeit}
\label{sub:15}

\subsubsection{Definition}
\label{ssub:128}
Sei $M^n$ eine abstrakte Mannigfaltigkeit, $x\in M$. Ein \Index{Tangentialvektor} von $M$ in $x$ ist eine Äquivalenzklasse von Kurven $\alpha : I\to M$ mit $\alpha(0) = x$. Dabei heißen $\alpha:I\to M$ und $\beta:J\to M$ äquivalent, falls eine Karte $(U,\varphi)$ um $x$ existiert mit
\[
\enbrace{\varphi \circ \alpha}'(0) = \enbrace{\varphi \circ \beta}'(0)
\]
\missingfigure{Veranschaulichung der Definition}
\minisec{Notation}
$T_x M =$ Raum der Tangentialvektoren
% subsubsection 128(end)

\subsubsection{Bemerkung}
\label{ssub:129}
\begin{enumerate}[(i)]
\item Gilt die oben genannte Äquivalenzbedingung für eine Karte, dann gilt sie für alle.
\minisec{Beweis}
Sei $\enbrace{V,\psi}$ eine andere Karte. Es gilt
\begin{align*}
\enbrace{\psi \circ \alpha}'(0) &= \enbrace{\psi \circ \varphi^{-1} \varphi \circ \alpha}'(0)\\
&= \mathd_{\varphi(\alpha(0))} (\psi \circ \varphi^{-1})\enbrace[\big]{(\varphi \circ \alpha)'(0)}\\
&= \mathd_{\varphi(\beta(0))} (\psi \circ \varphi^{-1})\enbrace[\big]{(\varphi \circ \beta)'(0)} = \enbrace{\psi \circ \beta}'(0)
\end{align*}
\item $T_x M$ ist $n$-dimensionaler Vektorraum$\vert_{\mathds{R}}$.
\minisec{Beweis}
Sei $\enbrace{U,\varphi}$ eine Karte um $x$. Sei $\Gamma_{(U,\varphi,x)}: T_x M\to \mathds{R}^n$, $\underset{\in T_x M}{[\alpha]} \mapsto (\varphi \circ \alpha)' (0)$. Aus der Äquivalenzrelation folgt, das $\Gamma$ wohldefiniert und bijektiv. Also können wir die Vektorraumstruktur auf $T_x M$ durch $\Gamma_{(U,\varphi,x)}$ erklären. Dies ist wohldefiniert, da $\mathd_{\varphi(x)} (\psi \circ \varphi^{-1})$ für jede weitere Karte $(V,\psi)$ um $x$ ein linearer Isomorphismus ist. (Da $\psi \circ \varphi^{-1}$ Diffeomorphismus.) 
\end{enumerate}
% subsubsection 129(end)

\subsubsection{Definition}
\label{ssub:130}
Sei $M$ eine Mannigfaltigkeit, $x\in M$, $(U,\varphi)$ und $(V,\psi)$ zwei Karten um $x$. Seien $u,v\in \mathds{R}^n$. Als Äquivalenzrelation ist definiert:
\[
u \sim v \xLeftrightarrow{\text{Def.}} v = \mathd_{\varphi(x)}(\psi \circ \varphi^{-1})(u)
\]
Ein \Index{Tangentialvektor} in $x$ auf $M$ ist eine solche Äquivalenzklasse $[(U,\varphi,u)]$
\missingfigure{}
% subsubsection 130(end)

\subsubsection{Bemerkung}
\label{ssub:131}
Sei $x\in M$, $(U,\varphi)$ Karte um $x$. So ist
\[
\Theta_{(U,\varphi,x)}: \underset{\in \varphi(x) + \mathds{R}^n}{u} \to [(U,\varphi,u)]\in T_x M
\]
eine Bijektion, die wohldefiniert $\text{VR}\vert_{\mathds{R}^n}$ auf $T_x M$ induziert.
% subsubsection 131(end)

\subsubsection{Definition}
\label{ssub:132}
$TM := \bigcup_{x\in M} T_x M$ ist das \Index{Tangentialbündel} von $M$.
% subsubsection 132(end)

\subsubsection{Theorem}
\label{ssub:133}
$TM^n$ trägt eine natürliche Differentialstruktur und hat die Dimension $2n$. (Beweis benutzt \ref{ssub:156})
\minisec{Beweis}
Definiere $\pi: TM\to M,\ \xi \in T_x M \mapsto \pi(\xi) = x\in M$ "Fußpunkt". $\pi$ ist surjektiv. Wir wollen nun \ref{ssub:157} anwenden. Für Karte $(U,\varphi)$ um $x\in M$ definiere $(\pi^{-1}(U),\Phi_{(U,\varphi)})$ durch 
\marginnote{$\Theta$ aus \ref{ssub:154}. Wieso genügt es nicht, $\Phi$ als $\Theta^{-1}$ zu erklären? Ich denke um zu gewährleisten, dass $\Phi$ bijektiv ist.}
\[
\Phi_{(U,\varphi)} (\xi) = \enbrace{\varphi \circ \pi(\xi),\Theta^{-1}_{(U,\varphi,\pi(\xi))}(\xi)}
\]
Es gilt $\bigcup_{(U,\varphi)} \pi^{-1}(U) = TM$ und $\Phi_{(U,\varphi)}(\pi^{-1}(U)) = \varphi(U)\times \mathds{R}^n \overset{\text{offen}}{\subset} \mathds{R}^n\times \mathds{R}^n$ ist eine Bijektion, da $\Theta$ bijektiv.\\
Verifiziere (i) und (ii):
\begin{enumerate}[(i)]
\item 
\[
\underbrace{\Phi_{(V,\psi)}^{-1}(\underset{\in \varphi(v)}{a},v)}_{T_{\psi^{-1}(a)} M} = \Theta_{(V,\psi,\psi ^{-1}(a))} (v) \xrightarrow{\Phi_{(V,\psi)}} \enbrace*{\psi\enbrace*{\psi^{-1}(a)}, \Theta^{-1}_{(V,\psi,\psi ^{-1}(a))}\enbrace*{\Theta_{(V,\psi,\psi ^{-1}(a))}(v)}}
\]
Berechne Übergangsfunktion:
\[
\Phi_{(U,\varphi)} \circ \Phi_{(V,\psi)}^{-1} (a,v) = \Phi_{(U,\varphi)}\enbrace*{\Theta_{(V,\psi,\psi ^{-1}(a))} (v)} = \enbrace{\varphi \circ \psi^{-1} (a), \mathd_{\phi(a)}(\varphi \circ \psi^{-1})(v)}
\]
(Beachte hierbei, das $T_{\varphi^{-1}(a)} \ni \big[(V,\psi,v)\big] = [(U,\varphi,\mathd_{\psi(a)}(\varphi \circ \psi^{-1})(v))]$)
\item Seien $\xi,\eta \in TM$. Falls $\pi(\xi) = \pi(\eta)$, so sind $\xi,\eta \in \pi^{-1}(U)$ für eine Karte um $x$. Falls $\pi(\xi) = x,\ \pi(\eta) = y,\ x\neq y$, so ist $\xi \in \pi^{-1}(U),\ \eta \in \pi^{-1}(V)$ für Karten $U,V$ von $x, y$ mit $U\cap V = \emptyset$, (diese existieren, da $M$ Hausdorff). Also $\pi^{-1}(U)\cap \pi^{-1}(V) = \emptyset$.
\end{enumerate}
% subsubsection 133(end)

\subsubsection{Lemma}
\label{ssub:134}
Sei $M$ eine Menge mit Überdeckung $\set{U_i\subset M}_{i\in I}$, das heißt $\bigcup_{i\in I} U_i = M$. Seien für alle $i\in I$ $\varphi_i:U_i\to \mathds{R}^n$ Bijektionen mit offene Mengen $\varphi(U_i)\subset \mathds{R}^n$. Falls
\begin{enumerate}[(i)]
\item $\forall\ i,j\in I\ \varphi_i \circ \varphi_j^{-1}$ ist glatt
\item $\forall\ x,y\in M$ gilt $\exists i\in I\ x,y\in U_i$ oder $\exists\ i,j$ und $U_i\cap U_j = \emptyset$ mit $x\in U_i$ und $y\in U_j$
\end{enumerate}
dann existiert eine eindeutig bestimmte Hausdorff-Topologie auf $M$ und eine differenzierbare Struktur, für die $\set{U_i,\varphi_i}$ ein Atlas ist.
\minisec{Beweis}
\begin{enumerate}[(i)]
\item $\leadsto$ Übergangsfunktion sind Diffeomorphismen
\item $\leadsto$ Hausdorff
\end{enumerate}
Siehe auch O'Neill.
% subsubsection 134(end)

\subsubsection{Bemerkung}
\label{ssub:135}
Ist $M$ eine sogenannte $C^p$-Mannigfaltigkeit, das heißt die Übergangsfunktionen sind $C^p$, so folgt mit dem Beweis von \ref{ssub:156}, das $TM$ eine $C^{p-1}$.Mannigfaltigkeit ist.\\
\\
$TM$ ist ein spezielles Beispiel der folgenden Klasse an Mannigfaltigkeiten.
% subsubsection 135(end)

\subsubsection[Definition: Vektorbündel]{Definition}
\label{ssub:136}
Seien $E,B$ zwei Mannigfaltigkeiten und $\pi: E\to B$ glatt. $\enbrace{E,B,\pi}$ heißt \Index{Vektorbündel} vom Rang $r$ über $B$, falls
\begin{enumerate}[(i)]
\item $\pi$ surjektiv
\item Es existiert eine offene Überdeckung $\set{U_i}_{i\in I}$ von $B$ und Diffeomorphismen $\Phi_i: \pi^{-1}(U_i) \to U_i\times \mathds{R}^r$, sodass $\forall x\in U_i:\ \Phi_i(\pi^{-1}(x)) = \set{x}\times \mathds{R}^r$.
\item Die "Übergangsfunktionen"
\[
\Phi_{ji} = \Phi_j \circ \Phi_i^{-1}: U_{ji}\times \mathds{R}^r \to U_{ji}\times \mathds{R}^r
\]
sind von der Form $\Phi_{ji} (x,v) = \enbrace*{x,g_{ji}(x) v}$, wobei $g_{ji}: U_{ji} \to GL(r,R)$, das heißt $g_{ji} = \enbrace*{(u_{kl}(x))_{k,l=1}^r}$
\begin{figure}[H]
	\centering{
\begin{tikzcd}[column sep=small]
\pi^{-1}(U_i) \arrow{rr}{\Phi_i} \arrow{dr}{\pi} & & U_i\times \mathds{R}^n \arrow{dl}{pr_1}\\
& U_i   &
\end{tikzcd}
	}
\end{figure}
Insbesondere ist $\Phi_{ji}(x,.):\mathds{R}^r \to \mathds{R}^r$ ein linearer Isomorphismus gegeben durch $g_{ji}(x)\in GL(r,\mathds{R}^)$	
\end{enumerate}
\minisec{Notation}
\begin{itemize}
\item $E$ Totalraum
\item $B$ Basis
\item $\pi^{-1}(x)$ Faser über $x$
\end{itemize}
\missingfigure{Vektorbündel}
% subsubsection 136(end)

\subsubsection{Beispiel}
\label{ssub:137}
\begin{enumerate}[(i)]
\item Sei $M$ eine Mannigfaltigkeit. $E = M^n\times \mathds{R}^r \xrightarrow{\pi = pr_1} M$ surjektiv und glatt, so ist dies ein "triviales Bündel", $\Phi: \pi^{-1}(M) = M^n\times \mathds{R}^r \xrightarrow{\id} M^n\times \mathds{R}^r,\ \Phi(x,.) = \id \in GL(r,\mathds{R})$.\\
Allgemein sagen wir, das $(E,\pi,B)$ trivial ist, falls ein Diffeomorphismus $\Phi: E^r \to B\times R^r$ existiert mit
\marginnote{$r =$ Rang von $E$}
\begin{figure}[H]
	\centering{
\begin{tikzcd}[column sep=small]
E \arrow{rr}{\Phi_i} \arrow{dr}{\pi} & & B\times \mathds{R}^n \arrow{dl}{pr_1}\\
& B &
\end{tikzcd}
	}
\end{figure}
kommutativ und $\forall x\in B$ ist $\Phi\vert_{\pi^{-1}(x)}$ ein linearer Isomorphismus.
\item $TM \to M$ ist ein Vektorbündel vom Rang $n = \dim M$. Vegleiche dazu Karten $\Phi_{(U,\varphi)}$ aus \ref{ssub:155}.
\end{enumerate}
% subsubsection 137(end)

\subsubsection{Bemerkung}
\label{ssub:138}
\begin{enumerate}[(i)]
\item Insbesondere gilt für alle $x\in B$ ist $\pi^{-1}(x)$ Faser über $x$ ein $r$-dim $VR\vert_{\mathds{R}}$ $\leadsto$ "Vektorbündel". So sei $\xi,\eta \in \pi^{-1}(x) \simeq \mathds{R}^r$. Es gilt
\[
\xi+\eta = \Phi^{-1}\enbrace[\big]{\overbrace{\underset{=(x,u)}{\Phi(\xi)} + \underset{=(x,v)}{\Phi(\eta)})}^{= (x,u+v)}}
\]
\item Ist $B$ zusammenhängend, so ist $E$ zusammenhängend. (Dies folgt aus der Tatsache das Fasern zusammenhängend sind und allgemeinen topologischen Argumenten.)
\end{enumerate}
% subsubsection 138(end)
% subsection 15(end)

\subsection{Das Differential einer glatten Abbildung}
\label{sub:16}

\subsubsection{Definition}
\label{ssub:139}
Das \Index{Differential} einer glatten Abbildung $f:M\to N$ in $x\in M$ ist definiert durch
\[
\mathd_x :T_x M \to T_{f(x)} N,\ d_x f \enbrace[\big]{[(U,\varphi,u)]} = [(V,\psi,\mathd_{\varphi(x)} \enbrace*{\psi \circ f \circ \varphi^{-1}}(u))]
\]
(kurz: $\mathd_x f(u) = v\in T_x N$ für Karten $(U,\varphi)$ um $x$ und $(V,\psi)$ um $f(x)$)\\
Dies ist wohldefiniert, denn sei $(\tilde{U},\tilde{\varphi},\tilde{u}) \sim (U,\varphi,u)$, so gilt mit der Kettenregel
\[
\mathd_{\tilde{\varphi}(x)} \enbrace*{\psi \circ f \circ \tilde{\varphi}^{-1}}\enbrace*{\tilde{u}} = 
\mathd_{\tilde{\varphi}(x)} \enbrace*{\psi \circ f \circ \tilde{\varphi}^{-1}}\enbrace*{\mathd_{\varphi(x)}(\tilde{\varphi} \circ \varphi^{-1})(u)} = 
\mathd_{\varphi(x)} \enbrace*{\psi \circ f \circ \varphi^{-1}}\enbrace*{u}
\]
Genauso hängt das Bild nicht von $\psi$ und $\varphi$ ab, das heißt das Differential von $f$ ist intrinsisch definiert.
\minisec{Definition des Differentials nach \ref{ssub:151}}
\missingfigure{Veranschaulichung der Definition des Differentials nach \ref{ssub:151}}
Sei $f:M\to N$ glatt, $[\alpha]\in T_x M$ im Sinne von \ref{ssub:151}, so gilt für das \Index{Differential}
\[
\mathd_x f([\alpha]) = [f\circ \alpha]\in T_{f(x)} N
\]
% subsubsection 139(end)

\subsubsection[Theorem:Kettenregel für Differentiale]{Theorem}
\label{ssub:140}
Seien $f:M\to N$ und $g:N\to P$ zwei glatte Abbildungen in $x\in M$ bzw. $f(x)\in N$. Dann ist $g\circ f: M\to P$ glatt in $x$ und 
\[
\mathd_x (g\circ f) = \mathd_{f(x)} g \circ \mathd_x f: T_x M \to T_{f(x)} N \to T_{(g\circ f)(x)} P
\]
\minisec{Beweis}
Übung
% subsubsection 140(end)

% subsection 16(end)

\subsection{Vektorfelder}
\label{sub:17}

Sei $E \xrightarrow{\pi} B$ ein Vektorbündel. Eine glatte Abbildung $\sigma:M\to E$ heißt Schnitt von $E$, falls $\pi\circ \sigma = \id_M$. ($\sigma: U\overset{\text{offen}}{\subset} M \to E$ glatt heißt lokaler Schnitt, falls $\pi \circ \sigma = \id_U$.)\\
$\Gamma(U,E)$ = Raum der lokalen Schnitte, $\Gamma(M,E) = \Gamma(E)$ = Raum der (globalen) Schnitte.
\missingfigure{Veranschaulichung eines Schnittes}

\subsubsection[Definition: Vektorfeld]{Definition}
\label{ssub:141}
Ein \Index{Vektorfeld} $X$ ist ein Schnitt von $TM$, das heißt $X\in \Gamma(TM) (=: \mathcal{X}(M))$
\marginnote{Das X muss anders aussehen..}
% subsubsection 141(end)

\subsubsection[Bemerkung: Vektorraumstruktur des $\Gamma(E)$]{Bemerkung}
\label{ssub:142}
$\Gamma(TM)$ (allgemein $\Gamma(E)$ für $E$ Vektorbündel) ist ein $\mathds{R}$-Vektorraum. Es gilt 
\[
(X + Y)(x) := X(x) + Y(x) \in T_x M \xrightarrow{\pi} \set{x}
\]
% subsubsection 142(end)

\subsubsection{Beispiel}
\label{ssub:143}
Betrachte $TS^n$:
\missingfigure{Vektorfelder auf $S^1$}
$X\in \Gamma(TS^n)$ ist $X:S^n \to \mathds{R}^{n+1}$ glatt mit $g_0(X(x),x) = 0$, wobei $g_0$ das Skalarprodukt.
% subsubsection 143(end)

\minisec{Übungsaufgabe} $E$ ist genau dann trivial, wenn $r$ Schnitte $\sigma_1(x),\ldots,\sigma_r(x)\in \Gamma(E)$ existieren, die in jeden Punkt linear unabhängig sind, das heißt $\forall\ x\in M:\sigma_1(x),\ldots,\sigma_r(x)$ Basis von $\pi^{-1}=: E_x$

\subsubsection[Satz: Parallelisierbarkeit des Tangentialbündels]{Satz}
\label{ssub:144}
$TM^n$ ist genau dann trivial, wenn $n$ Vektorfelder $X_1,\ldots,X_n$ existieren, die in jeden Punkt linear unabhängig sind. Man nennt $M^n$ dann auch \Index{parallelisierbar}.
% subsubsection 144(end)

\subsubsection{Bemerkung}
\label{ssub:145}
Parallelisierbarkeit ist eine sehr starke Bedingung! So ist zum Beispiel der $S^1$ parallelisierbar, der $S^2$ jedoch nicht.
% subsubsection 145(end)
\\
\\
Wir beginnen nun den 3. Ansatz zur Definition des Tangentialraums vorzustellen, die Definition des Algebraikers.
% subsubsection 145(end)

\subsubsection{Definition}
\label{ssub:146}
Seien $x \in M$, $U$ und $V$ offene Umgebungen um $x$. Sei ferner $f:U\to \mathds{R}$ und $g:V\to \mathds{R}$ zwei glatte Funktionen. Wir sagen, das $f$ äquivalent zu $g$ ist, falls offene Umgebung $W\subset U\cap V$ existiert mit $f\vert_W = g\vert_W$. Die Äquivalenzklasse heißt \Index{Keim} von $f$ in $x$ und wird $[f,x]$ geschrieben (oder einfach nur $[f]$ falls $x$ klar, oder sogar nur $f$). Die Menge der Keime in $x$ wird mit $C^\infty_x (M)$ oder kurz $C^\infty_x$ bezeichnet.
% subsubsection 146(end)

\subsubsection{Bemerkung}
\label{ssub:147}
\begin{enumerate}[(i)]
\item $C^\infty_x$ ist ($\infty$-)dimensionaler reeller Vektorraum: $[f]+[g]=[f+g]$. Es ist sogar eine Algebra.
\item 
\marginnote{In meiner Mitschrift steht $[f\circ \varphi]$ und nicht $[f\circ \varphi^{-1}]$, das macht jedoch keinen Sinn. Meinungen sind erwünscht.}
Es gilt $C_x^\infty (M) \simeq C_0^\infty (\mathds{R}^n)$. Sei $(U,\varphi)$ Karte um $x$ mit $\varphi(x) = 0$. So ist $[f]\in C_x^\infty (M) \mapsto [f\circ \varphi^{-1}]\in C_0^\infty (\mathds{R}^n)$ der gewünschte Isomorphismus.
Aber: Dieses Isomorphismus ist nicht kanonisch, das heißt er hängt von $\varphi$ ab!
\item $U\subset M$ offen und $x\in U$, dann ist $C_x^\infty(U) = \set[[f]]{f:U\to \mathds{R}\ \text{glatt}} \simeq C_x^\infty (M)$.\\
Klar: $[f]\in C_x^\infty (M) \hookrightarrow [f\vert_U]\in C_x^\infty (U)$. Bleibt die Surjektivität zu zeigen: Sei $[g]\in C_x^\infty (M)$, $g:U\to \mathds{R}$. Sei $\psi: M \to \mathds{R}$ eine Cut-Off Funktion,also mit $\supp \psi \subset U$ und $\psi \equiv 1$ auf Umgebung $W$ von $x$. So ist $\psi\vert_U \cdot g: U \to \mathds{R}$ glatt und $(\psi\vert_U \cdot g)\vert_{U\cap \enbrace{\supp \psi}^c} \equiv 0$. Setze $\psi\vert_U \cdot g$ durch $0$ auf $U^c$ fort, es ergibt sich $\tilde{g}:M\to \mathds{R}$ glatt mit $\tilde{g}\vert_W = \psi\vert_W \cdot g\vert_W = g\vert_W$, also $[\tilde{g}]\in C_x^\infty (M) \to [\tilde{g}\vert_U] = [g] \in C_x^\infty(U)$.
\end{enumerate}
% subsubsection 147(end)

\subsubsection{Definition}
\label{ssub:148}
Eine \Index{Derivation} auf $C_x^\infty (M)$ ist eine lineare Abbildung $\delta: C_x^\infty \to \mathds{R}$ mit 
\[
\delta(fg) = f(x)\delta(g) + \delta(f)g(x) \tag{Leibnitz-Regel}
\]
\minisec{Notation}
$D_x =$ Raum der Derivationen auf $C_x^\infty$ 
% subsubsection 148(end)

\subsubsection{Satz}
\label{ssub:149}
Jede Derivation $\delta \in D_0(\mathds{R}^n)$ ist von der Form
\marginnote{Nochmal zur Erinnerung: $v^j$ ist keine Potenz sondern ein Indize.}
\[
\delta(f) = \sum_{j=1}^n v^j \frac{\partial f}{\partial x^j}(0)\quad \text{für}\ v^j\in \mathds{R}
\]
\minisec{Beweis}
\begin{itemize}
\item Sei $U\subset \mathds{R}^n$. $\underset{\in C^\infty(U)}{f}	 \mapsto \frac{\partial}{\partial x^j}f(0)$ ist eine Derivation, da $\frac{\partial}{\partial x^j} f(0)$ nur von $[f]$ abhängt.
\item $\delta(c) = c\cdot \delta(1) = 0$, da $\delta(1) = \delta(1\cdot 1) = \delta(1) + \delta(1) = 2\delta(1)$. Es folgt $\delta(f-f(0)) = \delta(f)$, wobei
\[
f-f(0) = \int_0^1 \frac{\mathd}{\mathd t} f(tx) dt = \sum_{j=1}^n x^j \underbrace{\int_0^1 \frac{\partial}{\partial x^j} f(tx) dt}_{h_j(x)}
\]
Somit gilt
\[
\delta(f) = \delta(f-f(0)) = \delta\enbrace*{\sum_{j=1}^n x^j h_j(x)} = \sum_{j=1}^n \underset{=: v^j}{\delta(x^j)} \underset{= \frac{\partial}{\partial x^j}}{h_j(0)} + \underset{=0}{x^j(0)} \delta(h_j)
\]
\end{itemize}
% subsubsection 149(end)

\subsubsection*{Einschub}
\marginnote{$\delta$ lin. Abb. und $\delta:C_x^\infty \to \mathds{R}$, also $\delta \in C_x^\infty(M)^*$}
Wie schaut $\delta\in C_x^\infty (M)^*$ Derivation in der Karte $\varphi$ aus?\\
Sei $\varphi(x) = \enbrace*{x^1(x),\ldots, x^n(x)}$. Dann definiert
\[
\frac{\partial}{\partial x^j}(x): C^\infty_x (M) \to \mathds{R},\ \frac{\partial f}{\partial x^j}(x) := \frac{\partial}{\partial x^j}\enbrace*{f\circ \varphi^{-1}}\enbrace*{\varphi(x)}
\]
eine Derivation.
\missingfigure{Derivationen in der Karte}
Insbesondere gilt dann (vgl. \ref{ssub:179}), wobei $\varphi = (x^1,\ldots,x^n)$, also $x^j:U\to \mathds{R}$:
\[
\underset{\in C^\infty_x (M) \simeq C^\infty_0 (\mathds{R}^n}{\delta f} = \sum_j \delta(x^j) \frac{\partial}{\partial x^j}(x)
\]
Definiere für $v\in T_x M$ die Derivation:
\[
\mathcal{L}_v f := \mathd_x f(v) \in \mathds{R}
\]
das heißt $\mathcal{L}_v: C^\infty_x (M) \to \mathds(R)$.

\subsubsection{Theorem}
\label{ssub:150}
$v\in T_x M \to \mathcal{L}_v \in \mathcal{D}_x(M)$ ist ein linearer Isomorphismus.
\minisec{Beweis}
\begin{itemize}
\item Linearität klar
\item Surjektivität: Sei $v = [(U,\varphi,u)] \in T_x M$, $f\in C^\infty_x (M)$. Es gilt
\[
\mathcal{L}_v f = \mathd_x f(v) = \mathd_{\varphi(x)}\enbrace*{f\circ \varphi^{-1}}(u) = \sum_{j=1}^n u^j \frac{\partial}{\partial x^j}\enbrace*{f\circ \varphi^{-1}}\enbrace*{\varphi(x)} = \sum_{j=1}^n u^j \frac{\partial}{\partial x^j} f(x) 
\]
Ist also $\delta \in \mathcal{D}_x (M)$, dann $\delta = \mathcal{L}_{v_\delta} f$ mit $v_\delta = \benbrace*{U,\varphi,\enbrace*{\delta(x^1),\ldots,\delta(x^n)}}$
\item Injektivität: Angenommen $\mathcal{L}_v f = 0$ für alle $f\in C_x^\infty (M)$ und $v = \benbrace*{\enbrace*{U,\varphi,u}}$ mit oBdA $u = \enbrace*{u^1,\ldots, u^n},\ u^1\neq 0$. 
\marginnote{$\partial_i = \frac{\partial}{\partial x^i}$}
Sei $g$ auf $\varphi(U)$ $C^\infty$ mit $\partial_1 g(\varphi(x)) \neq 0$, $g \equiv g(x^1)$. Es folgt $\partial_j g(\varphi(x)) = 0,\ j \geq 2$. So ist 
\[
\mathcal{L}_v (g\circ \varphi) = \sum_{j=1}^n u^j \frac{\partial}{\partial x^j} g(\varphi(x)) = u^1 \frac{\partial}{\partial x^1} g(\varphi(x)) \neq 0 \lightning
\]
Also folgt aus $\mathcal{L}_v f = 0$ für alle $f\in C^\infty_x (M)$, das $v = 0$, das heißt Injektivität.
\end{itemize}
% subsubsection 150(end)
Insbesondere definiert $\frac{\partial}{\partial x^j} (x)$ für $\enbrace*{U,\varphi = (x^1,\ldots, x^n)}$ Karte einen Vektor in $T_x M$.
\minisec{Übungsaufgabe}
\begin{enumerate}[(i)]
\item $x\in U \mapsto \frac{\partial}{\partial x^j}(x)\in T_x M$ glatt, das heißt $\frac{\partial}{\partial x^j} \in \Gamma(U,TM\vert_U)$ "Schnitt von TM über $U$" oder ein Vektorfeld auf $U$.
\item $\forall\ x \in U$ ist $\enbrace*{\partial_{x^1}, \ldots, \partial_{x^n}}$ eine Basis von $T_x M$.
\missingfigure{}
\end{enumerate}
Insbesondere folgt: Sind $a^1,\ldots,a^n$ glatte Funktionen auf $U\subset M$, $(U,\varphi)$ Karte. So folgt
\[
X(x) = \sum_i a^i(x) \partial_{x^j}(x) \in T_x M
\]
ist ein glattes Vektorfeld auf $U$. Umgekehrt ist jedes $X\in \Gamma(U,TM)$ von dieser Form, das heißt ist $X\in \Gamma(TM)$, so gilt
\[
X\vert_{(U,\varphi)} = \sum a^j(x) \partial_j(x) = 
\begin{pmatrix}
a^1(x)\\
\vdots \\
a^n(x)\\
\end{pmatrix}\quad \text{Vektor mit glatten Einträgen}
\]
beschreibt Vektorfeld lokal.

\subsubsection[Definition: Derivation auf M]{Definition}
\label{ssub:151}
Eine \Index{Derivation} auf $M$ ist eine Abbildung $\delta:C^\infty (M) \to C^\infty (M)$ mit
\[
\delta(fg) = f\delta(g) + \delta(f)g
\]
\minisec{Beispiel}
Sei $X\in \Gamma(TM)$, so definiert
\[
\enbrace*{\mathcal{L}_X \underset{\in C^\infty (M)}{f}}(x) = \mathcal{L}_{X(x)} [f,x] = \mathd_x f(X(x)) =: \mathd f(X)(x)
\]
eine Derivation auf $M$.
% subsubsection 151(end)

\subsubsection{Bemerkung}
\label{ssub:152}
Eine Derivation auf $M$ ist ein lokaler Operator $\delta: C^\infty (M) \to C^\infty (M)$, das heißt ist $f\in C^\infty (M)$ mit $f\vert_U = 0$, so folgt $(\delta f)\vert_U = 0$. Sei nämlich $x\in U$. Dann nehme $\psi \in C^\infty (M)$ mit $\psi(x) = 1$, $\supp \psi \subset U$, also $\psi \vert_{U^c} \equiv 0$. So gilt $\psi \cdot f \equiv 0$ auf $M$, das heißt $\delta(\psi f) = 0$. Es gilt 
\[
\delta(\psi f)(x) = \delta(\psi)(x)\underset{= 0}{f(x)} + \underset{= 1}{\psi(x)}\delta(f)(x) = \delta(f)(x)
\]
für alle $x\in U$. Also $(\delta f)\vert_U = (\delta(\psi f))\vert_U = 0$.
% subsubsection 152(end)

\subsubsection{Theorem}
\label{ssub:153}
$\Gamma(TM) \to \mathcal{D}$, $X \mapsto \mathcal{L}_X$ ist ein Linearer Isomophismus.
\minisec{Beweis}
\begin{itemize}
\item Abbildung ist wohldefiniert:\\
Sei $X\vert_U = \sum_j a^j \delta X_j$, so gilt $\mathcal{L}_{X(x)}[f,x] = \sum_j \underbrace{a^j(x)}_{C^\infty} \underbrace{\delta_j f(x)}_{C^\infty}$. Also $\mathcal{L}_X: C^\infty(M) \to C^\infty(M)$.
\item Abbildung injektiv:\\
Angenommen $\mathcal{L}_X f = 0$ für alle $f\in C^\infty(M)$, das heißt $\mathcal{L}_X = 0$. Es ist zu zeigen, das $X = 0$. Dazu:
\[
\mathcal{L}_X f(x) = \mathcal{L}_{X(x)} [f,x] = 0\quad \forall\ [f,x]\in C^\infty_x (M)
\]
In \ref{ssub:150} wurde gezeigt, das $T_x M \simeq D_x$, also aus $\mathcal{L}_v [f,x] = 0$ für alle $[f,x]\in C^\infty_x (M)$ folgt $v = 0$. Somit gilt $\mathcal{L}_X f(x) = 0\ \forall f\in C^\infty (M)\Rightarrow X(x) = 0\ \forall x$, also $X = 0$. 
\item Abbildung surjektiv:\\
Sei $\delta \in \mathcal{D}(M)$. Gesucht ist dann $X\in \Gamma(TM)$ mit $\mathcal{L}_X = \delta$. Dazu sei $\underbrace{X(x)}_{\in T_x M \simeq \mathcal{D}_x} = \delta_x$ definiert durch $\delta_x [f,x] = (\delta f)(x)$. Es muss gezeigt werden ,das $\delta_X$ wohldefiniert ist. Dafür nehmen wir $[f,x] = [g,x]$ an. Nach der Definition \ref{ssub:176} der Keime existiert eine Umgebung $V$ von $x$ mit $f\vert_V = g\vert_V$, also $(f-g)\vert_V$. Mit Bemerkung \ref{ssub:152} folgt $\delta(f-g)\vert_V = 0$ und insbesondere $\delta(f)(x) = \delta(g)(x)$. Also ist die Definition von $X$ wohldefiniert und es gilt $X: M \to TM$ mit $\pi \circ X = \id_M$. Es bleibt zu zeigen, das $X$ glatt ist. Dazu gilt für geeignete Karte $(U,\varphi)$ 
\[
\delta_x(f) \overset{\ref{ssub:179}}{=} \sum_j \delta_x (x^j) \delta_{x^j} f(x) = \sum_j \delta(x^j)(x)\cdot \delta_x^j f(x) 
\]
Also
\[
X(x) = \sum_j \delta_x (x^j) \delta_{x^j} (x) = \sum_j \underbrace{\delta(x^j)(x)}_{C^\infty}\cdot \delta_x^j (x)
\]
das heißt $X$ ist $C^\infty$.
\end{itemize}
% subsubsection 153(end)
Im nächsten Abschnitt wollen wir uns anschauen, wie $\mathcal{L}_X$ und $\mathcal{L}_Y$ sinnvoll verkettet werden können. So ist 
\[
\mathcal{L}_X \circ \mathcal{L}_Y: C^\infty (M) \xrightarrow{\mathcal{L}_Y} C^\infty (M) \xrightarrow{\mathcal{L}_X} C^\infty (M)
\]
im Allgemeinen falsch, aber

\subsubsection{Definition und Satz}
\label{ssub:154}
Sei $X,Y \in \Gamma(TM) \to \mathcal{L}_X,\mathcal{L}_Y \in \mathcal{D}$. So ist 
\[
\benbrace*{\mathcal{L}_X,\mathcal{L}_Y}:= \mathcal{L}_X \mathcal{L}_Y - \mathcal{L}_Y \mathcal{L}_X 
\]
der \Index{Kommutator} bzw. die \Index{Lie-Klammer} von $X$ und $Y$ eine Derivation von $M$. Das zugehörige Vektorfeld wird $\benbrace{X,Y}$ geschrieben, das heißt $\mathcal{L}_{\benbrace{X,Y}} = \benbrace*{\mathcal{L}_X,\mathcal{L}_Y}$
\minisec{Beweis}
Es ist zu zeigen, das $\benbrace*{\mathcal{L}_X,\mathcal{L}_Y}$ die Leibniz-Regel für $f,g\in C^\infty (M)$ erfüllt.
\begin{align*}
\benbrace*{\mathcal{L}_X,\mathcal{L}_Y}(fg) 
=& \mathcal{L}_X \enbrace{f\mathcal{L}_Y g + g\mathcal{L}_X f} - \mathcal(L)_Y \enbrace{f\mathcal{L}_X g}\\ 
=& \mathcal{L}_X f \mathcal{L}_Y g + f \mathcal{L}_X \mathcal{L}_Y g + \mathcal{L}_X g \mathcal{L}_Y f + g \mathcal{L}_X \mathcal{L}_Y f \\
-& \mathcal{L}_Y f \mathcal{L}_X g - f \mathcal{L}_Y \mathcal{L}_X g - \mathcal{L}_Y g \mathcal{L}_X f - g \mathcal{L}_Y \mathcal{L}_X f\\
=& f\benbrace*{\mathcal{L}_X,\mathcal{L}_Y}(g) + g\benbrace*{\mathcal{L}_X,\mathcal{L}_Y}(f)
\end{align*}
% subsubsection 154(end)
% subsection 17(end)

\subsection{Der Fluß eines Vektorfeldes}
\label{sub:18}
\begin{wrapfigure}[8]{R}[1.5cm]{4cm}
\centering{\begin{tikzpicture}[xscale=1.3, yscale=1.1]
		%\draw[help lines] (0,0) grid (4,3);
		\coordinate (p) at (2.1,1.6);
		\coordinate (s) at (0.5,0.7);
		\coordinate (e) at (3.2,3);
		%\draw[<->, thick] (4,0) node[right]{$x$} -- (0,0) -- (0,3) node[above]{$y$};
		\foreach \x/\angle in {0.5/20, 1.5/40, 2.5/60, 3.5/80} {
		    \foreach \y in {0.5, 1.5, 2.5} {
		        %\fill (\x,\y) circle[radius=1pt];
		        \draw[->,thick]  (\x, \y) -- ++(\angle:1);
		    }
		}
		\draw[very thick,OliveDrab3] (s) to[out=20,in=220] (p) to[out=40, in=245] (e) node[below, xshift=0.1cm, yshift=-0.1cm]{$\gamma_p$};
		\draw[fill, Tomato3] (p) circle (0.03) node[below]{$p$}; 
	\end{tikzpicture}
	\caption{Integralkurve in einem Vektorfeld des $\mathds{R}^2$}}
\end{wrapfigure}
Frage: Existiert $c:I\to \mathds{R}^2$ mit $c(0) = x,\ \frac{\mathd}{\mathd t}c(t) = c'(t) = X(c(t))$\\
Wir wissen aus Analysis (I-IV), das Lösung mit Startbedingung $c(0) = x$ lokal um $0$ (das heißt $I = (-\epsilon,\epsilon)$) existieren und $\psi_X:U\times (-\epsilon,\epsilon) \to \mathds{R}^n$ mit $\psi(x,t) = c_x(t)$ glatte Abbildung. Wie sieht das auf Mannigfaltigkeiten aus?

\subsubsection{Satz und Definition}
\label{ssub:155}
Sei $X\in \Gamma(TM)$, $x\in M$. Dann existiert eine Kurve $c_X:0\in I\to M$ mit $c_x(0) = x$, $c_x'(t) := \mathd_t c_x \enbrace*{\frac{\mathd}{\mathd t}} = X(c(t))$, die sogenannte \Index{Integralkurve} von $X$ durch $x$.
\minisec{Beweis}
\marginnote{Ausführlich wurde Picard - Lindelöf für Mfkt. in der Mitschrift zu Diffformen und Mfkt. aus dem letzten Semester behandelt}
Lesen in Karte + Standard ODE Theorie ($\leadsto$ Berger-Gostiaux)
% subsubsection 155(end)

\subsubsection{Bemerkung}
\label{ssub:156}
Es folgt mit ODE - Theorie: Sind $c_x:I\to M,\ \tilde{c}_x:\tilde{I}\to M$ zwei Integralkurven von $X$ durch $x$, dann gilt $\tilde{c}_x = c_x$ auf $\tilde{I}\cap I$. Es existiert ein maximales Intervall $I_x$ mit Integralkurve $c_x: I_X\to M$
% subsubsection 156(end)

\subsubsection{Satz und Definition}
\label{ssub:157}
Für $x\in M$ existiert Umgebung $U$ um $x$ und $I$ Intervall um $0$, sodass die Abbildung
\[
\psi_X: U\times I\to M,\ \psi_X(x,t) := c_x(t)
\]
glatt ist. $\psi_X$ heißt \Index{lokaler Fluß} von $X$. ($c_x$ kann auf $I_X$ mit $I \subsetneq I_X$ definiert sein.)
\begin{figure}[H]
	\centering{
	\begin{tikzpicture}[scale = 0.8]
		%\draw[help lines] (0,0) grid (12,6);
		\draw[dashed] (1,3) .. controls +(80:1) and +(180:1) .. (3,5) node[above]{$U$}
		.. controls +(0:1) and +(90:1) .. (6,3.5)
		.. controls +(270:1.5) and +(0:1.5) .. (3,1.5)
		.. controls +(180:1) and +(270:1) .. (1,3);
		\draw[dashed] (5,3) .. controls +(90:1.5) and +(180:0) .. (7,4.5)
		.. controls +(0:1.5) and +(90:1) .. (11,3)  node[above,right]{$\psi_X(U)$}
		.. controls +(270:1.5) and +(0:2) .. (8,1)
		.. controls +(180:1.5) and +(270:2) .. (5,3);
		\foreach \x/\angle in {2.5/30, 4/15, 5.5/00, 7/345, 8.5/330}{
			\foreach \y in {2.5, 3, 3.5, 4}{
				\draw[->,thick] (\x,\y) -- ++(\angle:1);
			}
		}
		\draw[thick,OliveDrab3] (2.5,2.5) .. controls +(30:2) and +(150:2) .. (10,2.5) node[right]{$c_y$};
		\draw[thick,OliveDrab3] (2.5,3.25) .. controls +(30:2) and +(150:2) .. (10,3.25) node[right]{$c_x$};
	\end{tikzpicture}
	\caption{lokaler Fluß eines Vektorfeldes}}
\end{figure}
% subsubsection 157(end)

\subsubsection{Beispiel}
\label{ssub:158}
\begin{wrapfigure}[6]{L}[1.5cm]{4cm}
	\begin{tikzpicture}[xscale=0.6, yscale=0.4]
		%\draw[help lines] (0,0) grid (6,4);
		\draw [->] (3,2.05) -- (3,4);
		\draw [->] (3.05,2) -- (6,2);
		\draw [->] (3,1.95) -- (3,0);
		\draw [->] (2.95,2) -- (0,2);
		\draw (3,2) circle (0.05);
			\foreach \x in {3.5,2.5,1.5,0.5}
				\foreach \y in {4.5,3.5,2.5,1.5,0.5}
					\draw [->] (\y, \x) -- (\y +0.7, \x);
			\foreach \y in {4.2,3.2,2.2,1.2,0.2}
				\draw[color=Sienna3,->] (\y, 2) -- (\y +0.7, 2);
				\node[color=Sienna3,below] at (5,2) {$y = 0$};
	\end{tikzpicture}
	\caption{Vektorfeld auf $\mathds{R}^2$}
\end{wrapfigure}
Sei $X(x,y) = \partial_{x^1}(x,y)$ Vektorfeld auf $\mathds{R}^2$, so ist $c_{(x,y)} = (x,y) + (t,0)$ eine Integralkurve von $X$, denn $c_{(x,y)}$ und $c'(t) = \partial_{x^1}(c(t)) = X(c(t))$.\\
Für $y\neq 0$ ist $c_{(x,y)}(t)$ auf ganz $\mathds{R}^2$ definiert. Falls $y = 0$ ist $c_{(x,0)}:(0,\infty)\to \mathds{R}^2$ für $x>0$ und $c_{(x,0)}:(-\infty,0)\to \mathds{R}^2$ für $x<0$.\\.
\\.
% subsubsection 158(end)

\subsubsection{Definition}
\label{ssub:159}
Sei $X\in \Gamma(TM)$, definiere
\[
\mathcal{D}_X := \set[(x,y)]{x\in M, t\in I_X} \overset{\text{offen}}{\subset}
\]
wobei $I_X$ max. Intervall von $c_x$. Die glatte Abbildung $\psi_X:\mathcal{D}_X \to M,\ (x,t)\mapsto \psi_X(x,t) = c_x(t)$ heißt \Index{globaler Fluss}. $X$ heißt vollständig, falls $\mathcal{D}_X = M\times \mathds{R}$.
% subsubsection 159(end)

\subsubsection{Proposition}
\label{ssub:160}
Für einen vollständigen Fluss $\Psi$ gilt. (Schreibweise $\Psi_t (x) = \Psi(x,t), \Psi:M\to M$)
\begin{enumerate}
\item $\Psi_0 = \id_M$
\item $\Psi_s \circ \Psi_t = \Psi_{s+t}$
\item $\Psi_t$ ist ein Diffeomorphismus und $\enbrace{\Psi_t}^{-1} = \Psi_{-t}$
\end{enumerate}
Das heit $\Psi$ definiert einen Gruppenhomomorphismus $\enbrace{\mathds{R},+} \to \text{Diff}(M)$. So etwas nennt man eine 1-Parametergruppe.
% subsubsection 160(end)

\subsubsection{Bemerkung}
\label{ssub:161}
\ref{ssub:186} gilt für lokale Flüsse, sofern die Ausdrücke definiert sind. Ein lokaler Fluss: $\Psi_t:U\to M$ ist ein Diffeomorphismus auf das Bild.
% subsubsection 161(end)

\subsubsection{Satz}
\label{ssub:162}
Jedes Vektorfeld mit kompakten Träger ist schon vollständig. Insbesondere ist jedes Vektorfeld auf kompakten Mannigfaltigkeiten vollständig.
\minisec{Beweis}
\begin{enumerate}[(i)]
\item 
\marginnote{da $c_x'(t) = X(c(t)) = 0$}
Sei $K = \supp K = \overline{\set[x\in M]{X(x)\neq 0}}$. Für $x\in K^c$ ist $c_x(t) = x\ \forall t\in \mathds{R}$ die Integralkurve von $X$ um $x$.
\item Es gibt eine endliche Überdeckung $K\subset \bigcup_j U_j$, sodass der lokale Fluss $\Psi^{(j)}: U_j\times (-\epsilon_j,\epsilon_j) \to M$ definiert ist. Setze $\epsilon = \min_j e_j > 0$, dann ist $M\times (-\epsilon,\epsilon)\subset \mathcal{D}_x$.
\item Angenommen $c < \infty$, sodass $c_x(s) = \Psi_s (x)$ nur für $|s| < c$ existiert. Dann betrachte
\[
c_{\Psi_s(x)} (t) = \Psi\enbrace{\Psi_s (x), t} = \Psi_t(\Psi_s(x)) =\Psi_{t+s} (x)
\]
definiert für $|t| < \epsilon$ (nach (ii)). Wähle $s, t$, sodass $s+t > c$. $\lightning$ Also existiert $c_x(s)$ für alle $s\in \mathds{R}$.
\end{enumerate}
% subsubsection 162(end
% subsection 18(end)

\subsection{Verhalten von Vektorfelden unter Diffeomorphismen}
\label{sub:19}
\missingfigure{Abbildung zwischen Tangentialräumen}

\subsubsection[Proposition und Definition: Push-Forward]{Proposition und Definition}
\label{ssub:163}
Sei $\phi:M\to N$ ein Diffeomorphismus. $X\in \Gamma(TM)$, dann ist durch $y\mapsto \phi_{*} X(y) = \mathd_{\phi^{-1}(y)} \phi\enbrace{X(\phi^{-1}(y))}$ ein Vektorfeld auf $N$ gegeben, der \Index{Push-Forward} von $X$. Die assoziierte Derivation ist $\enbrace{\mathcal{L}_{\phi_* X} f}(x) = \mathcal{L}_X (f\circ \phi)(\phi^{-1}(y))$ für $f\in C^\infty (N)$
\minisec{Beweis}
\begin{enumerate}[(i)]
\item Es ist zu zeigen, das $\phi_* X$ ein Vektorfeld ist. Da $\phi_* X$ ein Schnitt auf $\Gamma(TN)$ ist, bleibt Glattheit von $\phi_* X$ zu zeigen. Wähle Karten $\enbrace*{U, \Phi = \enbrace*{x^1,\ldots,x^n}}$ um $\phi^{-1}(y)\in M$ und $\enbrace*{\phi(U),\Psi = \Phi \circ \phi = \enbrace*{y^,\ldots,y^n}}$ um $y\in N$.
\begin{enumerate}[Step 1:]
\item Zeige zunächst $\phi_* \delta_{x^j} = \delta_{y^j}$. Es ist $\delta_{x^j} = \benbrace{U,\Phi,e_j}$, $\delta_{y^j} = \benbrace{\phi(U),\Psi,e_j}$ wobei $e_j$ Standardbasis des $\mathds{R}^n$. Nach Definition gilt
\begin{align*}
\phi_* \delta_{x^j}\enbrace*{\phi^{-1}(y)} &= \mathd_{\phi^{-1} (y)} \phi\enbrace*{\benbrace*{U,\Phi,e_j}}\\
&= \benbrace[\big]{\phi(U),\Psi, \mathd_{\Phi\enbrace*{\phi^{-1}(y)}} \underbrace{\enbrace*{\Psi \circ \phi \circ \Phi^{-1}}}_{\id_{\mathds{R}^n}}(e_j)}\\
&= \benbrace{\phi(U),\Psi,e_j} = \delta_{y^j}(y)
\end{align*}
\item Sei $X = \sum_j a^j \delta_{x^j}$, dann gilt
\[
\phi_* X(y) = \sum_j a^j(\phi^{-1}(y))\phi_* \delta_{x^j}(y) = \sum_j \underbrace{a^j\enbrace{\phi^{-1}(y)}}_{\text{glatt}} \delta_{y^j} (y)
\]
Also ist $\phi_* X$ ein glattes Vektorfeld.
\end{enumerate}
\item Mit der lokalen Form der Lie-Ableitung und der oben gezeigten Darstellung von $\phi_* X$ folgt
\[
\mathcal{L}_{\phi_* X} f(y) = \sum_j a^j\enbrace*{\phi^{-1}(y)} \partial_{y^j} f(y)
\]
für $f\in C^\infty (N)$. Nach \ref{ssub:179} gilt
\begin{align*}
\partial_{y^j}(y) &= \sum_k \partial_{y^j}y^k(y)\partial_k \enbrace*{f\circ \Psi^{-1}}(\Psi(y))\\
&= \partial_j \enbrace*{f\circ \Psi^{-1}}(\Psi(y)) = \partial_j \enbrace*{f\circ \phi\circ \Phi^{-1}} = \partial_{x^j}(f\circ \phi)(\phi^{-1}(y))
\end{align*}
Somit folgt
\[
\mathcal{L}_{\phi_* X} f(y) = \sum_j a^j \enbrace*{\phi^{-1}(y)}\partial_{x^j}(f\circ \phi) (\phi^{-1}(y)) = \mathcal{L}_X (f\circ \phi)(\phi^{-1}(y))
\]
\end{enumerate}
% subsubsection 163(end)

\subsubsection{Bemerkung}
\label{ssub:164}
\begin{enumerate}[(i)]
\item Wir schreiben $\phi^* f = f\circ \phi$ für den \Index{Pullback} von $f$. Dann ist
\[
\mathcal{L}_{\phi_* X} = \enbrace{\phi^{-1}}^* \circ \mathcal{L}_X \circ \phi
\]
\item Es ist $(\phi\circ \psi)^* f = f\circ \phi\ psi = \psi^*(f\circ \phi) = \psi^* \phi^* f$, das heißt der Pullback $(\phi\circ \psi)^* = \psi^* \circ \phi^*$ ist kontravariant und der Push-Forward $(\phi\circ \psi)_* X = \phi_* \circ \psi_* X$ ist kovariant.
\end{enumerate}
% subsubsection 164(end)

\subsubsection{Korollar}
\label{ssub:165}
Sei $\phi: M\to N$ ein Diffeomorphismus, $X,Y$ Vektorfelder auf $M$. Dann gilt
\[
\phi_*\benbrace*{X,Y} = \benbrace{\phi_* X,\phi_* Y}
\]
\minisec{Beweis}
\begin{align*}
\mathcal{L}_{\benbrace{\phi_* X, \phi_* Y}} &= \mathcal{L}_{\phi_* X} \mathcal{L}_{\phi_* Y} - \mathcal{L}_{\phi_* Y} \mathcal{L}_{\phi_* X}\\
&= \enbrace{\phi^{-1}}^* \circ \mathcal{L}_{X} \circ \phi^* \circ \enbrace{\phi^{-1}}^* \circ \mathcal{L}_{Y} \circ \phi^* - \enbrace{\phi^{-1}}^* \circ \mathcal{L}_{Y} \circ \phi^* \circ \enbrace{\phi^{-1}}^* \circ \mathcal{L}_{X} \circ \phi^*\\
&= \enbrace{\phi^{-1}}^* \circ \mathcal{L}_{X} \circ \mathcal{L}_{Y} \circ \phi^* - \enbrace{\phi^{-1}}^* \circ \mathcal{L}_{Y} \circ \mathcal{L}_{X} \circ \phi^*\\
&= \enbrace{\phi^{-1}}^* \circ (\mathcal{L}_{X} \circ \mathcal{L}_{Y} - \mathcal{L}_{Y} \circ \mathcal{L}_{X}) \circ \phi^*\\
&= \enbrace{\phi^{-1}}^* \circ \mathcal{L}_{\benbrace{X,Y}} \circ \phi^* = \mathcal{L}_{\phi_* \benbrace{X,Y}}
\end{align*}
% subsubsection 165(end)

\subsubsection{Proposition}
\label{ssub:166}
Für (lokale) Flüsse gilt:
\[
\Psi_{\phi_* X,t} = \phi \circ \Psi_{X,t} \circ \phi^{-1}
\]
\minisec{Beweis}
Es ist zu zeigen,das für alle $s$, für die $\Psi_{X,s}(x)$ definiert ist, gilt:
\[
\frac{\mathd}{\mathd t} \phi \circ \Psi_{X,t} \circ \phi^{-1}\vert_{t=s} = \phi_* X\enbrace{\phi \circ \Psi_{X,s} \circ \phi^{-1}}
\]
\begin{enumerate}[(i)]
\item Sei $s = 0$. Es gilt
\begin{align*}
\frac{\mathd}{\mathd t} \phi \circ \Psi_{X,t} \circ \phi^{-1}(y) \Big\vert_{t=0}& = \frac{\mathd}{\mathd t} \phi \enbrace*{\Psi_{X,t} \circ \phi^{-1}(y)} \Big\vert_{t=0}\\
&= \mathd_{\phi^{-1}(y)} \phi\enbrace[\bigg]{\frac{\mathd}{\mathd t} \Psi_{X,t}\enbrace*{\phi^{-1}(y)}\Big\vert_{t=0}}\\
&= \mathd_{\phi^{-1}(y)} \phi \enbrace*{X \enbrace*{\phi^{-1}(y)}}\\
&= \phi_* X(y)
\end{align*}
\item Setze $r = t - s$, $\Psi_{X,t} = \Psi_{X,r} \circ \Psi_{X,s}$. Es gilt
\[
\frac{\mathd}{\mathd t} \phi \circ \Psi_{X,t} \circ \phi^{-1}(y) \Big\vert_{t=s} = \frac{\mathd}{\mathd r} \phi \circ \Psi_{X,r} \circ \Psi_{X,s} \circ \phi^{-1}(y) \Big\vert_{r=0} = \phi_* X\enbrace*{\phi \circ \Psi_{X,s}\enbrace*{\phi^{-1}(y)}}
\]
\end{enumerate}
% subsubsection 166(end)

\subsubsection{Theorem}
\label{ssub:167}
Seien $X,Y \in \Gamma(TM)$, $\Psi_t$ sei der Fluss von $Y$. Dann gilt
\[
\frac{\mathd}{\mathd t} \enbrace{\Psi_t^{-1}}_* X \Big\vert_{t=0} = \benbrace{Y,X}
\]
\minisec{Beweis}
siehe Skript
% subsubsection 167(end)

\subsubsection[Korollar: Jacobi-Identität]{Korollar (Jacobi-Identität)}
\label{ssub:168}
Seien $X,Y,Z\in \Gamma(TM)$. Dann gilt
\[
\benbrace{X,\benbrace{Y,Z}} + \benbrace{Z,\benbrace{X,Y}} + \benbrace{Y,\benbrace{Z,X}} = 0
\]
\minisec{Beweis}
Übungsaufgabe
% subsubsection 168(end)

\subsubsection{Bemerkung}
\label{ssub:169}
Sei $V$ ein Vektorraum (in der Regel $\mathds{R}$ oder $\mathds{C}$). Sei $\benbrace{.,.}:V\times V \to V$ eine schiefsymmetrische bilineare Abbildung mit Jacobi-Identität, zum Beispiel $V = \Gamma(TM)$ mit $\benbrace{.,.} = $ Lie-Klammer. Dann heißt $\enbrace{V,\benbrace{.,.}}$ Lie-Algebra. Ein weiteres Beispiel ist $\mathds{R}^{n\times n}$ mit $\benbrace{A,B} = AB-BA$ "Matrixkommutator". Die Lie-Algebra ist ein zentrales Objekt der Darstellungstheorie von Lie-Gruppen, zum Beispiel $\GL(n,\mathds{R}),\ O(n),\ \SU(n)$
% subsubsection 169(end)

% subsection 19(end)

\subsection{Tensorbündel}
\label{sub:110}

\subsubsection[Wiederholung:Tensorprodukt von Vektorräumen]{Wiederholung}
\label{ssub:170}
Seien $V,W$ (endlich) dimensionale VR$\vert_{\mathds{R}}$. Dann existiert ein VR$\vert_{\mathds{R}}$ mit folgender Eigenschaft. Ist $U$ VR$\vert_{\mathds{R}}$, dann
\[
\underbrace{L_2\enbrace{V\times W, U}}_{\text{bilineare Abb.} V\times W\to U} \underbrace{\cong L\enbrace{V\otimes W, U}}_{\text{lineare Abb.}V\otimes W \to U}
\]
Genauer:
\begin{figure}[H]
	\centering{
\begin{tikzcd}[column sep=small]
& V\otimes W \arrow{d}{\hat{B}}\\
V\times W \arrow{r}{B} \arrow{ur}{L_2\enbrace{V\times W,V\otimes W} \ni \otimes} & U
\end{tikzcd}
	\caption{"universelle Eigenschaft" von $\otimes$}}
\end{figure}
wobei $B\in L_2\enbrace{V\times W, U} \simeq L \enbrace{V\otimes W, U}$ gegeben durch $B\mapsto \hat{B}$. $\hat{B}$ ist eindeutig bestimmt, das heißt $B(v,w) = \hat{B}(v\otimes w)$.\\
Im Allgemeinen:
\[
\underbrace{L_n\enbrace{V_1\times \ldots \times V_n, U}}_{\text{n-lineare Abbildung}} \simeq L\enbrace{V_1\otimes \ldots \otimes V_n, U}
\]
Es gilt $V\otimes W\cong W\otimes V$ und $\enbrace{V\otimes W} \otimes U \cong V\otimes \enbrace{W\otimes V}$, das heißt $\otimes$ ist kommutativ und assoziativ. 
\minisec{Konkret}
 Ist $\set{e_i}_{i = 1}^n,\ \set{f_i}_{i = 1}^n$ Basis von $V$ beziehungsweise $W$, so ist die Basis von $V\otimes W$ gegeben durch $\set{e_i\otimes f_j}_{i,j = 1}^{n,m}$(, wobei das Bild von $(e_i,f_j)$ unter $\otimes$: $\otimes(e_i,f_j) = e_i \otimes f_j$). Insbesondere ist $\dim V\otimes W = \dim V \cdot \dim W$. Dann gilt für $v = \sum v^i e_i,w = \sum w^j f_j$.
\[
\otimes(v,w) = v\otimes w = \otimes\enbrace{\sum v^i e_i, \sum w^j e_j} = \sum_{i,j} v^i w^j e_i\otimes f_j
\]
\minisec{Bemerkung}
\begin{enumerate}[(i)]
\item Seien $v,v'\in V, w\in W$, so ist $v\otimes w\in V\otimes W$ und 
\[
(v+v')\otimes w = \otimes (v+v',w) = \otimes (v,w) + \otimes (v',w) = v\otimes w + v'\otimes w
\]
\item In der Regel lässt sich $t\in V\otimes W$ nicht als $v\otimes w$ für $v\in V,w\in W$ schreiben, denn sei $\set{e_i\otimes f_j}$ Basis von $V\otimes W$, das heißt $t = \sum t^{ij} e_i\otimes f_j$ für $t^{ij}\in \mathds{R}$. Ist $t = v\otimes w = \sum_{i,j} \enbrace{v^i e_i}\otimes \enbrace{w^j f_j} \overset{(i)}{=} \sum v^i w^j e_i\otimes f_j$, das heißt $t^{ij} = v^i\cdot w^j,\ \forall i,j$ ist im Allgemeinen falsch.\\
Aber $\set[v\otimes w]{v\in V, w\in W}$ ist ein Erzeugendensystem von $V\otimes W$
\end{enumerate}
% subsubsection 170(end)

\subsubsection{Beispiel}
\label{ssub:171}
Seien $V,W$ VR$\vert_{\mathds{R}}$, also $V^*,W^*$ der zugehörige duale Vektorraum. Dann ist
\[
V^* \otimes W^* \cong L_2\enbrace{V\times W,\mathds{R}} \cong L\enbrace{V\otimes W,\mathds{R}} = \enbrace{V\otimes W}^*
\]
\[
\alpha \otimes \beta \in V^* \otimes W^* \mapsto \enbrace*{(v,w)\in V\times W \to \alpha(v)\beta(w) \in \mathds{R}} \in L_2 \enbrace{V\times W,\mathds{R}}
\]
\minisec{Zum Beispiel}
Sei $V = \mathds{R}^2 = \lenbrace*{\begin{pmatrix}1\\0\end{pmatrix},\begin{pmatrix}0\\1\end{pmatrix}}, 
W = \mathds{R}^3 = 
\lenbrace*{\begin{pmatrix}1\\0\\0\end{pmatrix},\begin{pmatrix}0\\1\\0\end{pmatrix},\begin{pmatrix}0\\0\\1\end{pmatrix}
}
$. So ist 
\[L_2\enbrace{\mathds{R}^2\times \mathds{R}^3,\mathds{R}} = \set[
\begin{pmatrix}
x_1 & x_2 & x_3\\
y_1 & y_2 & y_3
\end{pmatrix}]{x_j,v_j \in \mathds{R}},\quad
(\alpha_1,\alpha_2)\begin{pmatrix}
x_1 & x_2 & x_3\\
y_1 & y_2 & y_3
\end{pmatrix}\begin{pmatrix}b_1\\b_2\\b_3\end{pmatrix} = B(v,w)
\]
$B = x_1 e^1\otimes f^1 + x_2 e^1\otimes f^2 + x_3 e^1\otimes f^3 + \sum_{j=1}^3 y^j e^2\otimes f^j,\quad e^i,f^j$ duale Basis. Zum Beispiel
\[
B\enbrace*{\begin{pmatrix}1\\0\end{pmatrix},\begin{pmatrix}0\\1\\0\end{pmatrix}} = \sum_{j=1}^3 \enbrace*{x_j e^1\otimes f^j + y_j e^2\otimes f^j}(e_1,f_2) = \sum_{j=1} x_j \underbrace{e^1(e_1)}_{1}\cdot \underbrace{f^j(f_2)}_{0} + y_j e^2(e_1) f^j(f_2) = x_2
\]
Übungsaugabe: $L(V,W) \simeq W\otimes V^*$
% subsubsection 171(end)

\subsubsection{Tensorprodukt der linearen Abbildung}
\label{ssub:172}
Erinnerung: Sei $V, W$ Vektorräume und $V^*,W^*$ der zugehörige duale Vektorraum. Sei $f:V\to W$ linearer Morphismus, so ist $f^T: W^*\to V^*$ die Transponierte von $f$. Seien $a\in L(V,V'), b\in L(W,W')$, so ist $a\otimes b \in L(V\otimes W,V'\otimes W')$ das Tensorprodukt von $a$ und $b$ definiert durch $a\otimes(v\otimes w) = a(v)\otimes b(w)\in V'\otimes W'$ und setzen es linear fort.
\[
a\otimes b (v\otimes w +v'\otimes w') = a\otimes b(v\otimes w) + a\otimes b(v'\otimes w')
\]
etc. Zum Beispiel
\[
a\otimes b (v\otimes w +v\otimes w') = a(v)\otimes b(w) + a(v)\otimes b(w') = a(v)\otimes (b(w)+b(w')) = a(v)\otimes b(w+w')
\]
% subsubsection 172(end)

\subsubsection{Tensorbündel}
\label{ssub:173}
Wir betrachten Tensoren nun über Mannigfaltigkeiten $M$. Definiere
\[
\bigotimes^p TM = \bigcup_{x\in M} \bigotimes^p T_x M\times \set{x}
\]
wobei $\otimes^p T_x M = \underbrace{T_x M \otimes \ldots \otimes T_x M}_{\text{p-mal}} = T_x M^{\otimes^p}$, das p-te Tensorprodukt von $TM$. Dies ist ein Vektorbündel $\pi_{\otimes^p TM}: \otimes^p TM\to M$ von Rang $(\dim M)^p$
\begin{itemize}
\item $v\in T_x M^{\otimes^p}\times \set{x} \mapsto x = \pi_{\otimes^p TM}(v) \in M$
\item Sei $(U,\varphi)$ Karte von $M$, $\enbrace{\pi_{\otimes^p TM}^{-1}(U), \Theta_{U,\varphi}^{-1}}$ Karten von $TM$, so ist $\enbrace{\pi_{\otimes^p TM}^{-1}(U), \Theta^{-1}_{U,\varphi}\otimes \ldots \otimes \Theta^{-1}_{U,\varphi}}$ Karte von $\otimes^p TM$
\end{itemize} 
% subsubsection 173(end)

\subsubsection{Kotangentialbündel von M}
\label{ssub:174}
Idee: Ersetze $T_x M$ durch $T_x^* M$ und erhalte das \Index{Kotangentialbündel} $T^* M = \bigcup_{x\in M} T_x^* M \times \set{x}$. Atlas ist gegeben durch $\enbrace{\pi_{\otimes^p TM}^{-1}(U), \enbrace{\Theta_{U,\varphi}}^{T}}$. \ref{ssub:173} und \ref{ssub:174} können allgemein für jedes Vektorfeld durchgeführt werden, das heißt sind $E\to M, F\to M$ Vektorbündel über $M$, so können $E^*\to M,\ E\oplus F\to M,\ E\otimes F\to M$ etc. konstruiert werden.
Insbesondere gilt: Ist $x^1,\ldots,x^n$ Koordinatensystem auf $U\subset M$, $\partial_{x^1},\ldots, \partial_{x^n}$ Basis von $TM\vert_U$, so ist $\partial_{x^{i_1}}\otimes \ldots \otimes \partial_{x^{i_p}}$ Basis von $TM^{\otimes^p}$ und $\enbrace*{\partial_{x^i}}^* =: \mathd x^i$ Basis von $T^* M\vert_U$
% subsubsection 174(end)

\subsubsection{Tensorbündel}
\label{ssub:175}
Sei $\enbrace{p,q}\in \mathds{N}\times \mathds{N}$, so ist
\[
\mathcal{T}^p_q M := TM^{\otimes^p}\otimes T^* M^{\otimes^q}\to M
\]
Bündel der $(p,q)$-Tensoren, glatte Schnitte aus $\Gamma(\mathcal{T}^p_q M)$ heißen $(p,q)$- Tensor(felder).
\minisec{Spezialfall}
\begin{itemize}
\item $\mathcal{T}_0^0 M := M\times \mathds{R}$, $\Gamma(\mathcal{T}_0^0) = C^\infty (M)$
\item $\mathcal{T}_0^1 M := TM$, $\Gamma(TM)$ Vektorfelder
\item $\mathcal{T}_1^0 M := T^*M$, $\Gamma(T^* M)$ 1-Formen
\item $\mathcal{T}_0^p M := \mathcal{T}_0^p M$ kovariante Tensoren
\end{itemize}
% subsubsection 175(end)

\subsubsection{Operatoren für Tensoren}
\label{ssub:176}
Eine kleine Zeitreise zurück in die lineare Algebra.
\begin{enumerate}[(i)]
\item \Index{Tensorprodukt}: $\mathcal{T}_q^p V \times \mathcal{T}_{q'}^{p'} V \to \mathcal{T}_{q+q'}^{p+p'	} V$ bilinear, definiert durch:
\begin{align*}
&\enbrace*{v_1\otimes \ldots \otimes v_p \otimes v^1 \otimes \ldots \otimes v^p, w_1\otimes \ldots \otimes \w_{p'} \otimes w^1 \otimes \ldots \otimes w^{p'}}\\
\mapsto\ &v_1\otimes \ldots \otimes v_p \otimes w_1\otimes \ldots \otimes \w_{p'}\otimes v^1 \otimes \ldots \otimes v^p \otimes w^1 \otimes \ldots \otimes w^{p'}
\end{align*}
\item \Index{Kontraktion} ("\Index{Tensorverjüngung}"): $C_j^i: \mathcal{T}_q^p V \to \mathcal{T}_{q-1}^{p-1}$, definiert durch:
\marginnote{$\hat{v}_i$ bzw. $\hat{v}^j$ heißt, das dies auslassen}
\begin{align*}
&C_j^i\enbrace*{v_1\otimes \ldots \otimes v_i \otimes \ldots \otimes v_p \otimes v^1\otimes \ldots \otimes v^i \otimes \ldots \otimes v^q}\\
&= v^j(v_i) v_1\otimes \ldots \otimes \hat{v}_i \otimes \ldots \otimes v_p \otimes v^1\otimes \ldots \otimes \hat{v}^i \otimes \ldots \otimes v^q
\end{align*}
\minisec{Beispiel}
Spurbildung eines Endomorphismus $A\in \End(V)$.\\
Ist $e_1,\ldots, e_n$ Basis von $V$,so können wir $A$ also als $A = (A_j^i)\in \mathds{R}^{n\times n}$ betrachten und die Spur definieren $\Tr(A) = \sum A_i^i$. Da $\Tr(\Gamma A \Gamma^{-1}) = \Tr(A)$ für $\Gamma \in \GL(n,\mathds{R})$, ist die Definition der Spur unabhängig von der Basis.\\
Alternativ:\\
$\End(V) \cong V\otimes V^*$, dann $\Tr(A) = C_1^1(A)$. Zum Beispiel $v\otimes \alpha \cong A\cdot \w \mapsto \alpha(w)\cdot v$, $\Tr(A) = \alpha(v) \in \mathds{R}$.
\item Ein Spezialfall von (ii), das \Index{innere Produkt} für kovariante Tensoren und Vektoren. Sei $v\in V,\ \alpha\in \mathcal{T}_q^0 V = \mathcal{T}_q V$.
\[
i_v \alpha := C_1^1\enbrace{\underbrace{v\otimes \alpha}_{\in \mathcal{T}_q^1 V}} \in \mathcal{T}_{q-1} V = \underbrace{V^*\otimes \ldots \otimes V^*}_{q-1 \text{- mal}} \cong L_{q-1}\enbrace{V\times \ldots \times V,\mathds{R}}
\]
Konkret $i_v \alpha\enbrace{v_1,\ldots,v_{q-1}} = \alpha\enbrace{v,v_1,\ldots,v_{q-1}}$.
\end{enumerate}
Dies induziert entsprechende Operationen auf Mannigfaltigkeiten indem man Tensoren durch Tensorfelder ersetzt und faserweise definiert. Sei zum Beispiel $S\in \Gamma(\mathcal{T}_q^p M), S'\in \Gamma(\mathcal{T}_{q'}^{p'} M)$, so ist $S\otimes S' \in \Gamma(\mathcal{T}_{q+q'}^{p+p'} M)$ definiert durch $(S\otimes S')_{x\in M} = S_x\otimes S'_x \in \mathcal{T}_{q+q'}^{p+p'} T_x M$. Es bleibt zu zeigen, das $S\otimes S'$ wirklich glatt ist. Dies ist klar, denn lokal gilt 
\[S\vert_U = \sum S_{j_1,\ldots,j_q}^{i_1,\ldots,i_p} \partial_{i_1}\otimes \ldots \otimes \partial_{i_p} \otimes \mathd x^{i_1}\otimes \ldots \otimes \mathd x^{i_q}
\]
\[
S'\vert_U = \sum (S')_{k_1,\ldots,k_{q'}}^{l_1,\ldots,l_{p'}} \partial_{l_1}\otimes \ldots \otimes \partial_{l_{p'}} \otimes \mathd x^{k_1}\otimes \ldots \otimes \mathd x^{k_{q'}}
\]
und somit folgt lokal für das Tensorprodukt $S\otimes S'$
\[
S\otimes S'\vert_U = \sum \underbrace{S_{j_1,\ldots,j_q}^{i_1,\ldots,i_p} \sum (S')_{k_1,\ldots,k_q}^{l_1,\ldots,l_p}}_{\text{glatt}} \partial_{l_1}\otimes \ldots \otimes \partial_{i_p'+l_{p'}} \otimes \mathd x^{k_1}\otimes \ldots \otimes \mathd x^{j_q + k_{q'}}
\]
oder sei $X\in \Gamma(TM) = \Gamma(\mathcal{T}_0^1 M)$, $\alpha\in \Gamma(\mathcal{T}_1^0 M) = \Gamma(T^* M)$ mit 
\[
i_x(\alpha) = C_1^1(X,\alpha) = \alpha(X) \in C^\infty (M),\quad a(X)\underset{(x)}{\in M} = \underset{\alpha_x}{\in T^*_x M}(X(x))\in \mathds{R}
\]
so ist $\tilde{\alpha}:\Gamma(TM) \to C^\infty(M)$, definiert durch $X\mapsto \alpha(X) = i_X(\alpha)$ $\mathds{R}$-linear und insbesondere auch $C^\infty (M)$-linear.
\begin{figure}[H]
\centering{
\begin{minipage}{0.9\textwidth}
$\Gamma (TM)$, allgemein $\Gamma(\mathcal{T}_q^p)$ sind
\begin{enumerate}[(a)]
\item $\mathds{R}$-Vektorräume (punktweise Addition + Skalarmultiplikation)
\item $\underbrace{C^p (M)}_{\text{Ring}}$-Modul: $f\in C^\infty (M), S\in \Gamma(\mathcal{T}_q^p M)$, so gilt $\enbrace{fS}_x = f(x)\cdot S_x$
\end{enumerate}
\end{minipage}}
\end{figure}
Dann $\alpha(fX)(x) = \alpha_x(f(x)\cdot X(x)) \underset{\mathds{R}\text{-linear}}{=} f(x) \underbrace{\alpha_x(X(x))}_{\alpha(X)(x)} = \enbrace{f\cdot \alpha(X)}(x)$.\\
Allgemein: Sei $S\in \Gamma(\mathcal{T}_q^p M) = \Gamma(\bigotimes^p TM \otimes \bigotimes^q T^* M)$, so ist 
\[
\tilde{S}:\Gamma(TM)\times \ldots \times \Gamma(TM)\times \Gamma(T^* M)\times \ldots \times \Gamma(T^* M)\to C^\infty (M)
\]
mit $\tilde{S}\in L_{q+p}\enbrace*{\bigtimes^q \Gamma(TM)\times \bigtimes^p \Gamma(T^* M),C^\infty (M)}$ definiert durch
\[
\tilde{S}\enbrace*{X_1,\ldots,X_q,\alpha^1,\ldots, \alpha^p}(x) = \underset{\in \mathds{R}}{S_x}\enbrace*{\alpha^1(x),\ldots, \alpha^p(x), X_1(x),\ldots, X_q(x)}
\]
\todo{Hier fehlt noch etwas}
%subsubsection 176(end)

\subsubsection{Beispiel}
\label{ssub:177}
Betrachte Lie-Klammer als Abbildung
\[
\Gamma(TM)\times \Gamma(TM)\times \Gamma(T^* M) \to C^\infty (M),\quad (X,Y,\alpha)\mapsto \alpha\enbrace{\benbrace{X,Y}} = i_{\benbrace{X,Y}} \alpha
\]
Wobei beachte $L_r\enbrace{V\times\ldots \times V,W}\simeq L_r\enbrace{V\times \ldots \times W^*,\mathds{R}}$, zum Beispiel 
\[
L(V,W) \simeq L(V\times W^*,\mathds{R}) = V\otimes W^*,\quad V\to W \simeq V\times W^* \to \mathds{R}
\]
für Mannigfaltigkeiten :
\[
\Gamma(TM)\times \Gamma(TM) \xrightarrow{[.,.] \mathds{R}-\text{linear}} \Gamma(TM) \cong \Gamma(TM)\times \Gamma(TM)\times \Gamma(T^* M) \to C^\infty (M)
\]
Dann ist $\benbrace{.,.}$ nicht $C^\infty (M)$-linear (aber natürlich $\mathds{R}$-linear). Zum Beispiel sei $f\in C^\infty (M)$
\begin{align*}
\enbrace{f\cdot X, Y,\alpha} \mapsto \alpha\enbrace{\underbrace{\benbrace{fX,Y}}_{
fX(Y)-Y(fX) = fX(Y) - Y(f)X - fY(X)
}} =& \alpha\enbrace{-\mathd f(Y)\cdot X + f(\benbrace{X,Y})}\\
=& -\mathd f(Y)\alpha(X) + f\alpha\enbrace{\benbrace{X,Y}}\\
\neq& f\alpha(\benbrace{X,Y})
\end{align*}
% subsubsection 177(end)

\subsubsection{Satz}
\label{ssub:178}
Sei $\Lambda \in L_{p+q} \enbrace{\bigtimes^p \Gamma(T^* M)\times \bigtimes^q \Gamma(TM),C^\infty (M)}$. Dann ist $\Lambda$ $C^\infty$-linear $\Leftrightarrow \Lambda = \tilde{S}$ für $S\in \Gamma(\mathcal{T}_q^p M)$. Hier ist $\tilde{S}:\Gamma(T^* M)\times \ldots \times \Gamma(T^* M)\times \Gamma(TM)\times \ldots \times \Gamma(TM)\to C^\infty(M)$ definiert durch
\[
\tilde{S}\enbrace{\alpha^1,\ldots,\alpha^p,X_1,\ldots,X_q}(x) = S_x\enbrace{\alpha_x^1,\ldots,\alpha_x^p,X_1(x),\ldots,X_p(x)}
\]
\minisec{Beweis}
\begin{itemize}
\item "$\Leftarrow$":
\item "$\Rightarrow$": Sei $\Lambda$ $C^\infty$-linear gegeben. Definiere $S_x\in \mathcal{T}_q^p T_x M\cong L_{p+q}\enbrace*{\bigtimes^p T^*_x M \times \bigtimes^q T_x M,\mathds{R}}$ durch
\[
S_x\enbrace{\alpha^1,\ldots,\alpha^p,v_1,\ldots,v_q} = \underbrace{\Lambda\enbrace{\hat{\alpha}^1,\ldots,\hat{\alpha}^p,\hat{v}_1,\ldots,\hat{v}_q}}_{C^\infty(M)}(x) \in \mathds{R}
\]
mit $\hat{\alpha}_j \in \Gamma(T^* M), \hat{v}_j \in \Gamma(TM)$ glatte Fortsetzung nach \ref{ssub:179}
\begin{enumerate}[(a)]
\item Es ist zu zeigen, das $S$ ein wohldefinierter Schnitt ist,das heißt unabhängig von der Fortsetzung. Dazu oBdA $\Lambda:\Gamma(TM)\to C^\infty (M)$ (allgemeiner Fall läuft genauso). Dann gilt $\Lambda$ lokaler Operator. Stärker noch:
\[
X(x) = 0\Rightarrow \Lambda(X)(x) = 0
\]
,denn
\[
\Lambda(X)\underset{\in U\ \text{Karte}}{(x)} = \Lambda(X)\vert_U (x) \underset{\text{lok. Operator}}{=} \underbrace{\Lambda(X\vert_U)}_{\sum X^j \partial_j,\ X^j\in C^\infty (M)}(x) = \sum X^j(x) \Lambda(\partial_j)(x) = 0
\]
da $X(x) = \sum X^j(x)\partial_j(x) = 0 \Leftrightarrow X^j(x) = 0$. Sei nun $X,X' \in \Gamma(TM)$ mit $X(x) = X'(x) = v$. So gilt $S_x(v) = \Lambda(X)(x) = \Lambda(X')(x)$, da $\Lambda(X-X')(x) = 0 = \Lambda(X)(x) - \Lambda(X')(x)$.
\item $S$ also definiert als Schnitt mit:  Ist $X\in \Gamma(TM)$, dann $S_x(X(x)) = \Lambda(X)(x)$. Also gilt lokal: $S\vert_U = \sum_j \Lambda(\partial_j)\mathd x^j$ glatt, da $\Lambda$ glatt.
\end{enumerate}
\end{itemize}
% subsubsection 178(end)

\subsubsection{Lemma}
\label{ssub:179}
Sei $E\to B$ ein Vektorbündel, $x\in E_x$. Dann existiert ein glatter Schnitt $s_v\in \Gamma(E)$ mit $s_v(x) = v$. Anders gesagt: Jeder Schnitt $s:\set{x}\to E_x$ über einen Punkt kann glatt auf ganz $M$ fortgesetzt werden.
\minisec{Beweis}
Übung
% subsubsection 179(end)

\subsubsection{Bemerkung}
\label{ssub:180}
In diesem Sinne sind $C^\infty (M)$-lineare Abbildung nicht nur lokal (wie zum Beispiel $\benbrace{.,.}$), sondern punktweise bestimmte Objekt, das heißt zum Beispiel für $\Lambda:\Gamma(TM)\to C\infty(M)$ $C^\infty(M)$-linear hängt $\Lambda(X)(x)$ nur von $X(x)$ ab.
% subsubsection 180(end)

\subsubsection{Definition}
\label{ssub:181}
Sei $f:M\to N$ glatte Abbildung, $S\in \Gamma(\mathcal{T}_q N)$. Der \Index{Pullback} von $S$ unter $f$ ist der Tensor $f^* S \in \Gamma(\mathcal{T}_q M)$ definiert durch
\[
\underset{\in L_q \enbrace*{\bigtimes^q T_x M,\mathds{R}}}{\enbrace*{f^* S}_x} \enbrace{v_1,\ldots,v_q} := \underset{\in L_q \enbrace*{\bigtimes^q T_{f(x)} N,\mathds{R}}}{S_{f(x)}} \enbrace{\mathd_x f(v_1),\ldots,\underbrace{\mathd_x f(v_p)}_{T_x M \to T_{f(x)} N}}
\]
% subsubsection 181(end)

\subsubsection[Satz:Funktionelle Eigenschaften]{Satz}
\label{ssub:182}
\begin{enumerate}[(i)]
\item Ist $f:M\to N,g:M\to N$ glatt, so ist $(g\circ f)^* = f^*\circ g^*$. ($f^*:\Gamma(\mathcal{T}_q N)\to \Gamma(\mathcal{T}_q M),g^*:\Gamma(\mathcal{T}_q P)\to \Gamma(\mathcal{T}_q N)$)
\item Ist $S\in \Gamma(\mathcal{T}_q N), T\in \Gamma(\mathcal{T}_q N)$, so gilt
\[
f^*(\underbrace{S\otimes T}_{\in \mathcal{T}_{q+q'} N}) = f^* S \otimes f^* T
\]
\end{enumerate}
% subsubsection 182(end)

\subsubsection{Bemerkung}
\label{ssub:183}
\begin{enumerate}[(i)]
\item Ist $f:M\to N$ ein Diffeomorphismus, so kann man den \Index{Pullback} auf kontravarianten Tensoren via dem \Index{Push-Forward} definieren, nämlich ist zum Beipsiel $X\in \Gamma(TN)$, so ist: $f^*X:= (f^{-1})_* X \in \Gamma(TM)$. Also $f^*:\Gamma(\mathcal{T}_q^p N)\to \Gamma(\mathcal{T}_q^p M)$.
\item $f^*$ kommutiert mit $C_j^i$-Kontraktion
\end{enumerate}
% subsubsection 183(end)

\subsubsection{Verkleben von Tensoren}
\label{ssub:184}
Sei $\enbrace{U_i,\phi_i}$ Karten und $S\in \Gamma\enbrace*{\mathcal{T}_q^p M}$ Tensorfeld. Dann gilt für $S_i = \enbrace*{\phi_i^{-1}}^* S\in \Gamma\enbrace*{\mathcal{T}_q^p U_i}$
\[
\enbrace*{\phi_i \circ \phi_j^{-1}}^* S_j \Big\vert_{\phi_j(U_j\cap U_i)} = S_i \Big\vert_{\phi_i(U_j\cap U_i)}
\]
Sind Tensorfelder $S_i$ gegeben, die diese Bedingung erfüllen, so lassen sie sich eindeutig zu einem Tensorfeld auf $M$ verkleben.
% subsubsection 184(end)

Setzt damit nun $\mathcal{L}_X$ "Lie-Ableitung" (definiert für $f\in C^\infty(M), Y\in \Gamma(TM)$) auf $\Gamma(\mathcal{T}_q^p M)$ fort.

\subsubsection{Definition}
\label{ssub:185}
Sei $X\in \Gamma(TM)$. Dann ist
\[
\mathcal{L}_X:\Gamma(\mathcal{T}_q^p M) \to \Gamma(\mathcal{T}_q^p M)
\]
definiert durch
\[
\mathcal{L}_X S:= \frac{\mathd}{\mathd t} \underbrace{\Psi^*_{X,t}}_{\text{Fluss von } X} S\vert_{t=0}
\]
(Beachte $\Psi_{X,t} S$ Kurve in $\Gamma(\mathcal{T}_q^p) M$)
% subsubsection 185(end)

\subsubsection[Satz:Eigenschaften von $\mathcal{L}_X$]{Satz}
\label{ssub:186}
\begin{enumerate}[(i)]
\item $\mathcal{L}_X f = \mathd f(X) = \underset{f\text{ in Richtung }X\text{ ableiten}}{X. f} \in C^\infty(M),\quad f\in C^\infty (M)$
\item $\mathcal{L}_X Y = \benbrace{X,Y},\quad Y\in \Gamma(TM)$
\item $\mathcal{L}_X (S\otimes T) = \mathcal{L}_X S \otimes T + S\otimes \mathcal{L}_X T$
\item $\mathcal{L}_X(C_j^i(S)) = C_j^i(\mathcal{L}_X S)$
\item $\underbrace{\enbrace{\mathcal{L}_X \alpha}}_{\in \Gamma(T^* M)}(Y) = \underbrace{\mathcal{L}_X(\alpha(Y))}_{\in C^\infty (M)} - \underbrace{\alpha(\mathcal{L}_X Y)}_{\in C^\infty (M)} = X.(\alpha(Y)) - \alpha\enbrace{\benbrace{X,Y}}$
\end{enumerate}
\minisec{Beweis}
\begin{enumerate}[(i)]
\item $\enbrace*{\mathcal{L}_X f}(x) = \frac{\mathd}{\mathd t}\enbrace*{\Psi_{X,t}^* f}(x)\vert_{t = 0} = \frac{\mathd}{\mathd t} f\enbrace*{\Psi_{X,t}(x)}\vert_{t = 0} = \mathd f\enbrace*{\frac{\mathd}{\mathd t} \Psi_{X,t}(x)\vert_{t = 0}} = \mathd f\enbrace[\big]{X\enbrace*{\Psi_{X,0}(x)}} = \mathd f(X)(x)$
\item Theorem \ref{ssub:167}
\item Leite die Identifikation $\Psi_{X,t}^*(S\otimes T) = \Psi_{X,t}^* S \otimes \Psi_{X,t}^* T$ und $C_j^i(\Psi_{X,t}^* S) = \Psi_{X,t}^* C_j^i(S)$ ab.
\item Siehe (iii)
\item Sei $\alpha\in \Gamma(T^* M)$, Dann ist $C_1^1(Y\otimes \alpha) = \alpha(Y) \in C^\infty (M)$.
\begin{align*}
X.(\alpha(Y)) &= \mathcal{L}_X (\alpha(Y)) =\mathcal{L}_X \enbrace*{C_1^1(Y\otimes \alpha)} = C_1^1 \mathcal{L}_X \enbrace{Y\otimes \alpha}\\
&= C_1^1 \enbrace{\mathcal{L}_X Y \otimes \alpha} + C_1^1\enbrace{Y\otimes \mathcal{L}_X \alpha} = \alpha(\mathcal{L}_X Y) + (\mathcal{L}_X \alpha)(Y)
\end{align*}
\end{enumerate}
% subsubsection 186(end)

\subsubsection{Korollar}
\label{ssub:187}
\begin{enumerate}
\item $\mathcal{L}_X$ ist lokaler Operator, der durch die Eigenschaften (i) - (iv) in \ref{ssub:187} eindeutig bestimmt ist.
\item $\mathcal{L}_{\benbrace{X,Y}} = \benbrace{\mathcal{L}_X,\mathcal{L}_Y}$ auf $\Gamma\enbrace*{\mathcal{T}_q^p M}$ (wissen dies bereits auf $\Gamma\enbrace*{\mathcal{T}_0^0 M} = C^\infty (M)$)
\end{enumerate}
% subsubsection 187(end)
% subsection 110(end)

\subsection{Äußere Formen (Differentialformen)}
\label{sub:111}

\subsubsection{Äußeres Produkt von Vektorräumen}
\label{ssub:188}
Sei $V$ ein $n$-dimensionaler Vektorraum. Auf $\otimes^p V$ definieren wir den \Index{Antisymmetrisierungsoperator} $A_p$ durch
\[
A_p \enbrace{v_1\otimes \ldots \otimes v_p} = \sum_{\sigma \in S_p} \sgn(\sigma) v_{\sigma(1)}\otimes \ldots \otimes v_{\sigma(p)}
\]
Zum Beispiel:
\begin{align*}
A_2(v\otimes w) = &v\otimes w - w\otimes v\\
A_3(v\otimes w\otimes u) = &v\otimes w\otimes u + u\otimes v \otimes w + w\otimes u \otimes u\\
 -& v\otimes u \otimes w - u\otimes w \otimes v - w\otimes v \otimes u
\end{align*}
Definiere $\Lambda^p V = \im A_P \subset \otimes^p V$. Sei $(e_i)$ eine Basis von $V$, so hat $\Lambda^p V$ die Basis $\enbrace{\underbrace{e_{i_1}\wedge \ldots \wedge e_{i_p}}_{= A_p\enbrace{e_{i_1}\otimes \ldots \otimes e_{i_p}}}}_{1\leq i_1 < \ldots < i_p \leq n}$
\minisec{Spezialfälle}
\begin{itemize}
\item $\Lambda^0 V = \mathds{R}$
\item $\Lambda^1 V = V$
\end{itemize}
\marginnote{$0 \equiv (v+w)\otimes (v+w) = v\otimes v + w\otimes v + v\otimes w + w\otimes w \equiv v\otimes w + w\otimes v$}
Alles auf einmal: $\Lambda^. V := \bigoplus\limits_p \Lambda^p V$, alternativ: $\Lambda^. V = \nicefrac{\mathcal{T}^. V}{\text{Ideal erzeugt von } v\otimes v}$. Auf $\Lambda^. V$ definieren wir das \Index{äußere (Wedge-, Dach-, Keil-,)Produkt}
\[
\wedge: \Lambda^p V\times \Lambda^q V \to \Lambda^{p+q} V,\quad (a,b) \mapsto a\wedge b := \frac{1}{p!q!} A_{p+q}(a\otimes b)
\]
Das Wedgeprodukt ist
\begin{itemize}
\item bilinear
\item assoziativ: $a\wedge (b\wedge c) = (a\wedge b)\wedge c$
\item graduiert kommutativ: $a\wedge b = (-1)^{pq} b\wedge a$, wenn $a \Lambda^p V, b \Lambda^q V$
\end{itemize}
\minisec{Beispiel}
Sei $\alpha \in \Lambda^1 V^*, \beta \in \Lambda^1 V^*$, so ist mit $(\alpha \otimes \beta)(v,w) = \alpha(v)\beta(w)$:
\begin{align*}
(\alpha \wedge \beta)(v,w) &= (\alpha\otimes \beta - \beta \otimes \alpha)(v,w)\\
&= \alpha(v)\beta(w) - \beta(v)\alpha(w)\\
&= -\enbrace{\alpha(w)\beta(v) - \beta(w)\alpha(v)}\\
&= - (\alpha \wedge \beta)(w,v)
\end{align*}
Für $\alpha: V\to V$ lineare Abbildung erhalten wir eine lineare Abbildung
\[
\wedge^p \alpha: \Lambda^p V \to \Lambda^p V\ \text{durch}\ \wedge^p \alpha\enbrace{v_1\wedge \ldots \wedge v_p} = \alpha(v_1)\wedge \ldots \wedge \alpha(v_p)\ \text{für}\ v_i\in V
\]
\minisec{Übungsaufgabe}
\begin{itemize}
\item $\Lambda^p V^* = (\Lambda^p V)^*$
\item $\wedge^n \alpha$ ist die Multiplikation mit der Determinante von $\alpha$
\end{itemize}
Für die Dimension von $\Lambda^p V$ gilt: $\dim \Lambda^p V = \binom{n}{b}$ (somit $\dim \Lambda^0 V = 1,\ \dim \Lambda^1 V = n, \dim \Lambda^n V = 0$).
% subsubsection 188(end)

\subsubsection{Definition}
\label{ssub:189}
Das Bündel der äußeren $p$-Formen (Differentialformen) ist das Vektorbündel $\Lambda^p T^* M$. Die glatten Schnitte heißen (äußere) $p$-Formen (Differentialformen), wir schreiben $\Omega^p (M) = \Gamma(\Lambda^p T^* M)$. Wir haben $\Lambda^0 T^* M = M\times \mathds{R},\ \Omega^0(M) = C^\infty (M)$. In lokalen Koordinaten $(U,x^i)$:
\[\alpha \in \Omega^p (M) \leadsto \alpha\vert_U = \sum_{1\leq i_1 < \ldots < i_p \leq n} \alpha_{i_1,\ldots,i_p} \mathd x^{i_1} \wedge \ldots \wedge x^{i_p}\]
% subsubsection 189(end)

\subsubsection{Bemerkung}
\label{ssub:190}
Sei $f:M\to N$, so ist $f^*:\Omega^p(N) \to \Omega^p(M)$ definiert durch
\[
(f^* \omega)(v_1,\ldots,v_p) = \omega\vert_{f(x)} \enbrace*{\mathd f\vert_{x} (v_1),\ldots, \mathd f\vert_{x} (v_p)}
\]
Es gilt:
\begin{itemize}
\item $\enbrace*{f\circ g}^* g^* \circ f^*$
\item $f^*(\alpha \wedge \beta) = (f^* \alpha) \wedge (f^* \beta)$ für $\alpha,\beta$ Differentialformen
\item $\mathcal{L}_X \alpha$ ist wieder eine Differentialformen
\[
\mathcal{L}_X (\alpha \wedge \beta) = \mathcal{L}_X \alpha \wedge \beta + \alpha \wedge \mathcal{L}_X \beta
\]
\item Es gilt das "innere Produkt": (innere Ableitung)
\[
i_X:\Omega^p (M)\to \Omega^{p-1} (M),\quad \omega \mapsto \omega (X,.)
\]
Für $\alpha \in \Omega^p (M), \beta \in \Omega^q (M)$ ist 
\[
i_X(\alpha \wedge \beta) = (i_X \alpha)\wedge + (-1)^p \alpha \wedge i_X(\beta)
\]
\end{itemize}
% subsubsection 190(end)

\subsubsection{Satz/Definition}
\label{ssub:191}
Auf $\Omega^p(M)$ gibt es einen lokalen linearen Operator $\mathd : \Omega^p (M) \to \Omega^{p+1} (M)$, der durch folgende Eigenschaften eindeutig bestimmt ist:
\begin{itemize}
\item Für $p = 0$ ist $\mathd: C^\infty (M) \to \Omega^1 (M)$ das gewöhnliche Differential
\item Für $f\in C^\infty (M)$ ist $\mathd(\mathd f) = 0$
\item Für $\alpha\in \Omega^p(M), \beta \in \Omega^p(M)$ ist:
\[
\mathd(\alpha \wedge \beta) = (\mathd \alpha) \wedge \beta + (-1)^p \alpha \wedge (\mathd \beta)
\]
\end{itemize}
Beweis:
\begin{itemize}
\item Eindeutigkeit: Sei $\enbrace{U,x^1,\ldots,x^n}$ Karte, so gilt
\begin{align*}
(\mathd \alpha)\vert_U \overset{\text{lokal}}{=} \mathd(\alpha\vert_U) &\overset{\text{linear}}{=} \sum_{1\leq i_1,\ldots,i_p,\leq n} \mathd \enbrace*{\alpha_{i_1,\ldots, i_p} \wedge \mathd x^{i_1} \wedge \ldots \wedge x^{i_p}}\\
&\overset{(i),(ii),(iii)}{=} \sum_{1\leq i_1,\ldots,i_p,\leq n} \underbrace{\mathd\alpha_{i_1,\ldots, i_p}}_{\text{bekannt (i)}} \wedge \mathd x^{i_1} \wedge \ldots \wedge x^{i_p}\quad (\star)
\end{align*}
\item Lokal definieren durch $(\star)$ (also auf $\Omega^p (V), V\overset{\text{offen}}{\subset} \mathds{R}^n$. Dann ist zu zeigen, dass dies die Eigenschaften erfüllt.
\item $\phi: V\subset \mathds{R}^n \to U \mathds{R}^n$ glatt $\leadsto$ $\phi^* \circ \mathd = \mathd \circ \phi^*$
\item Zusammenkleben mit Bemerkung \ref{ssub:184} um $\mathd$ auf ganz $M$ fortzusetzen.
\end{itemize}
% subsubsection 191(end)

\subsubsection{Korollar}
\label{ssub:192}
\begin{itemize}
\item $\mathd \circ \mathd = 0$
\item Für $f:M\to N$ ist $\mathd \circ f^* = f^* \circ \mathd$
\end{itemize}
% subsubsection 192(end)

\subsubsection{Bemerkung}
\label{ssub:193}
\begin{enumerate}
\item Wir erhalten eine Folge
\[
0 \to \Omega^0 (M) \xrightarrow{\mathd} \Omega^1 (M) \xrightarrow{\mathd} \ldots \xrightarrow{\mathd} \Omega^n (M) \xrightarrow{\mathd} 0
\]
$\leadsto$ Kettenkomplex
\item Es gibt auch eine explizite Formel ohne Koordinate für $\mathd$ für $\alpha \in \Omega^1 (M):$
\[
\mathd \alpha(X,Y) = X.\alpha(Y) - Y.\alpha(X) - \alpha\enbrace{\benbrace{X,Y}}
\]
\end{enumerate}
% subsubsection 193(end)

\subsubsection{Beispiel}
\label{ssub:194}
Für $\alpha \in \Omega^1 (M),\ \alpha = \sum \alpha_j \mathd x^j$ gilt
\[
(\mathd \alpha)\vert_U = \mathd(\alpha\vert_U) = \sum_{1\leq i_1 < \ldots <i_p \leq n } \sum_{j} \enbrace{\partial_j \alpha_{i_1,\ldots, i_p}} \mathd x^j \wedge \mathd x^{i_1} \wedge \ldots \wedge \mathd x^{i_p}
\]
% subsubsection 194(end)

\subsubsection[Proposition: Cartansche Formel]{Proposition}
\label{ssub:195}
Für $X\in \Gamma(TM)$:
\[
\mathcal{L}_X = \mathd \circ i_X + i_X \circ \mathd
\]
\minisec{Beweis}
Definiere $P_X = \mathd \circ i_X + i_X \circ \mathd$. Die ist eine linearer lokaler Operator, da Verknüpfung von lokalen linearen Operatoren. Sei $\alpha \in \Omega^p (M)$, so ist $\alpha$ von der Form:
\[
\alpha\vert_U = \sum \alpha_{i_1 \ldots i_p} \mathd x^1 \wedge \ldots \wedge \mathd x^p
\]
Es genügt zu zeigen:
\begin{enumerate}[(i)]
\item $\mathcal{L}_X = P_X$ gilt für $0$-Formen und exakte $1$-Formen 
\marginnote{$w$ exakt $\Leftrightarrow$ $\exists\ \alpha:\omega = \mathd \alpha$}
\item rechte Seite erfüllt Leibnizregel (selber Rechnen)
\end{enumerate}
\begin{enumerate}[(i)]
\item Für $f\in \Omega^0 (M)$: $P_X f = (i_X \circ \mathd)f = \mathd f (X) = \mathcal{L}_X f$
\item Für $\mathd f \in \Omega^1 (M)$: $P_X(\mathd f) = \mathd i_X \mathd f + i_X \underbrace{\mathd \mathd}_{0} f$.
Sei $Y\in \Gamma(TM)$, so gilt $(P_X \mathd f)(Y) = (\mathd i_X \mathd f)(Y) = \mathd \enbrace*{\mathd f(X)}(Y) = Y.X.f$. Außerdem gilt
\[
\enbrace*{\mathcal{L}_X \mathd f}(Y) = \mathcal{L}_X \enbrace*{\mathd f (Y)} = \mathd f \enbrace*{\mathcal{L}_X Y} = X.Y.f - \benbrace{X,Y}.f = Y.X.f
\]
\end{enumerate}
% subsubsection 195(end)

% subsection 111(end)

\subsection{Volumenformen und Integration}
\label{sub:112}

\subsubsection{Definition}
\label{ssub:197}
Eine \Index{Volumenform} ist ein $\omega \in \Omega^n (M)$ ($n = \dim M$) mit $\omega_x \neq 0\ \forall x\in M$
% subsubsection 197(end)

\subsubsection{Bemerkung}
\label{ssub:198}
Aus der Existenz einer Volumenform folgt $\Lambda^n T^* M = M\times \mathds{R}$. Das ist nicht immer der Fall, das heißt nicht auf jedes Mannigfaltigkeit existiert eine Volumenform.
% subsubsection 198(end)

\subsubsection{Integration auf Mannigfaltigkeiten}
\label{ssub:199}
Sei $\omega \in \Omega^n (M^n)$ Volumenformen, $\enbrace{U,\phi}$ Karte, $\phi^{-1*} \omega \vert_U = \omega_U \mathd x^1 \wedge \ldots \wedge \mathd x^n$. Sei $f\in C_c^0 (M)$ mit $\supp(f) \subset U$. Setze 
\[
\Lambda_\omega (f) = \int_{\phi(U)} \phi^{-1*} f \omega_U \underbrace{\mathd x^1 \wedge \ldots \wedge \mathd x^n}_{\mathd L^n}
\]
Ist dies unabhängig von der Wahl der Karte? Dazu sei $\enbrace{V,\psi}$ Karte, $\supp f \subset U\cap V$ mit $\psi^{-1*} \omega = \omega_V  \mathd x^1\wedge \ldots \wedge \mathd x^n$. Setze $F := \phi \circ \psi^{-1}\vert_{\psi(U\cap V)}$ und wir nehmen zunächst an, dass $\det \mathd F > 0$. Es gilt
\marginnote{Es gilt $\omega_V = \det \mathd F (\omega_U \circ F)$}
\begin{align*}
\int_{\psi(U\cap V)} (f\circ \psi^{-1}) \cdot \omega_V \mathd L^n &= \int_{\psi(U\cap V)} (f\circ \phi^{-1}\circ F)\cdot \det F \cdot \omega_U \circ F \mathd L^n\\
&= \int_{F\circ \psi(U\cap V)} (f\circ \phi^{-1}) \cdot \omega_U \mathd L^n\\
&= \int_{\phi(U\cap V)} (f\circ \phi^{-1}) \cdot \omega_U \mathd L^n
\end{align*}
Also ist $\Lambda_\omega (f)$ unabhängig von der Wahl der Karten. Für beliebige $f\in C_c^0 (M)$ wähle Atlas $\enbrace{U_i,\phi_i}$ und eine untergeordnete Teilung der Eins $(\rho_i)$, das heißt $\supp \rho_i \in U_i$. Dann setze
\[
\Lambda_\omega f := \sum_i \Lambda_w (\rho_i \cdot f)
\]
Die Unabhängigkeit von $\enbrace{\rho_i}$ und dem Atlas ist eine Übungsaufgabe. Mit dem Rieszschen Darstellungssatz folgt die Existenz des Radon-Maß $\omega$ auf $M$ mit $\Lambda_\omega f = \int_M f \mathd \omega$.
% subsubsection 199(end)

\subsubsection[Bemerkung: Volumen von kompakten Mannigfaltigkeiten]{Bemerkung}
\label{ssub:1100}
Wenn $M$ kompakt und alle Jacobideterminanten der Übergangsfunktion positiv sind, setze 
\[
\Vol_\omega (M) = \int_M \omega
\]
% subsubsection 192(end)

\subsubsection{Definition}
\label{ssub:1101}
$M$ heißt \Index{orientierbar}, falls es einen Atlas mit positiven Jacobideterminanten der Übergangsfunktionen gibt. So ein Atlas heißt \Index{Orientierungsatlas} von $M$.
% subsubsection 1101(end)

\subsubsection{Beispiel}
\label{ssub:1102}
$TM$ ist immer orientierbar. Der angegebene Atlas hat Übergangsfunktion 
\[
\Phi_{ij} (x,v) = \enbrace*{\phi_i \circ \phi_j^{-1} (x), \mathd _{\phi_j(x)}\enbrace*{\phi_i \circ \phi_J^{-1}}(v)}
\]
, das heißt
\[
\mathd_{(x,v)} \Phi_{ij} = \begin{pmatrix}
\mathd_x \phi_{ij} & *\\
0 & \mathd_x \phi_{ij}
\end{pmatrix}
\]
Somit folgt, dass $\det \mathd_{(x,v)} \Phi_{ij} > 0$.
% subsubsection 1102(end)

\subsubsection{Proposition}
\label{ssub:1103}
$M$ ist orientierbar, genau dann wenn eine Volumenform auf $M$ existiert.
\minisec{Beweis}
\begin{itemize}
\item "$\Rightarrow$": Sei $\omega$ Volumenform. Wähle Atlas mit $\enbrace{U_i,\phi_i}$ mit $\phi_i^{-1*} \omega = \omega_i \mathd x^1 \wedge \ldots \wedge \mathd x^n$, sodass $\omega_i > 0$. Das ist möglich, da $\omega_i(x) \neq 0\ \forall x\in U_i$ und man für $\omega_i < 0$ die Karte $\phi_i$ mit der Abbildung $\mathds{R}^n \to \mathds{R}^n, 
\begin{psmallmatrix}
-1&&&\\
&1&&\\
&&\ddots&\\
&&&1
\end{psmallmatrix}$ verknüpfen kann um eine neue Karte mit $\omega_i' > 0$ zu erhalten. Mit $\omega_V = \det \mathd F(\omega_U \circ F)$ folgt, dass die Jacobideterminante der Übergangsfunktion größer $0$ ist.
\item "$\Leftarrow$": Wähle Orientierungsatlas $\enbrace{U_i,\phi_i}$ und untergeordnete Teilung der Eins $\enbrace{\rho_i}$. Setze $\omega_i = \phi_i^*\enbrace{\mathd x^1 \wedge \ldots \wedge \mathd x^n}$. Dann ist $\omega := \sum_i \rho_i \omega_i$ Volumenform auf $M$, da
\begin{align*}
\phi_k^{-1*} \omega &= \sum_i \rho_i \phi_k^{-1*} \phi_i^* \enbrace{\mathd x^1 \wedge \ldots \wedge \mathd x^n}\\
&= \sum_i \rho_i \enbrace{\phi_i\circ \phi_k^{-1}}^* \enbrace{\mathd x^1 \wedge \ldots \wedge \mathd x^n}\\
&= \underbrace{\sum_i \rho_i \underbrace{\det \mathd \phi_{ik}}_{0}}_{\geq \min_{\rho_i \neq 0} \det \mathd \phi_{ik} \underbrace{\sum \rho_i}_{=1} > 0} \enbrace{\mathd x^1 \wedge \ldots \wedge \mathd x^n}\\
\end{align*}
\end{itemize}
% subsubsection 1103(end)

% subsection 112(end)

% section 1 (end)

\section{Riemannsche Geometrie}
\label{sec:2}

\subsection{Riemannsche Metriken}
\label{sub:21}

\subsubsection[Definition: Riemannsche Metrik]{Definition}
\label{ssub:21}
\Index{Riemannsche Metrik} $g$ ist ein symmetrischer positiv definiter $(0,2)$ Tensor auf $M$, das heißt $g\in \Gamma(\mathcal{T}_2 M)$ (beziehungsweise $g: \Gamma(TM)\times \Gamma(TM) \to C^\infty (M)$) mit
\begin{itemize}
\item $g(X,Y) = g(Y,X)\quad \forall X,Y\in \Gamma(TM)$ (Symmetrie)
\item $g_x(v,v) \geq 0$ mit Gleichheit genau dann wenn $v = 0$
\end{itemize}
das heißt $\enbrace{T_x M,g_x}$ ist ein euklidischer Vektorraum.
\begin{itemize}
\item $\enbrace{M,g}$ ist eine Riemannsche Mannigfaltigkeit
\item $\varphi \in \Diff (M) = \set{\varphi: M\to M\ \text{Diffeomorphismus}}$ heißt \Index{Isometrie} von $\enbrace{M,g}$, falls $\varphi^* g = g$
\item Ist $\varphi: M\to M$ lokaler Diffeomorphismus mit $\varphi^* g = g$ heißt $\varphi$ lokale Isometrie
\item $\Isom\enbrace{M,g} =$ Gruppe der Isometrien
\item $\enbrace{M,g},\enbrace{N,h}$ Riemannsche Mannigfaltigkeiten heißen isometrisch, falls ein Diffeomorphismus $\varphi: M\to N$ existiert mit $\varphi^* h = g$
(insbesondere ist dann $\enbrace{T_{\varphi(x)} N,h_{\varphi(x)}}$ isometrisch als euklidischer Vektorraum zu $\enbrace{T_x M,g_x}$
\end{itemize}
%subsubsection 21(end)


\subsubsection{Bemerkung}
\label{ssub:22}
Sind $x^1,\ldots,x^n$ lokale Koordinaten, so gilt lokal
\[
g\vert_U = \sum_{i,j} g_{ij} \mathd x^i \otimes \mathd x^j,\quad \enbrace{g_{ij}(x)}_{i,j = 1}^n\ \text{sym. pos. def. Matrix für alle}\ x\in U
\]
%subsubsection 22(end)

\subsubsection{Beispiele}
\label{ssub:23}
\begin{enumerate}[(i)]
\item $\enbrace{\mathds{R}^n,g_{euk} = \sum_i \mathd x^i \otimes \mathd x^i}$ "flache" oder "euklidische Metrik".
\item Sei $M\subset \mathds{R}^n$ Untermannigfaltigkeit, allgemein $M\subset (N,h)$ Untermannigfaltigkeit, so ist die \Index{induzierte Metrik}
\[
g = h\vert_U :\underset{\subset \Gamma(TN)\vert_M}{\Gamma(TM)} \times \underset{\subset \Gamma(TN)\vert_M}{\Gamma(TM)} \to C^\infty (M)
\]
definiert durch $g\enbrace{\underbrace{X,Y}_{\in \Gamma(TM)\times \Gamma(TM)}} = h\enbrace{X,Y}\in C^\infty (M)$. (Vergleiche $V\subset (W,g)$ euklidischer Vektorraum $\leadsto$ $\enbrace{V,g\Vert_V: V\times V \to \mathds{R}}$ euklidischer Vektorraum)\\
Zum Beispiel: $S^n \subset \mathds{R}^{n+1}, g_{euk}\vert_{S^n} = g_{rou}$ "runde" Metrik
\item Betrachte $S^n \times \mathds{R}_{>0} \simeq \mathds{R}^{n+1}\backslash \set{0}$. 
%\begin{figure}[H]
%\centering{
%\begin{tikzpicture}[scale = 0.6]
%	%\draw[help lines] (0,0) grid (7,6);
%	\draw[thick,->] (0,3) -- (3,3) -- (6,3);
%	\draw[thick,->] (3,0) -- (3,3) -- (3,6);
%	\draw[dashed] (3,3) circle[radius = 2];
%	\node[below right] at (1,5){$S^1$};
%	\draw (3,3) -- (5.5,5.5) node[below right] {$\mathds{R}_{>0}$};
%	\draw[red] (4,3) node[above left]{$\theta$} arc (0:45:1);
%	\draw[green] (3,3) -- (4,4) node[above]{$r$} -- (5,5);
%\end{tikzpicture}
%\caption{$S^1\times \mathds{R}_{>0} \simeq \mathds{R}^2\backslash\set{0}$}
%}
%\end{figure}
$T_{\theta,r} \enbrace{S^n\times \mathds{R}_{>0}} = T_\theta S^n \oplus T_r \mathds{R}_{>0}$ mit Metrik $r^2 g_{rou} + \mathd t^2$ $\cong$ $g_{euk}\vert_{\mathds{R}^{n+1}\backslash \set{0}}$
\item Sei $f:M\to N$ Immersion in $x\in M$, das heißt $\mathd_x f: T_x M \to T_{f(x)} N$ ist injektiv. Ist $h$ Metrik auf $N$, so ist $f^* h$ Metrik auf einer Umgebung von $x$. Nämlich
\begin{itemize}
\item Symmetrie: $f^* h \in \Gamma(\mathcal{T}_2 M)$ und $\enbrace{f^* h}_x (v,w) = h_{f(x)} \enbrace{\mathd f_x (v),\mathd f_x (w)} = h_{f(x)} \enbrace{\mathd f_x (w),\mathd f_x (v)} = \enbrace{f^* h}_x (w,v)$
\item Positiv Definitheit: $\enbrace{f^* h}_x (v,v) = h_{f(x)} \enbrace{\mathd f_x (v),\mathd f_x (v)} \leq 0$
\item $h_{f(x)} \enbrace{\mathd f_x (v),\mathd f_x (v)} = 0 \Rightarrow \mathd f_x (v) = 0 \xRightarrow{f\ \text{Immersion}} v = 0$
\end{itemize}
\end{enumerate}
%subsubsection 23(end)

\subsubsection{Allgemeine Konstruktion}
\label{ssub:24}
\begin{enumerate}[(i)]
\item Eine Abbildung $p:M\to N$ heißt Überlagerung von $N$, falls for alle $y\in N$ eine Umgebung $y\in V\subset N$ existiert, sodass $p^{-1}(V) = \dot{\bigcup_j}U_j$ disjunkte Vereinigung offener $U_j \subset M$ mit $p\vert_{U_j}$ ist Diffeomorphismus auf sein Bild.
\missingfigure{$S^n = \nicefrac{\mathds{R}}{2\pi \mathds{Z}}$}
Sei $h$ eine Metrik auf $N$. Dann exisiert eine eindeutig bestimmte Metrik $g$ auf $M$, sodass $p$ eine lokale Isometrie wird, nämlich $g = p^* h$. Im Allgemeinen nennt man eine Überlagerung $p:(M,g)\to (N,h)$ Riemannsch, falls $p$ lokale Isometrie. Insbesondere existiert eine kanonische Riemannsche Metrik auf der sogenannten "universellen Überlagerung" $\tilde{N}$ einer Riemannschen Mannigfaltigkeit. 
\item Sei $G\subset \Diff(M)$ Untergruppe.
\[
G\times M\to M, (f,x)\mapsto f(x)\quad \text{Gruppenwirkung}
\]
$G$ wirkt \Index{eigentlich}, falls für alle $x,y \in M$ gilt $x\not\in G.y = \set[f(y)]{f\in G}$ "Bahn von $y$ unter $G$". (zum Beispiel $S^1 \ni e^{i\theta}$ wirkt auf $y\in \mathds{C}\simeq \mathds{R}^2$ durch Multiplikation $e^{i\theta} y$
\missingfigure{Wirkung von $S^1$ auf $\mathds{R}^2$}
eine Umgebung $U$ von $x$ mit $U\cap G.y = \emptyset$
\minisec{Bemerkung} Diese Eigenschaft ist für nicht kompakte Gruppen interessant, zum Beispiel falls es eine dichte Bahn gibt.
\missingfigure{Dichte Bahn}
Die Wirkung ist \Index{diskontinuierlich}, falls $\forall x\in M$ eine offene Umgebung $x\in U_x$ existiert mit $f(U_x)\cap U_x = \emptyset$ außer für $f = \id_M$
\missingfigure{Diskontinuierliche Wirkung}
Wirkt $G$ auf $M$ eigentlich diskontinuierlich, dann ist $\nicefrac{M}{G}$ eine Mannigfaltigkeit, sodass $p:M\to \nicefrac{M}{G}, x\mapsto Gx$ eine Überlagerung ist.
\minisec{Beispiel}
$\mathds{Z}^n$ wirkt auf $\mathds{R}^n$ durch Diffeomorphismus, nämlich
\[
\enbrace{m^1,\ldots,m^n}\in \mathds{Z}^n, \enbrace{x^1,\ldots,x^n}\in \mathds{R}^n \mapsto \enbrace{x^1+m^1,\ldots,x^n+m^n}
\]
\missingfigure{$\nicefrac{\mathds{R}^2}{\mathds{Z}^2}$}
\minisec{Bemerkung}
Daraus, dass die Wirkung eigentlich ist, folgt, dass $\nicefrac{M}{G}$ Hausdorff ist.
\missingfigure{Darstellung der Hausdorffeigenschaft}
Aus der Diskontinuierlichkeit der Wirkung, folgt die Existenz der Überlagerung $\mathds{Z}^n \times \mathds{R}^n \to \mathds{R}^n$
\end{enumerate}
\todo{Hier fehlt noch etwas}
%subsubsection 24(end)

\subsubsection{Satz}
\label{ssub:25}
Auf jeder Mannigfaltigkeit existiert eine Riemannsche Metrik.
\minisec{Beweis}
Sei $\set{(U_i,\varphi_i)}$ Atlas, $\set{\theta_i:M\to\mathds{R}}$ untergeordnete Zerlegung der Eins, $g_i:=\varphi_i^*g_{euc}$ Metrik auf $U_i$. So ist $\theta_i g_i\in\Gamma(\mathcal{T}_2 M)$ durch 0 fortgesetzter symmetrischer (0,2)-Tensor und damit $g=\sum_i \theta_i g_i$ Riemannsche Metrik, denn 
\begin{itemize}
\item $g_x=\sum_{\text{endl.}} \theta_i g_i$ mit $x\in \supp \theta_i$, also ist $g$ wohldefiniert
\item g ist symmetrisch und nicht negativ, da $(\theta_i\underset{\geq 0}{\underbrace{g_i(v,v)}}\geq 0)$
\item $g_x(v,v)=0\Rightarrow \theta_ig_i(v)=0~\forall i\Rightarrow v=0$ (für j mit $v\in U_j\Rightarrow g_j$ Riemannsche Metrik)
\end{itemize}

Metrik induziert mehrere Objekte 
\begin{itemize}
\item kanonischen Isomorphismus $TM\simeq T^*M$
\item kanonische $\nabla^g$ kovariante Ableitung
\item kanonische Distanzfunktion $d^g$
\end{itemize}
%subsubsection 25(end)

\subsubsection{Musikalische Isomorphismen}
\label{ssub:26}
Sei $g$ positiv definit, also insbesondere nicht ausgeartet, d.h. ist $g_x(v,w)=0~\forall w\in T_xM\iff v=0$. So können wir definieren
\[
b:TM\to T^* M,\ b: v\in T_x M\mapsto \enbrace*{g_x (v,.):T_x M\to\mathds{R}\ \text{linear}}\in T_x^* M
\]
$b$ ist Bündelisomorphismus das heißt $b\vert_{T_x M}$ ist linearer Isomorphismus, denn angenommen $b(v)=0$, das heißt $b(v)(w) = g(v,w) = 0\ \forall w\in T_x M$, so folgt $v=0$, das heißt $b$ injektiv und damit $(\dim T_x^* M = \dim T_x M)$ Isomorphismus\\
Die Umkehrabbildung $\#:T^* M\to TM$ definiert durch $(\#\alpha_x)$, das definiert durch die Gleichung $g_x(\#\alpha_x,\omega)=\alpha_x(\omega)$ für $\omega\in T_x M$\\
\minisec{Bezeichnung}
Lokal ist Vextorfeld $X$ bzw. 1-Form $\alpha$ durch 
\[X\vert_U=\sum X^i\partial_i,\ \alpha\vert_U=\sum\alpha_idx^i\]
Also senkt $b$ die Indizes von $X$ ($b(X)$ ist $1-Form$) und $\#$ hebt die Indizes von $\alpha$ ($\#(\alpha)$ ist Vektorfeld).
In der Folge identifiziere $TM$ oft mit $T^*M$, falls Metrik fixiert ist. Zum Beispiel induziert Metrik $g$ eine Metrik $g^*$ auf $T^* M$, in den man $g^*(\alpha,\beta)=g(\#\alpha,\#\beta)$ setzt. 
%subsubsection 26(end)

% subsection 21(end)

\subsection{Kovariante Abbildung}

\subsubsection{Definition}
\label{ssub:27}
Eine \Index{kovariante Abbildung} auf $M$ ist eine $\mathds{R}$-lineare Abbildung
\[
\nabla:\Gamma(TM)\to\Gamma(TM\otimes T^* M) = \Gamma(\End(TM))=\Gamma(\mathcal{T}_1^1 M)
\]
mit $\nabla(f\cdot X) = X\otimes df + f\nabla X$ (Produkt - Regel)
\minisec{Schreibweise}
\[
\nabla_X Y = C_1^1\enbrace{\underbrace{X\otimes \nabla Y}_{\in \Gamma(TM\otimes TM\otimes T^* M) = \Gamma(\mathcal{T}_1^2)}} \in \Gamma (TM)
\]
Es gilt
\[
\nabla Y\in\Gamma(TM\otimes T^*M)\leftrightarrow\nabla Y:\Gamma(T^*M)\times\Gamma(TM)\to C^\infty(M)\leftrightarrow \nabla Y:\Gamma(TM)\to\Gamma(TM)
\]
das heißt $\nabla_XY\in\Gamma(TM)$. Punktweise gilt:
\[
\enbrace{\nabla Y}_x\in T_x M\otimes T^*_x M \simeq \End (T_x M),\ (\nabla_v Y)_x := (\nabla Y)_x(v)\in T_x M
\]
%subsubsection 27(end)

\subsubsection{Beispiel}
\label{ssub:28}
\begin{enumerate}[(i)]
\item Sei $\mathds{R}^n$ mit "flacher" kovarianter Ableitung. Sei $Y\in\Gamma(T\mathds{R}^n)$
\[
\sum Y^i \partial_i = \begin{pmatrix}Y^1\\ \vdots\\Y^n\end{pmatrix}, Y^i:\mathds{R}^n\to\mathds{R}
\]
So ist $\nabla_X Y\in \Gamma(T\mathds{R}^n)$ definiert durch 
\[
\nabla_X Y=\mathd Y(X)=\begin{pmatrix}dY^1(X)\\\vdots\\dY^n(X)\\\end{pmatrix}
\]
$dY(X) = $ Richtungsableitung von Y nach X 
\begin{itemize}
\item $\mathds{R}$-linear und $\nabla Y=\mathd Y(.):\Gamma(T\mathds{R}^n)\to\Gamma(T\mathds{R}^n)$
\item Leibniz-Regel: 
\[
\nabla_X(fY)=(Y\otimes \mathd f)(X) + f\nabla  Y(X) = \mathd f(X)\cdot Y + f\nabla_X Y = df(X)\cdot Y + f\mathd Y(X)
\]
Andererseits $\nabla_X(fY)=d(fY)(X)=df(X)\cdot Y+fdY(X)$ (Produktregel der Ableitung)
\end{itemize}
Intuition der konvarianten Ableitung $\nabla_XY$ als Richtungsableitung von Y nach X. 
\item $\nabla$ induziert kovariante Ableitung $\nabla:\Gamma(\mathcal{T}^p_qM)\to\Gamma(\mathcal{T}^p_{q+1})$ durch 
\begin{itemize}
\item $\nabla_Xf:=df(X),\ f\in C^\infty(M)$
\item $\enbrace{\nabla_X\alpha}(Y)=X.\alpha(Y) - \alpha(\nabla_X Y)$ bestimmt $\nabla_X\alpha\in\Gamma(T^*M)$ gesehen als $C^\infty$-lineare Abbildung $\Gamma(TM)\to C^\infty(M)$
\item $\nabla_X (S\otimes T)=\nabla_X S\otimes T + S\otimes\nabla_X T$
\end{itemize}
\end{enumerate}
%subsubsection 28(end)

\subsubsection{Bemerkung}
\label{ssub:29}
\begin{enumerate}[(i)]
\item Sei $\nabla Y\in\Gamma(TM\otimes T^* M)$ das heißt $\nabla Y: \Gamma(TM)\to \Gamma(TM)$. $\nabla_X Y$ ist $C^\infty (M)$-linear in $X$, das heißt $\nabla_{f X} Y = f\nabla_X Y\ \forall f\in C^\infty(M)$. Anders gesagt $\nabla_X Y$ ist \Index{tensoriell} in X (anders als $\mathcal{L}_X Y$). Insbesondere hängt $(\nabla_X Y)(x)\in T_xM$ nur von $X(x)$ ab! (Und nicht von der Umgebung von $X(x)$)
\item Abhängigkeit von Y? Betrachte lokale Darstellung von $\nabla_X Y:$
\[
X\vert_U=\sum X^i\partial_i, Y\vert_U=\sum Y^i\partial_i\ \text{für lokale Koordianten}
\]
\marginnote{$\Gamma_{ij}^k$ sog. Christoffel-Symbole von $\nabla$}
Es gilt:
\begin{align*}
\enbrace{\nabla_X Y}(x) &= \enbrace{\nabla_X Y\vert_U}(x)\\
&=\sum \enbrace{\nabla_X Y^i}\partial_i + Y^i\nabla_X \partial_i\\
&=\sum_{i,j}\enbrace{X^j \underbrace{\nabla_{\partial_j} Y^i}_{\mathd ^i(\partial_j)}}\partial_i+ X^j Y^i\underbrace{\nabla_{\partial_j}\partial_i(x)}_{=\sum \Gamma_{ij}^k \partial_k \in \Gamma(TM\vert_U)}\\
&=\sum_k \enbrace*{\sum_j\enbrace*{X^j \mathd Y^k(\partial_j)+\sum_i X^j Y^i \Gamma_{i,j}^k}}\partial_k(x)\\
&=\sum_k \enbrace*{\sum_j \enbrace*{\underbrace{X^j(x) \mathd_x Y^k(\partial_j)}_{\mathd_x Y^k\enbrace*{\sum_j X^j(x) \partial_j} = \mathd_x Y^k(X(x))} + \sum_i X^j(x)Y^i(x)\Gamma_{i,j}^k(x)}}\partial_k(x)
\end{align*}
Hängt nur von $X^j(x)$, $Y^j(x)$ und $\mathd_x Y^k(X(x)) = \frac{\mathd}{\mathd t} Y^k(\alpha(t))$, das heißt von $Y$ entlang einer Kurve $\alpha$ mit $\alpha(0) = x, \dot{\alpha}(0) = X(x)$.
\item Werden sehen: Kovariante Ableitungen existieren immer. Wie sieht der Raum der kovarianten Ableitungen aus? Angenommen $\nabla$ ist kovariante Ableitung 
\[
A\in \Gamma(\underbrace{TM\otimes T^* M\otimes T^* M}_{\End(TM\otimes T^* M)}) \hat{=} A:\Gamma(TM)\to\Gamma(\End TM), X\mapsto A_X
\]
$\Rightarrow \nabla^A=\nabla+A$ definiert durch 
\[
\nabla_X^A Y = \nabla_X Y+\underbrace{A_X Y}_{\in \Gamma(TM)}\ \text{, wobei}\ A_X = C_2^1(X\otimes A)
\]
kovariante Ableitung, da 
\[
\nabla_X^A(fY)=\nabla_X(fY) + \underset{\text{tensoriell in X und Y}}{A_X fY}=\mathd f(X)Y+f\nabla_X Y + f A_X Y = \mathd f(X)Y + f\nabla_X^AY
\]
Umgekehrt: Seien $\nabla,\tilde{\nabla}$ zwei kovariante Ableitungen 
\[\nabla-\tilde{\nabla}:\Gamma(TM)\to\Gamma(TM\otimes T^*M)=\Gamma(\End(TM))\]
$C^\infty$-linear? Falls ja, $A = \nabla-\tilde{\nabla}\in \Gamma(\End(TM)\otimes T^* M)$, das heißt $\tilde{\nabla}=\nabla + A = \nabla^A$
Also ist der Raum der kovarianten Ableitungen ist ein affiner Vektorraum $\nabla+\Gamma(\End(TM)\otimes T^* M)$.
Tatsächlich 
\begin{align*}
\nabla_X(fY) - \tilde{\nabla}_X(fY) &= df(X)Y+f\nabla_XY-df(X)Y-f\tilde{\nabla}_XY\\
&= f\nabla_X Y - f\tilde{\nabla}_X Y\\
&= f(\nabla_X Y-\tilde{\nabla}_X Y)
\end{align*}
\end{enumerate}
%subsubsection 29(end)

\subsubsection[Definition: kovariante Ableitung entlang $gamma$]{Definition und Satz}
\label{ssub:211}
Sei $\nabla$ kovariante Ableitung auf $M$, $\gamma:I\to M$ Weg und $\Gamma_\gamma(TM) = \set[Y:I\to TM]{Y(t)\in T_{\gamma(t)} M}$. Elemente aus $\Gamma_\gamma(TM)$ nennen wir "Vektorfelder entlang $\gamma$". Dann existiert eine Abbildung, die sogenannte \Index{kovariante Ableitung entlang $\gamma$} 
\[
\frac{\nabla}{\mathd t}:\Gamma_\gamma(TM) \to \Gamma_\gamma(TM)
\]
,die durch folgende Eigenschaften eindeutig bestimmt ist:
\begin{itemize}[(i)]
\item $\frac{\nabla}{\mathd t}(f Y) = \dot{f}(t)Y(t) + f(t)\frac{\nabla}{\mathd t} Y(t)$ für $f\in C^\infty(I)$
\item Ist $Y(t) = \bar{Y}(\gamma(t))$, wobei $\bar{Y} \in \Gamma(TM)$, so gilt $\enbrace{\frac{\nabla}{\mathd t} Y}(t) = \enbrace{\nabla_{\dot{\gamma}(t)} \bar{Y}}(\gamma(t))$
\end{itemize}
\minisec{Beweis}
\begin{itemize}[(a)]
\item Eindeutigkeit: Sei $Y \in \Gamma_\gamma(TM)$ und $\frac{\nabla}{\mathd t}$ eine kovariante Ableitung entlang $\gamma$ mit den oben genannten Eigenschaften. Sei $\gamma(t) \in U$ Karten mit lokalen Koordinaten $x_1,\ldots,x_n$, so können wir lokal $Y\vert_U = \sum Y^i (t) \partial_i (\gamma(t))$ schreiben. Für $J\subset I$ mit $\gamma(J) \subset U$ gilt mit (i)
\marginnote{Verstehe ich momentan leider nicht, wieso der zweite Summand so aussieht}
\begin{align*}
	\frac{\nabla}{\mathd t} Y\vert_J (t) &= \sum_i \dot{Y}^i(t)\partial_i(\gamma(t)) + Y^i(t) \enbrace{\frac{\nabla}{\mathd t} \partial_i}(t)\\
	& \overset{(ii)}{=} \sum_i \dot{Y}^i(t)\partial_i(\gamma(t)) + Y^i(t) \nabla_{\dot{\gamma}(t)} \partial(\gamma(t))
\end{align*}
Somit ist $\frac{\nabla}{\mathd t} Y$ durch $\nabla$ eindeutig bestimmt.
\item Existenz: Definiere $\enbrace{\frac{\nabla}{\mathd t} Y}(t)$ wie oben auf $J\subset I$, sodass $\gamma(J) \subset U$. Aufgrund der oben gezeigten Eindeutigkeit müssen die so definierten $\frac{\nabla}{\mathd t}$ auf $J_1\cap J_2$ übereinstimmen.\\
(Wobei man hierfür im Beweis der Eindeutigkeit noch die Unabhängigkeit der Kartenwahl zeigen müsste.)
\end{itemize}
%subsubsection 211(end)

\subsubsection{Beispiel}
\label{ssub:212}
Sei $\mathds{R}^n$ mit flacher kovarianter Ableitung gegeben, so folgt aus dem Beweis $\ref{ssub:210}$ direkt, dass für ein Vektorfeld $X:I\to \mathds{R}^n$ entlang $\gamma$ gilt:
\[
\frac{\nabla}{\mathd t} X = \frac{\mathd}{\mathd s} X(s)\vert_t = \dot{X}(t)
\]
%subsubsection 212(end)

\subsubsection{Bemerkung}
\label{ssub:213}
$\nabla$ wird durch die kovariante Ableitung $\enbrace*{\frac{\nabla}{\mathd t}}$ bereits bestimmt. Sei $X,Y$ Vektorfelder und $\alpha:I\to M$ ein Weg mit $\alpha(0) =  x$ und $\dot{\alpha}(0) = X(x)$, so gilt
\[
\enbrace{\nabla_X Y}(x) = \enbrace{\nabla_{X(x)} Y}(x) = \frac{\nabla}{\mathd t}\enbrace{Y\circ \alpha}(0)
\]
%subsubsection 213(end)

\subsubsection{Fortsetzung von $nabla$ auf $\Gamma(\mathcal{T}_q^p M)$}
\label{ssub:214}
Wurde schon in $\ref{ssub:28}$ behandelt.
%subsubsection 214(end)

\subsubsection{Definition}
\label{ssub:215}
Die \Index{Torsion} von $\nabla$ $T^\nabla \in \Gamma(\wedge^2 T^* M \otimes TM)$ ist definiert durch 
\[
T^\nabla (X,Y) = \nabla_X Y - \nabla_Y X - \benbrace{X,Y}
\]
(gesehen als $\underbrace{\Gamma(TM)\times \Gamma(TM)}_{\text{schiefsymmetrisch}} \to \Gamma(TM)$). $\nabla$ heißt \Index{torsionsfrei}, falls $T^\nabla = 0$.\\
Für $f\in C^\infty(M)$ gilt:
\[
T^\nabla (fX,Y) = T^\nabla (X,fY) = f T^\nabla(X,Y)
\]
Somit folgt $T^\nabla \in \Gamma(T^* M \otimes T^* M \otimes TM)$ und da $T^\nabla (X,Y) = - T^\nabla (Y,X)$ folgt $T^\nabla \in \Gamma(\wedge^2 T^* M \otimes TM)$.
%subsubsection 215(end)

\subsubsection{Bemerkung}
\label{ssub:216}
Lokal ist $T^\nabla (\partial_i,\partial_j) = \sum_k \Gamma_{ij}^{\nabla k} \partial_k$. Es gilt
\[
\underset{= \Gamma_{ij}^k}{\nabla_{\partial_i} \partial_j} - \underset{= \Gamma_{ji}^k}\nabla_{\partial_j} \partial_i - \underset{= 0}{\benbrace{\partial_i,\partial_j}}
\]
somit folgt, $\Gamma_{ij}^{\nabla k} = \Gamma_{ij}^k - \Gamma_{ji}^k$, also misst $T^\nabla$ die Kommutativität der Christoffelsymbol in i und j.
%subsubsection 216(end)

\subsubsection{Beispiel}
\label{ssub:217}
Die flache kovariante Ableitung $\nabla = \mathd$ auf $\mathds{R}^n$ ist torsionsfrei ($\Gamma_{ij}^k = 0$).
%subsubsection 217(end)

\subsubsection[Definition: Krümmung von $\nabla$]{Definition}
\label{ssub:218}
Die \Index{Krümmung von $\nabla$} $R^\nabla \in \Gamma(\mathcal{T}_3^1 M)$ (gesehen als $R^\nabla: \Gamma(TM)\times \Gamma(TM) \times \Gamma(TM) \to \Gamma(TM)$) ist definiert durch
\[
R^\nabla (X,Y) Z = \nabla_X \nabla_Y Z - \nabla_Y \nabla_X Z - \nabla_{\benbrace{Y,X}}Z = \enbrace{\nabla_{\benbrace{X,Y}} - [\nabla_X,\nabla_Y]}Z
\]
(Eigentlich $R(X,Y,Z)$, aber wir sehen das hier als $R(X,Y):\Gamma(TM)\to \Gamma(TM)$)
Hier ist $\benbrace{\nabla_X,\nabla_Y} = \nabla_X \circ \nabla_Y - \nabla_Y \circ \nabla_X$ der Abbildungskommutator. (vergleiche mit $\ref{ssub:187}: \mathcal{L}_{\benbrace{X,Y}} - \benbrace{\mathcal{L}_X,\mathcal{L}_Y} = 0$) $\nabla$ ist \Index{flach}, falls $R^\nabla = 0$
%subsubsection 218(end)

\subsubsection{Beispiel}
\label{ssub:219}
Die "flache" kovariante Ableitung auf $\mathds{R}^n$ ist wirklich flach.
%subsubsection 219(end)

\subsubsection{Fundamentallemma der Riemannschen Geometrie}
\label{ssub:220}
Sei $(M,g)$ eine Riemannsche Mannigfaltigkeit. Dann existiert eine eindeutig bestimmte kovariante Ableitung $\nabla^g$ die sogenannte \Index{Levi-Civita kovariante Ableitung} mit folgenden Eigenschaften:
\begin{itemize}
\item $\nabla^g$ ist torsionsfrei
\item $\nabla^g$ ist metrisch, das heißt $\nabla^g g = 0$
\end{itemize}
$\nabla^g_X Y$ wird durch die Koszul-Formel eindeutig bestimmt: Für alle $Z\in \Gamma(TM)$ gilt 
\[
2g\enbrace{\nabla^g_X Y,Z} = X.g(Y,Z) + Y.g(X,Z) - Z.g(X,Y) + g(\benbrace{X,Y},Z) - g(\benbrace{X,Z},Y) - g(\benbrace{Y,Z},X)
\]
(benutzt die Nicht-Degeneriertheit vgl. $\ref{ssub:26}$ Musikalische Isomorphismen).\\
$R^g = R^{\nabla^g}$ heißt Krümmungstensor bzgl. $\nabla^g$ und wir nennen $g$ flach, falls $\nabla^g$ flach.
\minisec{Beweis}
%subsubsection 220(end)

% section 2(end)

\cleardoubleoddemptypage
\pagenumbering{Alph}
\setcounter{page}{1}
\printindex
\listoffigures
\end{document}
